% La consommation de biens et services — Allemand 1ère A — Cours signé OnBuch+
% Programme officiel MINESEC. Compiler avec : tectonic ALA06-la-consommation-de-biens-et-services.tex
\newcommand{\DOCMATIERE}{Allemand}
\newcommand{\DOCNIVEAU}{1ère A}
\newcommand{\DOCMODULE}{Chapitre 2 : Vie économique}
\newcommand{\DOCLECON}{ALA06}
\newcommand{\DOCTITRE}{La consommation de biens et services}
% OnBuch+ — préambule « cours signés » ENRICHI (compilé avec tectonic / XeLaTeX)
% Système visuel « Pop » d'OnBuch+ : fond crème (jamais blanc), bordures encre
% épaisses, ombres « dures » décalées, titres Archivo Black en capitales.
% Version MATHÉMATIQUES : graphes (pgfplots/tikz), boîtes pédagogiques
% (définition, propriété, méthode, attention, le savais-tu, exercice résolu,
% exercice, TP, bilan), page de garde et signature OnBuch+ ancrée en bas.
%
% Macros attendues dans main.tex AVANT \input{preamble} :
%   \DOCMATIERE  \DOCNIVEAU  \DOCMODULE  \DOCLECON (numéro)  \DOCTITRE
\documentclass[11pt]{article}

\usepackage{fontspec}
\usepackage[ngerman,french]{babel} % français = langue principale ; l'allemand via \all{..} / textebox
\usepackage[a4paper,margin=2cm,top=3.4cm,bottom=2.8cm,headheight=1.4cm,footskip=1.3cm]{geometry}
\usepackage{amsmath,amssymb,amsfonts}
\usepackage{mathtools} % \xmapsto, \coloneqq... souvent utilisés par les modèles
\usepackage{cancel} % simplifications barrées (bilans de forces)
% \abs{z} (module, valeur absolue) et \norm{u} (norme), utilisés par les modèles
\providecommand{\abs}{}\let\abs\relax\DeclarePairedDelimiter{\abs}{\lvert}{\rvert}
\providecommand{\norm}{}\let\norm\relax\DeclarePairedDelimiter{\norm}{\lVert}{\rVert}
\usepackage[table]{xcolor} % [table] : fournit aussi \rowcolors (alternance de lignes)
\usepackage{pagecolor}
\usepackage{tikz}
\usetikzlibrary{arrows.meta,positioning,calc,shapes.geometric,shapes.misc,shapes.arrows,decorations.pathmorphing,decorations.pathreplacing,decorations.markings,patterns,shadows,mindmap,trees,fit,backgrounds,matrix,intersections,angles,quotes}
\usepackage{pgfplots}
\usepackage[version=4]{mhchem} % équations nucléaires \ce{^{235}_{92}U}
\usepackage[european,siunitx]{circuitikz} % schémas électriques
\pgfplotsset{compat=1.18}
\usepgfplotslibrary{fillbetween,groupplots} % aires entre courbes (calcul intégral)
\usepackage{enumitem}
\usepackage{fancyhdr}
% [most] charge toutes les bibliothèques tcolorbox (listings, minted, raster,
% documentation...) et gonfle inutilement la mémoire TeX sur un document long
% (~300 boîtes) : on ne charge que ce qui est réellement utilisé.
\usepackage[skins,breakable]{tcolorbox}
\usepackage{graphicx}
\usepackage{colortbl}
\usepackage{booktabs}
\usepackage{tabularx}
\usepackage{diagbox} % case d'en-tête diagonale des tableaux à double entrée
% Colonnes extensibles centrée / gauche / droite (C, L, R), souvent utilisées par les modèles
\newcolumntype{C}{>{\centering\arraybackslash}X}
\newcolumntype{L}{>{\raggedright\arraybackslash}X}
\newcolumntype{R}{>{\raggedleft\arraybackslash}X}
\usepackage{array}
\usepackage{multicol}
\usepackage{multirow}
\usepackage{siunitx}
% parse-numbers=false : accepte aussi \SI{2,5 \times 10^{-3}}{...}
\sisetup{locale=FR,output-decimal-marker={,},group-minimum-digits=4,parse-numbers=false}
\DeclareSIUnit\ohm{\ensuremath{\symup{\Omega}}} % U+2126 absent de la police
\DeclareSIUnit\division{div} % réglages d'oscilloscope (ms/div, V/div)
\DeclareSIUnit\tr{tr} % tours (vitesse de rotation, ex. \SI{1500}{\tr\per\minute})
\DeclareSIUnit\molar{M} % concentration molaire (ex. \SI{3}{\molar}, \SI{1}{\micro\molar})
\DeclareSIUnit\an{an} % année (le modèle confond parfois avec \year, un compteur TeX) — ex. \SI{5}{\kilo\meter\per\an}
\DeclareSIUnit\FCFA{FCFA} % franc CFA (ex. \SI{500}{\FCFA\per\kilo\gram})
\DeclareSIUnit\degreeBrix{\textdegree Brix} % teneur en sucre d'un sirop/jus (réfractométrie)
\DeclareSIUnit\month{mois} % le modèle utilise parfois \month, un compteur TeX, comme unité de durée
\usepackage{qrcode}
% \degree : commande fréquemment utilisée par les modèles pour un angle ou une
% température en dehors de \SI{}{} (non fournie par les paquets chargés).
\providecommand{\degree}{^{\circ}}
% \qed (fin de démonstration), utilisé par les modèles sans amsthm
\providecommand{\qed}{\hfill$\square$}
% \barre : double barre des valeurs interdites dans un tableau de variations (tkz-tab)
\providecommand{\barre}{\ensuremath{\|}}
% environnement proof (amsthm non chargé) utilisé par les modèles
\ifdefined\proof\else\newenvironment{proof}[1][Démonstration]{\par\noindent\textit{#1.}\ }{\qed\par}\fi
\usepackage{lastpage}
\usepackage{hyperref}
\hypersetup{hidelinks}

% ============================================================
% Palette « Pop » OnBuch+ — mêmes hex que la classe P (plus_ui.dart)
% ============================================================
\definecolor{poporange}{HTML}{F59321}
\definecolor{popdark}{HTML}{E07A0C}
\definecolor{popcream}{HTML}{FFF6E8}
\definecolor{popink}{HTML}{1C1714}
\definecolor{poppurple}{HTML}{7A5AE0}
\definecolor{popgreen}{HTML}{2E9E6A}
\definecolor{poppink}{HTML}{E0457B}
\definecolor{popblue}{HTML}{2F7DE1}
\definecolor{popgold}{HTML}{C98A12}
\definecolor{popmuted}{HTML}{8A7F76}
\definecolor{popline}{HTML}{EADFD2}
\definecolor{popgray}{HTML}{8A7F76}
\colorlet{popgrey}{popgray}
% Alias tolérants (le modèle nomme parfois les couleurs différemment)
\colorlet{popwhite}{white}
\colorlet{popblack}{popink}
\colorlet{popred}{poppink}
\colorlet{popyellow}{popgold}
% Teintes claires (fonds de tableaux/encadrés) — le modèle nomme parfois
% les couleurs avec un suffixe L (ex. popblueL) sans les avoir définies.
\colorlet{poporangeL}{poporange!20}
\colorlet{popdarkL}{popdark!20}
\colorlet{poppurpleL}{poppurple!20}
\colorlet{popgreenL}{popgreen!20}
\colorlet{poppinkL}{poppink!20}
\colorlet{popblueL}{popblue!20}
\colorlet{popgoldL}{popgold!20}
\colorlet{popredL}{popred!20}
\colorlet{popgrayL}{popgray!20}
% Teintes claires pour fonds de tableaux / zones
\colorlet{popblueL}{popblue!10!white}
\colorlet{poporangeL}{poporange!14!white}
\colorlet{popgreenL}{popgreen!12!white}
\colorlet{poppurpleL}{poppurple!12!white}
\colorlet{poppinkL}{poppink!12!white}
\colorlet{popgoldL}{popgold!14!white}

\pagecolor{popcream}

% ============================================================
% Typographie
% ============================================================
\defaultfontfeatures{Ligatures=TeX}
\setmainfont{PlusJakartaSans-Medium.ttf}[Path=../../../fonts/,BoldFont=PlusJakartaSans-Bold.ttf]
% Symboles, API et emojis absents de la police principale : repli sur DejaVu Sans (caractères actifs)
\newfontface\symfont{DejaVuSans.ttf}[Path=../../../fonts/]
\makeatletter
\newcommand\symactive[1]{\catcode#1=13 \begingroup\lccode`\~=#1 \lowercase{\endgroup\def~}{{\symfont\char#1}}}
\newcommand\symblank[1]{\catcode#1=13 \begingroup\lccode`\~=#1 \lowercase{\endgroup\def~}{}}
\newcommand\symhyph[1]{\catcode#1=13 \begingroup\lccode`\~=#1 \lowercase{\endgroup\def~}{\mbox{-}}}
\newcount\symc
\newcommand\symrange[2]{\symc=#1 \@whilenum\symc<#2 \do{\expandafter\symactive\expandafter{\the\symc}\advance\symc by 1 }}
\newcommand\symblankrange[2]{\symc=#1 \@whilenum\symc<#2 \do{\expandafter\symblank\expandafter{\the\symc}\advance\symc by 1 }}
\makeatother
\symrange{"0250}{"02FF}\symactive{"2713}\symactive{"2714}\symactive{"2717}\symactive{"2718}\symactive{"2610}\symactive{"2611}\symactive{"2612}
\symhyph{"2011}\symhyph{"2E1A}\symblankrange{"1F300}{"1FAFF}\symblank{"FE0F}
% Maths en sans-serif (Fira Math) avec chiffres et lettres droites de Plus
% Jakarta, pour que les formules se fondent dans le texte.
\usepackage[math-style=ISO,bold-style=ISO]{unicode-math}
\setmathfont{FiraMath-Regular.otf}[Path=../../../fonts/]
\setmathfont{PlusJakartaSans-Medium.ttf}[Path=../../../fonts/,range={up/num,up/{Latin,latin},\mathup}]
\usepackage{icomma} % virgule décimale sans espace en mode math
% Caractères mathématiques Unicode absents des polices de texte (Plus Jakarta,
% Archivo Black) : rendus par leur équivalent mathématique, en texte comme en
% formule — sinon ils s'affichent en carré vide (ex. « Arithmétique dans ℤ »).
\usepackage{newunicodechar}
\newunicodechar{ℝ}{\ensuremath{\mathbb{R}}}
\newunicodechar{ℤ}{\ensuremath{\mathbb{Z}}}
\newunicodechar{ℂ}{\ensuremath{\mathbb{C}}}
\newunicodechar{ℕ}{\ensuremath{\mathbb{N}}}
\newunicodechar{ℚ}{\ensuremath{\mathbb{Q}}}
\newunicodechar{𝔻}{\ensuremath{\mathbb{D}}}
\newunicodechar{α}{\ensuremath{\alpha}}
\newunicodechar{β}{\ensuremath{\beta}}
\newunicodechar{γ}{\ensuremath{\gamma}}
\newunicodechar{δ}{\ensuremath{\delta}}
\newunicodechar{ε}{\ensuremath{\varepsilon}}
\newunicodechar{θ}{\ensuremath{\theta}}
\newunicodechar{λ}{\ensuremath{\lambda}}
\newunicodechar{μ}{\ensuremath{\mu}}
\newunicodechar{σ}{\ensuremath{\sigma}}
\newunicodechar{τ}{\ensuremath{\tau}}
\newunicodechar{φ}{\ensuremath{\varphi}}
\newunicodechar{ψ}{\ensuremath{\psi}}
\newunicodechar{ω}{\ensuremath{\omega}}
\newunicodechar{Σ}{\ensuremath{\Sigma}}
% majuscules produites par \MakeUppercase dans les titres (α-aminés -> Α)
\newunicodechar{Α}{\ensuremath{\alpha}}
\newunicodechar{Β}{\ensuremath{\beta}}
\newunicodechar{Γ}{\ensuremath{\Gamma}}
\newunicodechar{Δ}{\ensuremath{\Delta}}
\newunicodechar{Ε}{\ensuremath{\varepsilon}}
\newunicodechar{Θ}{\ensuremath{\Theta}}
\newunicodechar{Λ}{\ensuremath{\Lambda}}
\newunicodechar{Μ}{\ensuremath{\mu}}
\newunicodechar{Τ}{\ensuremath{\tau}}
\newunicodechar{Φ}{\ensuremath{\Phi}}
\newunicodechar{Ψ}{\ensuremath{\Psi}}
\newunicodechar{Ω}{\ensuremath{\Omega}}
\newunicodechar{↦}{\ensuremath{\mapsto}}
\newunicodechar{∈}{\ensuremath{\in}}
\newunicodechar{∉}{\ensuremath{\notin}}
\newunicodechar{∧}{\ensuremath{\wedge}}
\newunicodechar{⇔}{\ensuremath{\Leftrightarrow}}
\newunicodechar{⟺}{\ensuremath{\Longleftrightarrow}}
\newunicodechar{⇒}{\ensuremath{\Rightarrow}}
\newunicodechar{⟹}{\ensuremath{\Longrightarrow}}
\newunicodechar{∘}{\ensuremath{\circ}}
\newunicodechar{∪}{\ensuremath{\cup}}
\newunicodechar{∩}{\ensuremath{\cap}}
\newunicodechar{≡}{\ensuremath{\equiv}}
\newunicodechar{⊂}{\ensuremath{\subset}}
\newunicodechar{⊥}{\ensuremath{\perp}}
\newunicodechar{∃}{\ensuremath{\exists}}
\newunicodechar{∀}{\ensuremath{\forall}}
\newunicodechar{∛}{\ensuremath{\sqrt[3]{\,}}}
\newunicodechar{∜}{\ensuremath{\sqrt[4]{\,}}}
\newunicodechar{⃗}{\ensuremath{{}^{\to}}}
\newunicodechar{ⁿ}{\ifmmode^{n}\else\textsuperscript{\ensuremath{n}}\fi}
\newunicodechar{ˣ}{\ifmmode^{x}\else\textsuperscript{\ensuremath{x}}\fi}
\newunicodechar{ᵇ}{\ifmmode^{b}\else\textsuperscript{\ensuremath{b}}\fi}
\newunicodechar{ʳ}{\ifmmode^{r}\else\textsuperscript{\ensuremath{r}}\fi}
\newunicodechar{ᵘ}{\ifmmode^{u}\else\textsuperscript{\ensuremath{u}}\fi}
\newunicodechar{ʸ}{\ifmmode^{y}\else\textsuperscript{\ensuremath{y}}\fi}
\newunicodechar{ᵃ}{\ifmmode^{a}\else\textsuperscript{\ensuremath{a}}\fi}
\newunicodechar{ᵉ}{\ifmmode^{e}\else\textsuperscript{\ensuremath{e}}\fi}
\newunicodechar{ᵏ}{\ifmmode^{k}\else\textsuperscript{\ensuremath{k}}\fi}
\newunicodechar{ᵖ}{\ifmmode^{p}\else\textsuperscript{\ensuremath{p}}\fi}
\newunicodechar{ᴺ}{\ifmmode^{N}\else\textsuperscript{\ensuremath{N}}\fi}
\newunicodechar{ᶜ}{\ifmmode^{c}\else\textsuperscript{\ensuremath{c}}\fi}
\newunicodechar{ʰ}{\ifmmode^{h}\else\textsuperscript{\ensuremath{h}}\fi}
\newunicodechar{ᵠ}{\ifmmode^{q}\else\textsuperscript{\ensuremath{q}}\fi}
\newunicodechar{ₙ}{\ifmmode_{n}\else\textsubscript{\ensuremath{n}}\fi}
\newunicodechar{ₐ}{\ifmmode_{a}\else\textsubscript{\ensuremath{a}}\fi}
\newunicodechar{ᵢ}{\ifmmode_{i}\else\textsubscript{\ensuremath{i}}\fi}
\newunicodechar{ᵣ}{\ifmmode_{r}\else\textsubscript{\ensuremath{r}}\fi}
\newunicodechar{ₖ}{\ifmmode_{k}\else\textsubscript{\ensuremath{k}}\fi}
\newunicodechar{ᵦ}{\ifmmode_{\beta}\else\textsubscript{\ensuremath{\beta}}\fi}
\newunicodechar{ₓ}{\ifmmode_{x}\else\textsubscript{\ensuremath{x}}\fi}
\newfontfamily\popdisplay{ArchivoBlack-Regular.ttf}[Path=../../../fonts/]
\newfontfamily\poplogo{SpaceGrotesk-Bold.ttf}[Path=../../../fonts/]
\newfontfamily\popsemi{PlusJakartaSans-SemiBold.ttf}[Path=../../../fonts/]
\newfontfamily\popxbold{PlusJakartaSans-ExtraBold.ttf}[Path=../../../fonts/]
\newfontfamily\popxboldls{PlusJakartaSans-ExtraBold.ttf}[Path=../../../fonts/,LetterSpace=3.0]

\setlist[enumerate]{leftmargin=1.7em,itemsep=5pt,topsep=4pt}
\setlist[itemize]{leftmargin=1.7em,itemsep=4pt,topsep=4pt,label={\color{poporange}\textbullet}}
\setlength{\parindent}{0pt}
\setlength{\parskip}{5pt}
\renewcommand{\arraystretch}{1.25}

% pgfplots : style Pop par défaut
\pgfplotsset{
  every axis/.append style={axis line style={popink,line width=1pt},tick style={popink},
    grid style={popline,line width=0.5pt},label style={font=\small},tick label style={font=\footnotesize}},
  popcurve/.style={poporange,line width=1.8pt,no markers,smooth},
}
% Même style disponible dans un \draw TikZ simple (hors pgfplots)
\tikzset{popcurve/.style={poporange,line width=1.8pt}}
\tikzset{popgrid/.style={popline,very thin}}
\colorlet{popcurve}{poporange} % le nom du style est parfois utilisé comme couleur

% ============================================================
% Pill / squiggle
% ============================================================
\newcommand{\poppill}[3]{%
  \tikz[baseline=-0.55ex]\node[draw=popink,line width=1.2pt,rounded corners=2.6pt,%
    fill=#2,inner xsep=6.5pt,inner ysep=3pt]{%
    \popxboldls\fontsize{7}{7}\selectfont\color{#3}\MakeUppercase{#1}};%
}
\newcommand{\popsquiggle}[1]{%
  \tikz[baseline=-0.4ex]{
    \draw[#1,line width=2.6pt,line cap=round]
      plot [smooth] coordinates {(0,0.09) (0.34,0.01) (0.68,0.21) (1.02,0.01) (1.36,0.10)};
  }%
}
% Mot-clé surligné (marqueur orange sous le texte)
\newcommand{\cle}[1]{{\popxbold\color{popdark}#1}}

% ============================================================
% En-tête / pied
% ============================================================
\pagestyle{fancy}
\fancyhf{}
\renewcommand{\headrulewidth}{0pt}
\renewcommand{\footrulewidth}{0pt}
\lhead{\raisebox{-0.15ex}{\poplogo\fontsize{13}{13}\selectfont\color{popink}OnBuch\color{poporange}+}}
\rhead{\poppill{\DOCMATIERE\ \textbullet\ \DOCNIVEAU\ \textbullet\ Leçon \DOCLECON}{white}{popink}}
\lfoot{\popxbold\fontsize{7.6}{7.6}\selectfont\color{popmuted}\MakeUppercase{Cours signé OnBuch+}\ \textbullet\ {\popsemi\DOCTITRE}}
\rfoot{\tikz[baseline=-0.6ex]\node[rounded corners=8.5pt,fill=popink,minimum height=17pt,minimum width=17pt,inner xsep=5pt,inner sep=0]{\popxbold\fontsize{8}{8}\selectfont\color{popcream}\thepage\,/\,\pageref*{LastPage}};}

% ============================================================
% Titres
% ============================================================
\newcommand{\chaptitre}[2]{%
  \begin{tcolorbox}[enhanced,colback=poporange,colframe=popink,boxrule=2.6pt,arc=11pt,
      drop shadow=popink,left=15pt,right=15pt,top=12pt,bottom=11pt,before skip=3pt,after skip=18pt]
    \poppill{#2}{popink}{popcream}\par\vspace{9pt}
    {\raggedright\hyphenpenalty=10000\exhyphenpenalty=10000\popdisplay\fontsize{22}{24}\selectfont\color{popcream}\MakeUppercase{#1}\par}
  \end{tcolorbox}
}
% Section numérotée : pastille orange avec le numéro + titre Archivo
\newcounter{popsec}
\newcommand{\coursec}[1]{%
  \stepcounter{popsec}\setcounter{popsub}{0}%
  \par\vspace{16pt}\noindent
  \begin{minipage}{\linewidth}
  \tikz[baseline=-0.7ex]\node[draw=popink,line width=1.6pt,fill=poporange,rounded corners=4pt,minimum size=22pt,inner sep=0]{\popdisplay\fontsize{12}{12}\selectfont\color{popcream}\thepopsec};%
  \hspace{8pt}{\popdisplay\fontsize{15}{17}\selectfont\color{popink}\MakeUppercase{#1}}\par
  \vspace{4pt}\noindent{\color{poporange}\rule{36pt}{3.6pt}}
  \end{minipage}\par\nopagebreak\vspace{8pt}
}
\newcounter{popsub}
\newcommand{\courssub}[1]{%
  \stepcounter{popsub}%
  \par\vspace{9pt}\noindent{\popxbold\fontsize{11.5}{13}\selectfont\color{popink}{\color{poporange}\thepopsec.\thepopsub}\hspace{6pt}#1}\par\nopagebreak\vspace{4pt}
}

% ============================================================
% Boîtes pédagogiques — même recette : fond blanc, bordure épaisse colorée,
% ombre dure encre, badge plein. Titre optionnel [#1].
% ============================================================
\tcbset{popbase/.style={breakable,enhanced,colback=white,boxrule=2.3pt,arc=9pt,
  shadow={3pt}{-3pt}{0pt}{popink},left=14pt,right=14pt,top=11pt,bottom=11pt,before skip=11pt,after skip=15pt,
  fontupper=\popsemi\fontsize{10.6}{14.5}\selectfont\color{popink},
  fontlower=\popsemi\fontsize{10.6}{14.5}\selectfont\color{popink}}}
% Coche dessinée (le glyphe ✓ n'existe pas dans les polices du document)
\newcommand{\popcheck}[1]{\tikz[baseline=-0.5ex]\draw[#1,line width=1.6pt,line cap=round,line join=round](-0.13,0.0)--(-0.04,-0.09)--(0.13,0.1);}
\providecommand{\coche}{\popcheck{popgreen}}
\providecommand{\resultat}[1]{\par\smallskip\noindent\fcolorbox{popgreen}{popgreenL}{\parbox{\dimexpr\linewidth-2\fboxsep-2\fboxrule\relax}{\textbf{Résultat.} #1}}\par\smallskip}
\newcommand{\popboxhead}[3]{\poppill{#1}{#2}{white}\ifx\relax#3\relax\else\hspace{6pt}{\popxbold\fontsize{10.6}{12}\selectfont\color{popink}#3}\fi\par\vspace{7pt}}

\newtcolorbox{aretenir}[1][]{popbase,colframe=popgreen,before upper={\popboxhead{À retenir}{popgreen}{#1}}}
% Texte en allemand (document de lecture, dialogue, extrait) : cadre
% d'étude. Un titre optionnel (source, auteur) s'affiche dans l'en-tête.
\newtcolorbox{textebox}[1][]{popbase,colframe=popblue,before upper={\popboxhead{Text}{popblue}{#1}\begin{otherlanguage}{ngerman}},after upper={\end{otherlanguage}}}
% Marques ✓ / ✗ (forme correcte / fautive) souvent utilisées par les modèles
\providecommand{\cross}{\textcolor{poppink}{\ensuremath{\times}}}
\providecommand{\xmark}{\cross}\providecommand{\crossmark}{\cross}\providecommand{\cmark}{\ensuremath{\checkmark}}
% Ligne de réponse / justification à compléter (inventée par les modèles)
\providecommand{\justification}[1]{\par\noindent\textit{Justification~:} \dotfill\par}
% \textipa (package tipa non chargé) : la police ne couvre pas l'API -> simple passage
\providecommand{\textipa}[1]{#1}
% Soulignement (ulem non chargé) : \uline, \ul
\providecommand{\uline}[1]{\underline{#1}}\providecommand{\ul}[1]{\underline{#1}}
% \glossaire{\item ...} inventée par les modèles : liste de glossaire
\providecommand{\glossaire}[1]{\begin{itemize}#1\end{itemize}}
% \alert (beamer) inventée par les modèles : texte mis en évidence
\providecommand{\alert}[1]{\textbf{\textcolor{poppink}{#1}}}
% variantes ASCII d'ordinaux écrites en macro
\providecommand{\iere}{\textsuperscript{re}}\providecommand{\ieme}{\textsuperscript{e}}\providecommand{\ere}{\textsuperscript{re}}\providecommand{\eme}{\textsuperscript{e}}\providecommand{\er}{\textsuperscript{er}}\providecommand{\ie}{\textsuperscript{e}}
% Ordinaux français écrits en macro par les modèles : 1\ière, 3\ième
\newcommand{\ière}{\textsuperscript{re}}\newcommand{\ième}{\textsuperscript{e}}
% \exemple (inventée par les modèles) : étiquette « Exemple : »
\providecommand{\exemple}{\textbf{Exemple~:}\ }
% \square utilisable aussi hors mode math (cases à cocher)
\AtBeginDocument{\let\squareMath\square\renewcommand{\square}{\ensuremath{\squareMath}}}
% Mot ou phrase en allemand à l'intérieur d'un paragraphe français
\newcommand{\all}[1]{\ifmmode\text{\textit{#1}}\else\begingroup\selectlanguage{ngerman}\itshape #1\endgroup\fi}
\let\esp\all % alias toléré
\newtcolorbox{exemplebox}[1][]{popbase,colframe=poporange,before upper={\popboxhead{Exemple}{poporange}{#1}}}
\newtcolorbox{definition}[1][]{popbase,colframe=popblue,before upper={\popboxhead{Définition}{popblue}{#1}}}
\newtcolorbox{propriete}[1][]{popbase,colframe=poppurple,before upper={\popboxhead{Propriété}{poppurple}{#1}}}
\newtcolorbox{methode}[1][]{popbase,colframe=popgold,before upper={\popboxhead{Méthode}{popgold}{#1}}}
\newtcolorbox{attention}[1][]{popbase,colframe=poppink,before upper={\popboxhead{Attention}{poppink}{#1}}}
\newtcolorbox{experience}[1][]{popbase,colframe=popblue,colback=popblueL,before upper={\popboxhead{Expérience}{popblue}{#1}}}
% « Le savais-tu ? » : carte violette pleine (contexte, culture, Cameroun)
\newtcolorbox{savaistu}[1][]{popbase,colback=poppurple,colframe=popink,
  fontupper=\popsemi\fontsize{10.4}{14.2}\selectfont\color{white},
  before upper={\poppill{Le savais-tu\,?}{white}{poppurple}\ifx\relax#1\relax\else\hspace{6pt}{\popxbold\color{white}#1}\fi\par\vspace{7pt}}}
% Exercice résolu : énoncé en haut, \tcblower puis la solution sur fond crème
\newcounter{popexo}\newcounter{popexr}
\newtcolorbox{exoresolu}[1][]{popbase,colframe=poporange,colbacklower=popcream,
  segmentation style={popink,line width=1.2pt,dashed},
  before upper={\stepcounter{popexr}\popboxhead{Exercice résolu \thepopexr}{poporange}{#1}},
  before lower={\poppill{Solution}{popink}{popcream}\par\vspace{6pt}}}
% Exercice (non résolu en place) — la correction est dans la partie Corrigés
\newtcolorbox{exercice}[1][]{popbase,colframe=popink,
  before upper={\stepcounter{popexo}\popboxhead{Exercice \thepopexo}{popink}{#1}}}
% Corrigé
\newtcolorbox{corrige}[1][]{popbase,colframe=popgreen,colback=popgreenL,
  before upper={\popboxhead{Corrigé}{popgreen}{#1}}}
% Objectifs / prérequis / bilan
\newtcolorbox{objectifs}{popbase,colframe=popink,colback=poporangeL,
  before upper={\popboxhead{Objectifs de la leçon}{popink}{}}}
\newtcolorbox{prerequis}{popbase,colframe=popink,colback=popblueL,
  before upper={\popboxhead{Prérequis}{popblue}{}}}
\newtcolorbox{bilan}{popbase,colframe=popink,colback=white,boxrule=2.8pt,
  before upper={\popboxhead{Fiche bilan}{poporange}{L'essentiel à connaître pour l'examen}}}
% Figure encadrée avec légende
\newtcolorbox{popfigure}[1][]{enhanced,breakable=false,colback=white,colframe=popline,boxrule=1.4pt,arc=8pt,
  left=10pt,right=10pt,top=10pt,bottom=8pt,before skip=10pt,after skip=12pt,halign=center,
  lower separated=false,halign lower=center,
  fontlower=\popsemi\fontsize{9}{11.5}\selectfont\color{popmuted},
  #1}
\newcommand{\legende}[1]{\par\vspace{6pt}{\popsemi\fontsize{9}{11.5}\selectfont\color{popmuted}#1}}

% ============================================================
% Signature OnBuch+
% ============================================================
\newcommand{\poplogoinline}[1]{{\poplogo\fontsize{#1}{#1}\selectfont\color{popink}OnBuch\color{poporange}+}}
\newcommand{\signatureonbuch}{%
  \begin{tcolorbox}[enhanced,colback=popink,colframe=popink,boxrule=2.2pt,arc=10pt,
      drop shadow=poporange,left=14pt,right=14pt,top=10pt,bottom=10pt,before skip=0pt,after skip=0pt]
    \tikz[baseline=-0.7ex]\node[circle,fill=popgreen,draw=popcream,line width=1.3pt,minimum size=22pt,inner sep=0]{\popcheck{white}};%
    \hspace{9pt}\begin{minipage}[c]{0.62\linewidth}
      {\popxbold\fontsize{12}{13}\selectfont\color{popcream}Signé\ \textbullet\ {\poplogo OnBuch\color{poporange}+}}\par\vspace{2pt}
      {\raggedright\popsemi\fontsize{8}{10.5}\selectfont\color{popline}\MakeUppercase{Équipe pédagogique OnBuch+}\\\MakeUppercase{Conforme au programme officiel MINESEC}\par}
    \end{minipage}\hfill
    {\poplogo\fontsize{22}{22}\selectfont\color{popcream}OnBuch\color{poporange}+}
  \end{tcolorbox}%
}

% ============================================================
% Page de garde
% #1 titre · #2 matière · #3 niveau · #4 module · #5 n° leçon · #6 durée ·
% #7 description / objectifs (texte ou itemize)
% ============================================================
\newcommand{\pagegarde}[7]{%
  \thispagestyle{empty}
  \newgeometry{a4paper,margin=1.7cm,top=1.6cm,bottom=1.4cm}%
  \begin{tikzpicture}[remember picture,overlay]
    \fill[poporange!18] ([xshift=-3.2cm,yshift=-3.6cm]current page.north east) circle (4.6cm);
    \fill[poppurple!14] ([xshift=2.4cm,yshift=7.6cm]current page.south west) circle (3.8cm);
  \end{tikzpicture}%
  {\poplogo\fontsize{30}{30}\selectfont\color{popink}OnBuch\color{poporange}+}\hfill
  \raisebox{0.6ex}{\poppill{Cours signé \textbullet\ Premium}{popink}{popcream}}\par
  \vspace{1pt}\popsquiggle{poppurple}\par
  \vspace{8pt}
  \begin{tcolorbox}[enhanced,colback=poporange,colframe=popink,boxrule=2.8pt,arc=13pt,
      drop shadow=popink,left=18pt,right=18pt,top=12pt,bottom=13pt,after skip=10pt]
    \poppill{#2 \textbullet\ #3}{popink}{popcream}\hspace{5pt}\poppill{#4}{white}{popink}\par\vspace{8pt}
    {\popdisplay\fontsize{14}{14}\selectfont\color{popink}LEÇON #5}\par\vspace{4pt}
    {\raggedright\hyphenpenalty=10000\exhyphenpenalty=10000\popdisplay\fontsize{25}{27}\selectfont\color{popcream}\MakeUppercase{#1}\par}
    \vspace{8pt}
    {\popxbold\fontsize{9.5}{9.5}\selectfont\color{popink}\MakeUppercase{Durée conseillée : #6}}
  \end{tcolorbox}
  \begin{tcolorbox}[enhanced,colback=white,colframe=popink,boxrule=2.2pt,arc=10pt,
      drop shadow=popink,left=16pt,right=16pt,top=10pt,bottom=10pt,after skip=8pt]
    \poppill{Dans cette leçon}{poppurple}{white}\par\vspace{5pt}
    \popsemi\fontsize{9.3}{12.6}\selectfont\color{popink}#7
  \end{tcolorbox}
  \noindent\begin{minipage}[t]{0.487\linewidth}
    \begin{tcolorbox}[enhanced,colback=popgreen,colframe=popink,boxrule=2.2pt,arc=10pt,
        drop shadow=popink,left=11pt,right=11pt,top=10pt,bottom=10pt,height=3.1cm,valign=center]
      \tcbox[colback=white,colframe=popink,boxrule=1.3pt,arc=5pt,boxsep=0pt,nobeforeafter,
        left=3pt,right=3pt,top=3pt,bottom=3pt,valign=center]{\qrcode[height=1.9cm]{https://play.google.com/store/apps/details?id=com.luvvix.onbuch&hl=fr}}%
      \hspace{8pt}\begin{minipage}[c]{0.5\linewidth}
        {\popdisplay\fontsize{11}{12.5}\selectfont\color{white}\MakeUppercase{Télécharge OnBuch}}\par\vspace{4pt}
        {\raggedright\popsemi\fontsize{8.4}{11}\selectfont\color{white}Gratuit \textbullet\ Google Play\\Tuteur IA, annales, concours, résultats\par}
      \end{minipage}
    \end{tcolorbox}
  \end{minipage}\hfill
  \begin{minipage}[t]{0.487\linewidth}
    \begin{tcolorbox}[enhanced,colback=poppurple,colframe=popink,boxrule=2.2pt,arc=10pt,
        drop shadow=popink,left=11pt,right=11pt,top=10pt,bottom=10pt,height=3.1cm,valign=center]
      \tikz[baseline=-0.7ex]\node[circle,fill=white,draw=popink,line width=1.3pt,minimum size=1.6cm,inner sep=0]{\popdisplay\fontsize{24}{24}\selectfont\color{poppurple}+};%
      \hspace{8pt}\begin{minipage}[c]{0.6\linewidth}
        {\popdisplay\fontsize{11}{12.5}\selectfont\color{white}\MakeUppercase{Passe à OnBuch+}}\par\vspace{4pt}
        {\raggedright\popsemi\fontsize{8.4}{11}\selectfont\color{white}Tous les cours signés, Tuteur IA illimité, fiches PDF — onglet OnBuch+ de l'app\par}
      \end{minipage}
    \end{tcolorbox}
  \end{minipage}
  \vspace{8pt}
  \popcontenu
  \vfill
  \signatureonbuch
  \restoregeometry
  \newpage
}

% Rangée « Ce cours contient » (page de garde)
\newcommand{\poptile}[3]{% #1 couleur #2 gros texte #3 légende
  \begin{tcolorbox}[enhanced,colback=#1,colframe=popink,boxrule=2pt,arc=9pt,shadow={2.5pt}{-2.5pt}{0pt}{popink},
      left=6pt,right=6pt,top=9pt,bottom=9pt,halign=center,valign=center,height=1.75cm,width=\linewidth,nobeforeafter]
    {\popdisplay\fontsize{12.5}{13}\selectfont\color{white}\MakeUppercase{#2}}\par\vspace{3pt}
    {\centering\popsemi\fontsize{7.8}{9.5}\selectfont\color{white}#3\par}
  \end{tcolorbox}}
\newcommand{\popcontenu}{%
  \par\noindent{\popxboldls\fontsize{7.5}{7.5}\selectfont\color{popmuted}CE COURS SIGNÉ CONTIENT}\par\vspace{7pt}
  \noindent\begin{minipage}[t]{0.235\linewidth}\poptile{popblue}{Cours}{complet, illustré, pas à pas}\end{minipage}\hfill
  \begin{minipage}[t]{0.235\linewidth}\poptile{poporange}{Méthodes}{et exercices résolus}\end{minipage}\hfill
  \begin{minipage}[t]{0.235\linewidth}\poptile{poppink}{Exercices}{type examen + corrigés}\end{minipage}\hfill
  \begin{minipage}[t]{0.235\linewidth}\poptile{popgreen}{Fiche bilan}{l'essentiel à retenir}\end{minipage}\par}

% Sommaire
\newtcolorbox{sommaire}{enhanced,colback=white,colframe=popink,boxrule=2.2pt,arc=10pt,
  drop shadow=popink,left=18pt,right=18pt,top=13pt,bottom=13pt,before skip=6pt,after skip=20pt,
  before upper={\poppill{Sommaire}{poppurple}{white}\par\vspace{9pt}\popsemi\fontsize{10.6}{14.5}\selectfont\color{popink}}}

% Signature de fin de cours — poussée tout en bas de la dernière page
\newcommand{\findecours}{%
  \par\vfill\nopagebreak
  \begin{center}\popsquiggle{poporange}\end{center}\vspace{4pt}
  \signatureonbuch
}
\AtBeginDocument{\def\llbracket{\mathopen{[\mkern-3.5mu[}}\def\rrbracket{\mathclose{]\mkern-3.5mu]}}} % intervalles d'entiers (glyphe ⟦ absent de la police)
% Glyphes absents de Fira Math : redéfinis à partir de glyphes disponibles
\AtBeginDocument{%
  \def\setminus{\mathbin{\backslash}}\let\smallsetminus\setminus
  \def\vdots{\mathord{\vbox{\baselineskip3.2pt\lineskiplimit0pt\kern4pt\hbox{.}\hbox{.}\hbox{.}}}}%
  \def\ddots{\mathinner{\mkern1mu\raise6.5pt\hbox{.}\mkern2mu\raise3.6pt\hbox{.}\mkern2mu\raise0.7pt\hbox{.}\mkern1mu}}%
  \def\triangle{\mathord{\scalebox{0.9}{$\symup{\Delta}$}}}%
  \def\nabla{\mathord{\raisebox{\depth}{\scalebox{1}[-1]{$\symup{\Delta}$}}}}%
  \def\checkmark{\popcheck{popdark}}%
  \def\star{\mathbin{\ast}}\def\bigstar{\mathord{\ast}}%
  \def\diamond{\mathbin{\scalebox{0.6}{$\Diamond$}}}%
  \def\bigcup{\mathop{\mathchoice{\raisebox{-0.3em}{\scalebox{1.7}{$\cup$}}}{\scalebox{1.25}{$\cup$}}{\cup}{\cup}}\displaylimits}%
  \def\bigcap{\mathop{\mathchoice{\raisebox{-0.3em}{\scalebox{1.7}{$\cap$}}}{\scalebox{1.25}{$\cap$}}{\cap}{\cap}}\displaylimits}%
  \def\lesseqqgtr{\mathrel{\lessgtr}}\def\bigoplus{\mathop{\oplus}\displaylimits}%
}
\providecommand{\ssi}{si et seulement si} % abréviation fréquente
\AtBeginDocument{\providecommand{\wideparen}{\overparen}} % arc de cercle

%% FIGFIX-BEGIN v3 — correctifs globaux des figures (docs/onbuch-generation/tools/figfix)
\usepackage{adjustbox}\usepackage{etoolbox}
% toute figure TikZ/pgfplots plus large que la ligne est réduite à la largeur disponible
\BeforeBeginEnvironment{tikzpicture}{\begin{adjustbox}{max width=\linewidth}}
\AfterEndEnvironment{tikzpicture}{\end{adjustbox}}
% légendes pgfplots lisibles par-dessus les courbes
\pgfplotsset{every axis/.append style={legend style={fill=white,fill opacity=0.92,text opacity=1,draw=black!15,inner sep=3pt}}}
% étiquettes d'arêtes (syntaxe edge["..."]) avec fond blanc
\tikzset{every edge quotes/.append style={fill=white,inner sep=1.2pt}}
% bulles de cartes mentales : texte à la ligne dans la bulle, bulle un peu plus grande
\tikzset{every concept/.append style={text width=2.0cm,align=center,minimum size=2.3cm,inner sep=2pt}}
%% FIGFIX-END
\begin{document}

\pagegarde{La consommation de biens et services}{Allemand}{1ère A}{Chapitre 2 : Vie économique}{ALA06}{2 h}{À partir d'un texte sur les habitudes de consommation d'une famille allemande, cette leçon développe le vocabulaire de l'achat et des services et consolide deux points majeurs de grammaire : la déclinaison de l'adjectif et le comparatif/superlatif, ainsi que les verbes à particule séparable.
\vspace{4pt}
\begin{multicols}{2}\small
\begin{itemize}
\item À la fin de cette leçon, je sais comprendre globalement et en détail un texte sur la consommation de biens et de services.
\item Je sais employer correctement le vocabulaire thématique (kaufen, der Preis, das Geschäft, die Ware, der Kunde, die Dienstleistung) avec articles et pluriels.
\item Je sais décliner les adjectifs après article défini, indéfini et sans article.
\item Je sais former et employer le comparatif et le superlatif de l'adjectif.
\item Je sais conjuguer et placer correctement les verbes à particule séparable dans la phrase.
\item Je sais répondre par écrit et à l'oral à des questions de compréhension sur le texte.
\end{itemize}\end{multicols}}

\begin{sommaire}
\begin{multicols}{2}
\begin{enumerate}
\item Texte support : Einkaufen in Deutschland
\item Compréhension du texte
\item Lexique thématique : Konsum und Dienstleistungen
\item Grammaire 1 : la déclinaison de l'adjectif
\item Grammaire 2 : comparatif et superlatif
\item Conjugaison : les verbes à particule séparable
\item Expression guidée et méthode du commentaire
\item Activité d'intégration
\item Méthodes et exercices résolus
\item Exercices
\item Corrigés des exercices
\item Fiche bilan
\end{enumerate}
\end{multicols}
\end{sommaire}

\begin{objectifs}
\begin{itemize}
\item À la fin de cette leçon, je sais comprendre globalement et en détail un texte sur la consommation de biens et de services.
\item Je sais employer correctement le vocabulaire thématique (kaufen, der Preis, das Geschäft, die Ware, der Kunde, die Dienstleistung) avec articles et pluriels.
\item Je sais décliner les adjectifs après article défini, indéfini et sans article.
\item Je sais former et employer le comparatif et le superlatif de l'adjectif.
\item Je sais conjuguer et placer correctement les verbes à particule séparable dans la phrase.
\item Je sais répondre par écrit et à l'oral à des questions de compréhension sur le texte.
\item Je sais résumer un texte et produire un court paragraphe guidé sur mes habitudes de consommation.
\item Je sais repérer et éviter les erreurs fréquentes des francophones (place de la particule, déclinaisons, faux amis).
\end{itemize}
\end{objectifs}

\begin{prerequis}
\begin{itemize}
\item Les articles définis et indéfinis (der/die/das, ein/eine) et les quatre cas (Seconde)
\item Le présent des verbes réguliers et irréguliers courants (kaufen, geben, nehmen)
\item Le vocabulaire de base de la vie quotidienne (maison, famille, argent) vu en Seconde
\item La structure de la phrase allemande : verbe conjugué en 2e position
\item Les adjectifs qualificatifs simples (gut, teuer, billig, schön)
\end{itemize}
\end{prerequis}

\newpage

\coursec{Texte support : Einkaufen in Deutschland}

C'est samedi matin au marché Mokolo à Yaoundé : Amina compare les prix des tomates, négocie avec la vendeuse et choisit les plus belles. En Allemagne, Lukas fait la même chose au supermarché ou dans une boutique (\all{das Geschäft}) : il regarde les prix (\all{die Preise}), compare la qualité des marchandises (\all{die Waren}) et paie à la caisse. Consommer, c'est acheter des biens et des services partout dans le monde — mais comment en parler en allemand ? Comment dire qu'un produit est « plus cher » ou « le meilleur » ? Cette leçon nous plonge dans le quotidien du consommateur allemand.

\textbf{Problématique :} Comment un texte allemand décrit-il l'acte d'achat et la consommation, et quels outils de langue (vocabulaire, comparaison, verbes) nous permettent de le comprendre et de le raconter à notre tour ?

\courssub{Objectifs de la séquence}

À la fin de cette section, tu seras capable de :
\begin{itemize}
\item lire et comprendre un texte allemand simple sur les achats et les services ;
\item repérer les informations globales : les personnes (\all{die Personen}), les lieux (\all{die Orte}), les actions (\all{die Handlungen}) ;
\item utiliser un glossaire pour lever les difficultés de vocabulaire ;
\item souligner le vocabulaire du thème et les comparatifs dans le texte ;
\item reformuler oralement le contenu du texte avec tes propres mots (\all{eine Zusammenfassung}).
\end{itemize}

\courssub{Méthode : comment lire un texte allemand}

\begin{methode}[Les cinq étapes de la lecture efficace]
\begin{enumerate}
\item \textbf{Lire le titre et anticiper} : le titre annonce le thème. Avant de lire, demande-toi : de quoi va parler ce texte ? Quels mots vais-je rencontrer ?
\item \textbf{Lecture globale silencieuse, sans dictionnaire} : lis tout le texte d'une traite pour saisir le sens général. Ne t'arrête pas sur les mots inconnus.
\item \textbf{Repérer les mots transparents} : beaucoup de mots allemands ressemblent au français ou à l'anglais : \all{der Supermarkt}, \all{das Radio}, \all{die Familie}, \all{der Preis}. Ils te donnent déjà la moitié du sens.
\item \textbf{Lecture détaillée avec le glossaire} : relis phrase par phrase, aide-toi du glossaire marginal et souligne le vocabulaire du thème.
\item \textbf{Reformuler} : raconte le texte avec tes propres mots, d'abord en français, puis en allemand avec des phrases simples.
\end{enumerate}
\end{methode}

\courssub{Première lecture : lecture globale}

Lis le texte suivant en silence, sans dictionnaire. Cherche seulement à répondre à trois questions : \emph{qui ? où ? quoi ?}

\begin{textebox}[Familie Berger geht einkaufen]
Am Samstagmorgen fährt Familie Berger in die Stadt. Zuerst gehen sie in den Supermarkt, denn die Lebensmittel sind dort billiger als auf dem Markt. Frau Berger vergleicht die Preise und sucht die besten Angebote. Sie kauft Obst, Gemüse, Brot und Milch. An der Kasse bezahlt sie mit Karte.

Der Vater kauft ein neues Radio im Elektrogeschäft : es ist teurer als das alte, aber viel besser. „Das neue Radio hat den besten Klang“, sagt der Verkäufer. Herr Berger ist ein zufriedener Kunde.

Am Nachmittag geht die Tochter zum Friseur — diese Dienstleistung kostet zwanzig Euro. Der Sohn geht zur Bank, denn er braucht Geld für den Ausflug der Schule. Am Abend sind alle müde, aber glücklich. Die Kunden sind zufrieden, und die Familie spart Geld für die Ferien. „Nächsten Samstag gehen wir auf den Markt“, sagt Frau Berger, „dort ist das Gemüse am frischesten !“
\end{textebox}

\textbf{Glossaire marginal}

\begin{center}
\begin{tabularx}{\linewidth}{|X|X|}
\hline
\rowcolor{popblueL} \textbf{Deutsch} & \textbf{Français} \\
\hline
\all{der Einkauf, die Einkäufe} & l'achat, les courses \\
\hline
\all{das Angebot, die Angebote} & l'offre, la promotion \\
\hline
\all{die Kasse, die Kassen} & la caisse \\
\hline
\all{der Kunde, die Kunden / die Kundin} & le client / la cliente \\
\hline
\all{die Dienstleistung, die Dienstleistungen} & la prestation de service \\
\hline
\all{die Ware, die Waren} & la marchandise \\
\hline
\all{der Preis, die Preise} & le prix \\
\hline
\all{das Geschäft, die Geschäfte} & le magasin, la boutique \\
\hline
\all{sparen} & économiser, épargner \\
\hline
\all{sich leisten} & se payer, s'offrir \\
\hline
\all{vergleichen (vergleicht, verglich, hat verglichen)} & comparer \\
\hline
\all{zufrieden} & satisfait, content \\
\hline
\end{tabularx}
\end{center}

\courssub{Repérage des informations globales : qui ? où ? quoi ?}

Après cette première lecture, complète mentalement le tableau d'identification du texte. C'est la première étape de tout commentaire de texte.

\begin{center}
\begin{tabularx}{\linewidth}{|X|X|}
\hline
\rowcolor{popgreenL} \textbf{Question} & \textbf{Réponse} \\
\hline
\all{Wer ?} (qui ?) & \all{Familie Berger : die Mutter, der Vater, die Tochter, der Sohn} \\
\hline
\all{Wo ?} (où ?) & \all{in der Stadt : im Supermarkt, im Elektrogeschäft, beim Friseur, bei der Bank} \\
\hline
\all{Wann ?} (quand ?) & \all{am Samstagmorgen, am Nachmittag, am Abend} \\
\hline
\all{Was ?} (quoi ?) & \all{die Familie kauft Lebensmittel und ein Radio, sie benutzt Dienstleistungen und spart Geld} \\
\hline
\end{tabularx}
\end{center}

\begin{definition}[Ware ou Dienstleistung ?]
\all{Die Ware} (pluriel \all{die Waren}) désigne une \textbf{marchandise}, un bien matériel que l'on peut toucher : le pain, le riz, un radio, un téléphone. \all{Die Dienstleistung} (pluriel \all{die Dienstleistungen}) désigne une \textbf{prestation de service} : le travail du coiffeur (\all{der Friseur}), de la banque (\all{die Bank}), du transport (\all{der Bus}, \all{das Taxi}). On achète une marchandise, mais on \emph{utilise} un service : \all{Ich kaufe ein Brot, aber ich gehe zum Friseur.}
\end{definition}

\begin{exemplebox}[Biens et services dans le texte]
\begin{itemize}
\item \all{Sie kauft Obst, Gemüse, Brot und Milch.} « Elle achète des fruits, des légumes, du pain et du lait. » (des \all{Waren})
\item \all{Der Vater kauft ein neues Radio.} « Le père achète un nouveau poste de radio. » (une \all{Ware})
\item \all{Die Tochter geht zum Friseur.} « La fille va chez le coiffeur. » (une \all{Dienstleistung})
\item \all{Der Sohn geht zur Bank.} « Le fils va à la banque. » (une \all{Dienstleistung})
\end{itemize}
\end{exemplebox}

\courssub{Deuxième lecture : lecture détaillée}

Relis maintenant le texte phrase par phrase avec le glossaire. Deux consignes de travail :

\begin{enumerate}
\item \textbf{Souligne} tout le vocabulaire du thème « consommation » : \all{der Supermarkt}, \all{die Lebensmittel}, \all{die Preise}, \all{die Angebote}, \all{die Kasse}, \all{bezahlen}, \all{der Kunde}, \all{die Dienstleistung}, \all{sparen}.
\item \textbf{Entoure} les adjectifs au comparatif et au superlatif : \all{billiger} (« moins cher »), \all{teurer} (« plus cher »), \all{besser} (« meilleur »), \all{am frischesten} (« le plus frais »), \all{die besten Angebote} (« les meilleures offres »), \all{den besten Klang} (« le meilleur son »).
\end{enumerate}

\begin{aretenir}[Observer avant d'apprendre la règle]
Observe : \all{billig} devient \all{billig\textbf{er}} (plus/moins cher), \all{teuer} devient \all{teur\textbf{er}}, \all{gut} devient \all{besser} et \all{am besten}. L'allemand forme le comparatif avec le suffixe \cle{-er} et le superlatif avec \cle{am \ldots -sten}, comme le français utilise « plus » et « le plus ». Nous étudierons ces formes en détail dans la partie grammaire de la leçon.
\end{aretenir}

\begin{attention}[Faux amis et pièges de lecture]
\begin{itemize}
\item \all{der Preis} signifie « le prix », mais aussi « le prix » au sens de récompense ; ne jamais traduire par « la presse » (\all{die Presse}).
\item \all{das Angebot} est « l'offre, la promotion », pas « l'ange ».
\item \all{bekommen} signifie « recevoir », jamais « devenir » (\all{werden}).
\item \all{der Kunde} est masculin : « le client » ; la cliente est \all{die Kundin}. Attention à l'article !
\item \all{billig} = « bon marché » ; ne pas confondre avec \all{billig} au sens péjoratif « de mauvaise qualité » selon le contexte.
\end{itemize}
\end{attention}

\courssub{Carte mentale du vocabulaire : Der Konsum}

\begin{popfigure}
\begin{tikzpicture}[
  mindmap,
  grow cyclic,
  every node/.style={concept, align=center},
  root concept/.append style={concept color=popblue, font=\large\bfseries, text=white},
  level 1/.append style={level distance=3.4cm, sibling angle=90},
  level 1 concept/.append style={concept color=poporange, text=white, font=\small\bfseries},
  level 2/.append style={level distance=2.4cm},
  level 2 concept/.append style={concept color=popgreen!70, text=white, font=\scriptsize}
]
\node[root concept]{Der Konsum}
  child { node {Waren kaufen}
    child { node {die Ware, -n} }
    child { node {der Einkauf, \"-e} }
    child { node {die Kasse, -n} }
  }
  child { node {Dienstleistungen}
    child { node {der Friseur} }
    child { node {die Bank} }
    child { node {der Transport} }
  }
  child { node {Preise vergleichen}
    child { node {der Preis, -e} }
    child { node {billiger als} }
    child { node {das Angebot, -e} }
  }
  child { node {Gesch\"afte}
    child { node {der Supermarkt} }
    child { node {der Markt} }
    child { node {das Elektro-gesch\"aft} }
  };
\end{tikzpicture}
\legende{Carte mentale : le champ lexical de la consommation}
\end{popfigure}

\courssub{Reformulation orale : die Zusammenfassung}

\begin{methode}[Résumer un texte en quatre phrases]
\begin{enumerate}
\item \textbf{Qui et quand} : \all{Am Samstag fährt Familie Berger in die Stadt.}
\item \textbf{Action principale} : \all{Sie kaufen Lebensmittel im Supermarkt und ein Radio im Elektrogeschäft.}
\item \textbf{Actions secondaires} : \all{Die Tochter geht zum Friseur, der Sohn geht zur Bank.}
\item \textbf{Conclusion} : \all{Die Familie ist zufrieden und spart Geld für die Ferien.}
\end{enumerate}
Règle d'or : une phrase = une idée ; utilise tes propres mots, pas de phrases copiées du texte.
\end{methode}

\begin{exoresolu}[Questions de compréhension]
Réponds en allemand par des phrases complètes :
\begin{enumerate}
\item \all{Wohin fährt Familie Berger am Samstagmorgen ?}
\item \all{Warum geht die Familie in den Supermarkt ?}
\item \all{Was kauft der Vater im Elektrogeschäft ?}
\item \all{Wie viel kostet die Dienstleistung beim Friseur ?}
\item \all{Wofür spart die Familie Geld ?}
\end{enumerate}
\tcblower
\textbf{Solutions :}
\begin{enumerate}
\item \all{Familie Berger fährt am Samstagmorgen in die Stadt.} « La famille Berger va en ville samedi matin. »
\item \all{Sie geht in den Supermarkt, denn die Lebensmittel sind dort billiger als auf dem Markt.} « Elle va au supermarché parce que les produits alimentaires y sont moins chers qu'au marché. »
\item \all{Der Vater kauft ein neues Radio.} « Le père achète un nouveau poste de radio. »
\item \all{Die Dienstleistung beim Friseur kostet zwanzig Euro.} « La prestation chez le coiffeur coûte vingt euros. »
\item \all{Die Familie spart Geld für die Ferien.} « La famille économise de l'argent pour les vacances. »
\end{enumerate}
\end{exoresolu}

\begin{exercice}[À ton tour]
\begin{enumerate}
\item Classe les mots suivants en deux colonnes \all{Waren} / \all{Dienstleistungen} : \all{das Brot, der Friseur, das Radio, die Bank, das Gemüse, der Transport, die Milch}.
\item Traduis en français : \all{Das neue Radio ist teurer als das alte, aber viel besser.}
\item Résume le texte en trois phrases allemandes avec tes propres mots.
\end{enumerate}
\end{exercice}

\begin{corrige}[Exercice « À ton tour »]
\begin{enumerate}
\item \all{Waren} : \all{das Brot, das Radio, das Gemüse, die Milch}. \all{Dienstleistungen} : \all{der Friseur, die Bank, der Transport}.
\item « Le nouveau poste de radio est plus cher que l'ancien, mais bien meilleur. »
\item Exemple : \all{Familie Berger macht am Samstag Einkäufe in der Stadt. Sie kauft Lebensmittel und ein Radio. Alle sind zufrieden und sparen für die Ferien.}
\end{enumerate}
\end{corrige}

\begin{savaistu}[Le dimanche, tout est fermé !]
En Allemagne, les magasins sont fermés le dimanche : c'est la loi du \all{Ladenschluss} (littéralement « fermeture des magasins »). Le samedi est donc le grand jour des courses, comme dans notre texte ! Seuls les boulangers, les stations-service et quelques commerces de gare peuvent ouvrir. Au Cameroun, c'est l'inverse : les marchés comme Mokolo ou le marché central de Douala sont animés presque tous les jours de la semaine.
\end{savaistu}

\coursec{Compréhension du texte}

Dans la section précédente, nous avons lu et découvert le texte support \emph{Einkaufen in Deutschland}, qui raconte la journée de courses de la famille Kamdem à Yaoundé et compare ses habitudes de consommation avec celles d'une famille allemande. Nous allons maintenant travailler ce texte en profondeur : d'abord une compréhension \cle{globale} (le sens général), puis une compréhension \cle{détaillée} (les informations précises), et enfin une compréhension \cle{interprétative} (ce que le texte suggère). C'est exactement la démarche exigée à l'examen : lire, repérer, citer, rédiger.

\courssub{Rappel du texte support}

Avant de répondre aux questions, relisons le texte. Le voici à nouveau, accompagné de son glossaire.

\begin{textebox}[Einkaufen in Deutschland — und in Kamerun]
Am Samstagmorgen fährt die Familie Kamdem mit dem Auto in die Stadt. Zuerst gehen sie in den Supermarkt, denn die Lebensmittel sind dort billiger als auf dem Markt. Die Mutter kauft Reis, Brot, Milch und Gemüse. Der Vater braucht ein neues Radio, denn das alte Radio ist kaputt. Das neue Radio ist teurer als das alte, aber es hat eine bessere Qualität. „Wir müssen die Preise vergleichen“, sagt der Vater, „denn wir wollen Geld sparen.“

Am Nachmittag geht der Sohn Junior zum Friseur. Der Haarschnitt kostet 1500 Franken. Das ist billig, findet Junior. Seine Schwester Marie kauft in einer Boutique ein schönes Kleid für die Hochzeit ihrer Cousine. Das Kleid ist schöner als die anderen Kleider, aber auch am teuersten.

In Deutschland machen viele Familien ihre Einkäufe auch am Samstag. Sie gehen in Supermärkte wie Aldi oder Lidl, denn dort sind die Waren oft am billigsten. Deutsche Kunden vergleichen die Preise im Internet, bevor sie etwas kaufen. Auch Dienstleistungen sind wichtig : der Friseur, der Arzt, die Reparatur des Autos. Diese Dienstleistungen sind in Deutschland viel teurer als in Kamerun. Ein Haarschnitt kostet dort oft zwanzig Euro oder mehr.

Am Abend ist die Familie Kamdem müde, aber zufrieden. Sie hat alles gekauft und auch ein bisschen Geld gespart. „Nächste Woche gehen wir wieder einkaufen“, sagt die Mutter und lacht.
\end{textebox}

\begin{definition}[Glossaire du texte]
\begin{itemize}
\item \all{der Einkauf, die Einkäufe} : l'achat, les courses ; \all{einkaufen} : faire les courses
\item \all{die Lebensmittel} (pluriel) : les denrées alimentaires
\item \all{der Friseur, die Friseure} : le coiffeur ; \all{der Haarschnitt, die Haarschnitte} : la coupe de cheveux
\item \all{die Ware, die Waren} : la marchandise ; \all{die Dienstleistung, die Dienstleistungen} : la prestation de service
\item \all{der Preis, die Preise} : le prix ; \all{die Preise vergleichen} : comparer les prix
\item \all{Geld sparen} : économiser de l'argent ; \all{kaputt} : cassé, en panne
\item \all{die Qualität, die Qualitäten} : la qualité ; \all{zufrieden} : satisfait
\item \all{bevor} : avant de (+ infinitif) ; \all{zuerst} : d'abord
\end{itemize}
\end{definition}

\courssub{La méthode de compréhension : Lire, Repérer, Citer, Rédiger}

Pour réussir les questions de compréhension à l'examen, il ne suffit pas d'avoir compris le texte : il faut le \emph{prouver} par une réponse bien construite. Voici la démarche en quatre étapes.

\begin{popfigure}
\begin{tikzpicture}[
  node distance=0.9cm,
  etape/.style={rectangle, rounded corners, draw=popblue, fill=popblueL, thick, minimum width=3.1cm, minimum height=1.1cm, align=center, font=\bfseries},
  desc/.style={rectangle, draw=popline, fill=popcream, rounded corners, align=center, minimum width=3.1cm, minimum height=1.4cm, font=\small},
  fleche/.style={-{Stealth[length=3mm]}, very thick, poporange}
]
\node[etape] (lire) {1. LIRE};
\node[desc, below=0.35cm of lire, align=center] (dlire){Lire le texte deux fois :\\ sens global, puis détails};
\node[etape, right=1.2cm of lire] (reperer) {2. REPÉRER};
\node[desc, below=0.35cm of reperer, align=center] (dreperer){Souligner dans le texte\\ les mots de la question};
\node[etape, right=1.2cm of reperer] (citer) {3. CITER};
\node[desc, below=0.35cm of citer, align=center] (dciter){Reprendre l'élément\\ exact du texte};
\node[etape, right=1.2cm of citer] (rediger) {4. RÉDIGER};
\node[desc, below=0.35cm of rediger, align=center] (drediger){Phrase complète :\\ verbe conjugué en 2e position};
\draw[fleche] (lire) -- (reperer);
\draw[fleche] (reperer) -- (citer);
\draw[fleche] (citer) -- (rediger);
\end{tikzpicture}
\legende{Organigramme de la méthode de compréhension : Lire $\rightarrow$ Repérer $\rightarrow$ Citer $\rightarrow$ Rédiger}
\end{popfigure}

\begin{methode}[Comment répondre à une question de compréhension]
\begin{enumerate}
\item \textbf{Lire la question} et identifier le mot interrogatif : \all{wohin} (où, direction), \all{was} (quoi), \all{warum} (pourquoi), \all{wie viel} (combien), \all{wer} (qui), \all{wann} (quand).
\item \textbf{Repérer} dans le texte la phrase qui contient la réponse : les mots de la question se retrouvent presque toujours dans le texte (\all{fährt}, \all{Supermarkt}, \all{Friseur}...).
\item \textbf{Citer} mentalement l'élément exact du texte qui répond à la question.
\item \textbf{Rédiger} une phrase complète en allemand : reprendre les mots de la question, ajouter l'élément cité, conjuguer le verbe et le placer en \cle{deuxième position}. Si la question commence par \all{warum}, la réponse commence par \all{weil} (et le verbe conjugué va à la fin).
\end{enumerate}
\end{methode}

\begin{exemplebox}[Une question, une réponse modèle]
\all{Warum gehen sie in den Supermarkt?} — \all{Weil die Lebensmittel dort billiger sind.}

« Pourquoi vont-ils au supermarché ? — Parce que les denrées y sont moins chères. »

Observez : la réponse reprend les mots de la question (\all{sie}, \all{Supermarkt}), cite l'information du texte (\all{die Lebensmittel sind dort billiger}) et forme une phrase complète. Avec \all{weil}, le verbe conjugué (\all{sind}) se place à la fin de la phrase.
\end{exemplebox}

\begin{attention}[Jamais de réponse en un seul mot !]
Au Bac comme à l'examen de Première, une réponse comme \all{In die Stadt.} ou \all{1500 Franken.} est \textbf{sanctionnée}, même si l'information est juste. Il faut toujours rédiger une \cle{phrase complète} : sujet + verbe conjugué en deuxième position + complément.

Faux (incomplet) : \all{In den Supermarkt.}

Correct : \all{Sie gehen in den Supermarkt.} « Ils vont au supermarché. »

Autre piège : après \all{weil}, le verbe conjugué se place \textbf{à la fin} : \all{Weil die Lebensmittel dort billiger \textbf{sind}.} et non \all{weil die Lebensmittel sind billiger}.
\end{attention}

\courssub{Compréhension globale : le sens général du texte}

Les questions globales portent sur les grandes informations : les personnages, le lieu, l'action principale. Elles se résolvent après une première lecture attentive.

\begin{exoresolu}[Questions de compréhension globale]
\textbf{Énoncé.} Répondez par des phrases complètes.

1. \all{Wohin fährt die Familie Kamdem?}

2. \all{Was kaufen sie?}

3. \all{Wann machen die Familien ihre Einkäufe?}
\tcblower
\textbf{Solution détaillée.}

1. Le mot interrogatif est \all{wohin} (vers quel lieu). Le texte dit : \all{fährt die Familie Kamdem mit dem Auto in die Stadt}. On reprend les mots de la question et on cite le lieu :

\all{Die Familie Kamdem fährt in die Stadt.} « La famille Kamdem va en ville. »

2. Le mot interrogatif est \all{was} (quoi). Le texte énumère les achats : \all{Reis, Brot, Milch und Gemüse}, puis \all{ein neues Radio}, \all{ein schönes Kleid}. On rédige :

\all{Sie kaufen Lebensmittel, ein neues Radio und ein schönes Kleid.} « Ils achètent des denrées alimentaires, une nouvelle radio et une belle robe. »

3. Le mot interrogatif est \all{wann} (quand). Le texte dit : \all{Am Samstagmorgen fährt die Familie...} et \all{viele Familien machen ihre Einkäufe auch am Samstag}.

\all{Die Familien machen ihre Einkäufe am Samstag.} « Les familles font leurs courses le samedi. »

Remarquez la place du verbe conjugué : toujours en deuxième position (\all{fährt}, \all{kaufen}, \all{machen}).
\end{exoresolu}

\courssub{Compréhension détaillée : les informations précises}

Les questions détaillées exigent de retrouver un élément précis : une cause, une comparaison, un prix. Il faut relire le passage concerné et citer exactement.

\begin{exoresolu}[Questions de compréhension détaillée]
\textbf{Énoncé.} Répondez par des phrases complètes.

1. \all{Warum gehen sie in den Supermarkt?}

2. \all{Was ist teurer als das alte Radio?}

3. \all{Wie viel kostet der Friseur?}

4. \all{Wo vergleichen die deutschen Kunden die Preise?}
\tcblower
\textbf{Solution détaillée.}

1. \all{Warum} appelle une réponse avec \all{weil}. Le texte dit : \all{denn die Lebensmittel sind dort billiger als auf dem Markt}.

\all{Sie gehen in den Supermarkt, weil die Lebensmittel dort billiger sind.} « Ils vont au supermarché parce que les denrées y sont moins chères. »

2. Le texte dit : \all{Das neue Radio ist teurer als das alte.} On reprend la structure de la question :

\all{Das neue Radio ist teurer als das alte Radio.} « La nouvelle radio est plus chère que l'ancienne. »

3. \all{Wie viel} demande un prix. Le texte dit : \all{Der Haarschnitt kostet 1500 Franken.}

\all{Der Friseur kostet 1500 Franken.} « Le coiffeur coûte 1500 francs. »

4. Le texte dit : \all{Deutsche Kunden vergleichen die Preise im Internet.}

\all{Die deutschen Kunden vergleichen die Preise im Internet.} « Les clients allemands comparent les prix sur Internet. »
\end{exoresolu}

\begin{aretenir}[Les mots interrogatifs et le type de réponse attendu]
\begin{center}
\begin{tabularx}{\linewidth}{|l|l|X|}
\hline
\rowcolor{popblueL} \textbf{Mot interrogatif} & \textbf{Français} & \textbf{Réponse attendue} \\
\hline
\all{Wohin?} & vers où ? & un lieu avec \all{in}, \all{nach}, \all{zu} \\
\hline
\all{Was?} & quoi ? & un objet ou une action \\
\hline
\all{Warum?} & pourquoi ? & une cause avec \all{weil} (verbe à la fin) \\
\hline
\all{Wie viel?} & combien ? & un prix ou une quantité \\
\hline
\all{Wer?} & qui ? & une personne \\
\hline
\all{Wann?} & quand ? & un moment (\all{am Samstag}, \all{am Abend}) \\
\hline
\all{Wo?} & où ? (position) & un lieu avec \all{im}, \all{auf dem} \\
\hline
\end{tabularx}
\end{center}
\end{aretenir}

\courssub{Compréhension interprétative : aller plus loin que le texte}

Les questions d'interprétation demandent de \emph{comprendre le sens profond} : une intention, une expression, une attitude. La réponse n'est pas écrite mot pour mot dans le texte ; il faut la déduire et l'expliquer avec ses propres mots, en s'appuyant sur le texte.

\textbf{Question 1 : Pourquoi la famille compare-t-elle les prix ?}

Le père dit : \all{Wir müssen die Preise vergleichen, denn wir wollen Geld sparen.} Comparer les prix, c'est regarder le prix du même produit dans plusieurs magasins pour choisir le moins cher. L'intention de la famille est donc d'acheter au meilleur prix et de ne pas dépenser trop.

Réponse rédigée : \all{Die Familie vergleicht die Preise, weil sie nicht zu viel Geld ausgeben will. Sie möchte die billigsten Waren kaufen.} « La famille compare les prix parce qu'elle ne veut pas dépenser trop d'argent. Elle veut acheter les marchandises les moins chères. »

\textbf{Question 2 : Que signifie „Geld sparen“ ?}

\all{Geld sparen} signifie littéralement « économiser de l'argent » : ne pas dépenser tout son argent, garder une partie pour plus tard. Dans le texte, la famille achète des produits moins chers et, à la fin, \all{sie hat ein bisschen Geld gespart} : il lui reste un peu d'argent.

Réponse rédigée : \all{„Geld sparen“ bedeutet : nicht alles Geld ausgeben und einen Teil für später behalten.} « „Geld sparen“ signifie : ne pas dépenser tout son argent et en garder une partie pour plus tard. »

\begin{attention}[Interpréter ne veut pas dire inventer]
Même pour une question d'interprétation, votre réponse doit \textbf{s'appuyer sur le texte}. On ne peut pas écrire que la famille est pauvre : le texte ne le dit pas ; il dit seulement qu'elle veut économiser. Citez toujours l'élément du texte qui justifie votre interprétation (\all{denn wir wollen Geld sparen}), puis expliquez-le avec vos mots.
\end{attention}

\courssub{Vrai ou faux : Richtig oder Falsch ?}

L'exercice \cle{Richtig/Falsch} est classique à l'examen. Pour chaque affirmation, vous devez écrire \all{richtig} ou \all{falsch} et \textbf{justifier par une citation du texte} entre guillemets allemands „…“. Une réponse sans justification ne vaut rien.

\begin{exoresolu}[Richtig oder Falsch ? Justifiez par une citation]
\textbf{Énoncé.} Dites si les phrases sont vraies (\all{richtig}) ou fausses (\all{falsch}) et justifiez par une citation du texte.

1. \all{Die Familie Kamdem fährt mit dem Bus in die Stadt.}

2. \all{Das neue Radio ist billiger als das alte Radio.}

3. \all{Marie kauft ein Kleid für eine Hochzeit.}

4. \all{Ein Haarschnitt ist in Deutschland billiger als in Kamerun.}

5. \all{Die Familie ist am Abend zufrieden.}
\tcblower
\textbf{Solution détaillée.}

1. \all{Falsch.} Le texte dit : „\all{fährt die Familie Kamdem mit dem Auto in die Stadt}“. Elle prend la voiture, pas le bus.

2. \all{Falsch.} Le texte dit : „\all{Das neue Radio ist teurer als das alte}“. La nouvelle radio est plus chère, pas moins chère. Attention au comparatif \all{teurer} (plus cher) qui est le contraire de \all{billiger} (moins cher).

3. \all{Richtig.} Le texte dit : „\all{Marie kauft in einer Boutique ein schönes Kleid für die Hochzeit ihrer Cousine}“.

4. \all{Falsch.} Le texte dit : „\all{Diese Dienstleistungen sind in Deutschland viel teurer als in Kamerun}“. Le coiffeur est plus cher en Allemagne (vingt euros ou plus contre 1500 francs à Yaoundé).

5. \all{Richtig.} Le texte dit : „\all{Am Abend ist die Familie Kamdem müde, aber zufrieden}“.
\end{exoresolu}

\begin{methode}[Réussir l'exercice Richtig/Falsch]
\begin{enumerate}
\item Lire l'affirmation et souligner ses mots-clés (\all{Bus}, \all{billiger}, \all{Hochzeit}).
\item Chercher ces mots-clés (ou leurs contraires) dans le texte.
\item Comparer attentivement : l'affirmation dit-elle exactement la même chose que le texte ? Méfiez-vous des contraires (\all{billiger} / \all{teurer}) et des petits changements (\all{Bus} / \all{Auto}).
\item Écrire \all{richtig} ou \all{falsch}, puis recopier entre guillemets la phrase du texte qui prouve votre réponse.
\end{enumerate}
\end{methode}

\courssub{Entraînement à l'expression écrite guidée}

Pour terminer, rédigeons des réponses complètes et variées. C'est le passage de la compréhension à la production : vous réutilisez le vocabulaire et les structures du texte pour écrire vos propres phrases.

\begin{exercice}[À vous de répondre]
Répondez par des phrases complètes en allemand.

1. \all{Was kauft die Mutter im Supermarkt?}

2. \all{Warum braucht der Vater ein neues Radio?}

3. \all{Wie findet Junior den Haarschnitt?}

4. \all{Wie ist das Kleid von Marie?} (utilisez le comparatif et le superlatif)

5. \all{Was machen die deutschen Kunden, bevor sie etwas kaufen?}

6. \all{Und Sie : Kaufen Sie gern auf dem Markt oder im Supermarkt ein? Warum?} (2 phrases)
\end{exercice}

\begin{corrige}[Exercice : expression écrite guidée]
1. \all{Die Mutter kauft Reis, Brot, Milch und Gemüse.} « La mère achète du riz, du pain, du lait et des légumes. »

2. \all{Der Vater braucht ein neues Radio, weil das alte Radio kaputt ist.} « Le père a besoin d'une nouvelle radio parce que l'ancienne est cassée. » (Verbe \all{ist} à la fin après \all{weil}.)

3. \all{Junior findet den Haarschnitt billig.} « Junior trouve la coupe de cheveux bon marché. »

4. \all{Das Kleid von Marie ist schöner als die anderen Kleider, aber es ist am teuersten.} « La robe de Marie est plus belle que les autres robes, mais c'est la plus chère. »

5. \all{Bevor sie etwas kaufen, vergleichen die deutschen Kunden die Preise im Internet.} « Avant d'acheter quelque chose, les clients allemands comparent les prix sur Internet. » (Après \all{bevor}, le verbe conjugué \all{kaufen} est à la fin de la subordonnée ; quand la subordonnée ouvre la phrase, le verbe de la principale, \all{vergleichen}, vient juste après la virgule.)

6. Modèle de réponse personnelle : \all{Ich kaufe gern auf dem Markt ein, weil die Preise dort billiger sind. Ich kaufe auch gern im Supermarkt ein, weil man dort alles findet.} « J'aime faire mes courses au marché parce que les prix y sont moins chers. J'aime aussi faire mes courses au supermarché parce qu'on y trouve tout. » (Attention au verbe à particule : \all{einkaufen} se sépare — \all{ich kaufe ... ein}.)
\end{corrige}

\begin{aretenir}[Les cinq règles d'or de la question de compréhension]
\begin{itemize}
\item Toujours une \textbf{phrase complète}, jamais un mot isolé.
\item Le verbe conjugué en \textbf{deuxième position} (sauf après \all{weil} : verbe à la fin).
\item Reprendre les \textbf{mots de la question} dans la réponse.
\item \textbf{Citer le texte} pour justifier (surtout au Richtig/Falsch).
\item Vérifier le sens : attention aux \textbf{contraires} (\all{billig/teuer}, \all{alt/neu}) et aux petits détails changés.
\end{itemize}
\end{aretenir}

Nous savons désormais lire le texte de trois façons : globalement, en détail et en interprétation. La section suivante nous permettra d'enrichir notre lexique : le vocabulaire thématique de la consommation et des services, \emph{Konsum und Dienstleistungen}.

\coursec{Lexique thématique : Konsum und Dienstleistungen}

Dans la section précédente, nous avons lu et compris le texte \all{Einkaufen in Deutschland}. Nous y avons rencontré des mots comme \all{der Preis}, \all{das Geschäft}, \all{der Kunde} ou \all{die Dienstleistung}. Il est temps maintenant de nous arrêter sur ce vocabulaire : c'est la condition indispensable pour parler de la consommation, comprendre un texte économique et réussir les questions de compréhension comme les productions écrites. En allemand, un nom ne s'apprend jamais seul : on l'apprend toujours avec son \cle{article} (\all{der}, \all{die}, \all{das}) et son \cle{pluriel}.

\courssub{Observer d'abord : un dialogue au magasin}

\begin{textebox}[Im Sportgeschäft — ein Dialog]
\textbf{Verkäufer:} Guten Tag! Kann ich Ihnen helfen?

\textbf{Kunde:} Ja, bitte. Ich suche neue Sportschuhe. Wie viel kostet dieses Paar?

\textbf{Verkäufer:} Diese Schuhe kosten 89 Euro. Aber wir haben heute ein Angebot: Sie bekommen 20 Prozent Rabatt.

\textbf{Kunde:} Das ist günstig! Und die Sportschuhe dort im Regal?

\textbf{Verkäufer:} Die sind billiger, sie kosten nur 45 Euro. Aber die Qualität ist nicht so gut.

\textbf{Kunde:} Hmm... Ich nehme die teuren Schuhe. Qualität ist wichtig. Kann ich mit Karte bezahlen?

\textbf{Verkäufer:} Natürlich. Bitte zahlen Sie an der Kasse. Hier ist Ihre Rechnung.

\textbf{Kunde:} Vielen Dank! Auf Wiedersehen!

\textbf{Verkäufer:} Auf Wiedersehen, und viel Spaß mit den neuen Schuhen!
\end{textebox}

\textbf{Glossaire (avec prononciation approximative) :}
\all{der Verkäufer} (-) [fèr-kèu-fèr] : le vendeur ; \all{das Paar} (-e) [paːɐ̯] : la paire ; \all{das Angebot} (-e) [an-ge-bot] : l'offre, la promotion ; \all{der Rabatt} (-e) [ra-bat] : la remise ; \all{günstig} [gyns-tiç] : avantageux, bon marché ; \all{das Regal} (-e) [re-gahl] : le rayon, l'étagère ; \all{die Qualität} (-en) [kwa-li-tɛːt] : la qualité ; \all{die Kasse} (-n) [kas-sə] : la caisse ; \all{die Rechnung} (-en) [rɛç-nuŋ] : la facture, l'addition.

\textbf{Questions de compréhension :}
\begin{enumerate}
\item \all{Was sucht der Kunde?} (Que cherche le client ?)
\item \all{Wie viel Rabatt bekommt er?} (Quelle remise obtient-il ?)
\item \all{Warum nimmt er die teuren Schuhe?} (Pourquoi prend-il les chaussures chères ?)
\item \all{Wo bezahlt er?} (Où paie-t-il ?)
\end{enumerate}

\courssub{Champ lexical 1 : Acheter — \all{kaufen, einkaufen, bestellen, bezahlen, sich leisten, sparen, ausgeben}}

Observons d'abord ces verbes dans des phrases :

\all{Ich kaufe ein Brot.} « J'achète un pain. »

\all{Am Samstag kaufen wir auf dem Markt ein.} « Samedi, nous faisons les courses au marché. »

\all{Sie bestellt ein Buch im Internet.} « Elle commande un livre sur Internet. »

\all{Er bezahlt die Rechnung an der Kasse.} « Il paie la facture à la caisse. »

\all{Ich kann mir kein Auto leisten.} « Je ne peux pas me permettre (m'offrir) une voiture. »

\all{Die Schüler sparen Geld für eine Reise.} « Les élèves économisent de l'argent pour un voyage. »

\all{Er gibt viel Geld für Kleidung aus.} « Il dépense beaucoup d'argent pour les vêtements. »

\begin{center}
\begin{tabularx}{\linewidth}{|l|X|X|l|}
\hline
\rowcolor{popblueL} \textbf{Verbe} & \textbf{Français} & \textbf{Remarque} & \textbf{Prononciation} \\
\hline
\all{kaufen} & acheter & verbe régulier : \all{ich kaufe, er kauft} & « kow-fen » \\
\hline
\all{einkaufen} & faire les courses & particule séparable : \all{ich kaufe ein} & « ayn-kow-fen » \\
\hline
\all{bestellen} & commander & régulier : \all{sie bestellt} & « be-chtèl-len » \\
\hline
\all{bezahlen} & payer & régulier : \all{er bezahlt} & « be-tsah-len » \\
\hline
\all{sich leisten} & se permettre, s'offrir & verbe pronominal : \all{ich leiste mir} & « lay-sten » \\
\hline
\all{sparen} & économiser, épargner & régulier : \all{wir sparen} & « chpa-ren » \\
\hline
\all{ausgeben} & dépenser & fort et séparable : \all{ich gebe aus, er gibt aus} & « ows-guh-ben » \\
\hline
\end{tabularx}
\end{center}

\begin{attention}[Verbes forts et séparables dans ce champ]
Deux pièges pour les francophones :
\begin{itemize}
\item \all{einkaufen} et \all{ausgeben} sont à \cle{particule séparable} : la particule va en fin de phrase (position de clôture). \all{Ich kaufe heute ein.} / \all{Er gibt sein Geld aus.}
\item \all{ausgeben} est en plus un verbe fort : voyelle radicale \all{e → i} au du/sg. (\all{geben → er gibt}), Participe passé : \all{ausgegeben} (Perfekt : \all{hat ausgegeben}). On dit toujours \all{Geld ausgeben} (dépenser de l'argent), jamais \all{Geld spendieren} dans ce sens au niveau A2 (\all{spendieren} = offrir un verre/un repas à qqn).
\end{itemize}
\end{attention}

\courssub{Champ lexical 2 : Prix et argent — \all{der Preis, das Angebot, die Rechnung, der Rabatt}}

\begin{center}
\begin{tabularx}{\linewidth}{|l|l|X|l|}
\hline
\rowcolor{popblueL} \textbf{Deutsch} & \textbf{Pluriel} & \textbf{Français} & \textbf{Prononciation} \\
\hline
\all{der Preis} & \all{die Preise} & le prix & « prays » \\
\hline
\all{das Angebot} & \all{die Angebote} & l'offre, la promotion & « an-ge-bot » \\
\hline
\all{die Rechnung} & \all{die Rechnungen} & la facture, l'addition & « rèkh-noung » \\
\hline
\all{der Rabatt} & \all{die Rabatte} & la remise, la réduction & « ra-bat » \\
\hline
\all{das Geld} & \all{die Gelder} (rare) & l'argent & « guèlt » \\
\hline
\all{billig} & — & bon marché & « bi-lig » \\
\hline
\all{teuer} & — & cher, chère & « toy-er » \\
\hline
\all{günstig} & — & avantageux, à bon prix & « gyns-tiç » \\
\hline
\end{tabularx}
\end{center}

\begin{exemplebox}[Exemples traduits]
\all{Der Preis ist zu hoch.} « Le prix est trop élevé. »

\all{Der Supermarkt hat heute ein gutes Angebot.} « Le supermarché a une bonne promotion aujourd'hui. »

\all{Die Rechnung bitte!} « L'addition, s'il vous plaît ! » (au restaurant)

\all{Schüler bekommen oft einen Rabatt.} « Les élèves obtiennent souvent une réduction. »

\all{Diese Ware ist billig, aber nicht gut.} « Cette marchandise est bon marché, mais pas bonne. »
\end{exemplebox}

\begin{attention}[\all{billig} ou \all{günstig} ?]
\all{billig} signifie « bon marché », parfois avec une nuance négative (de mauvaise qualité). \all{günstig} est toujours positif : « avantageux, intéressant ». Un prix \all{günstig} est une bonne affaire ; un produit \all{billig} peut être de la camelote. Attention aussi au faux ami : \all{die Rechnung} n'est pas seulement « l'addition » au restaurant, c'est aussi « la facture » (électricité, téléphone, achat en ligne).
\end{attention}

\courssub{Champ lexical 3 : Lieux et personnes — \all{das Geschäft, der Laden, der Supermarkt, der Markt, der Kunde}}

\begin{center}
\begin{tabularx}{\linewidth}{|l|l|X|l|}
\hline
\rowcolor{popblueL} \textbf{Deutsch} & \textbf{Pluriel} & \textbf{Français} & \textbf{Prononciation} \\
\hline
\all{das Geschäft} & \all{die Geschäfte} & le magasin, le commerce & « ge-chèft » \\
\hline
\all{der Laden} & \all{die Läden} & le magasin, la boutique & « lah-den » / « lè-den » \\
\hline
\all{der Supermarkt} & \all{die Supermärkte} & le supermarché & « zou-per-markt » \\
\hline
\all{der Markt} & \all{die Märkte} & le marché & « markt » / « mèrk-te » \\
\hline
\all{der Kunde} & \all{die Kunden} & le client & « koun-de » \\
\hline
\all{die Kundin} & \all{die Kundinnen} & la cliente & « koun-din » \\
\hline
\all{der Verkäufer} & \all{die Verkäufer} & le vendeur & « fèr-kèu-fèr » \\
\hline
\all{die Verkäuferin} & \all{die Verkäuferinnen} & la vendeuse & « fèr-kèu-fèr-in » \\
\hline
\all{die Kasse} & \all{die Kassen} & la caisse & « kas-se » \\
\hline
\end{tabularx}
\end{center}

\begin{attention}[\all{der Kunde} : un nom faible (N-Deklination) !]
\all{der Kunde} appartient à la \cle{déclinaison faible} (noms en \textbf{-e} désignant des personnes/animaux masculins) : il prend un \textbf{-n} à tous les cas \textbf{sauf au nominatif singulier}. Voici sa déclinaison complète au singulier :
\begin{center}
\begin{tabular}{|l|l|l|}
\hline
\rowcolor{popblueL} \textbf{Cas} & \textbf{Forme} & \textbf{Exemple avec verbe régissant le cas} \\
\hline
Nominatif & \all{der Kunde} & \all{Der Kunde kauft ein Brot.} (Sujet) \\
\hline
Accusatif & \all{den Kunden} & \all{Ich sehe den Kunden.} (Objet direct / \all{sehen}) \\
\hline
Datif & \all{dem Kunden} & \all{Ich helfe dem Kunden.} (Objet indirect / \all{helfen}) \\
\hline
Génitif & \all{des Kunden} & \all{Das Auto des Kunden ist neu.} (Possession) \\
\hline
\end{tabular}
\end{center}
Au pluriel, la forme \textbf{-n} s'ajoute aussi au datif : \all{den Kunden} (datif pluriel). \textbf{Ne dites jamais \all{den Kunde} ni \all{dem Kunde} !}
\end{attention}

\courssub{Champ lexical 4 : Biens et services — \all{die Ware, das Produkt, die Dienstleistung}}

\begin{center}
\begin{tabularx}{\linewidth}{|l|l|X|l|}
\hline
\rowcolor{popblueL} \textbf{Deutsch} & \textbf{Pluriel} & \textbf{Français} & \textbf{Prononciation} \\
\hline
\all{die Ware} & \all{die Waren} & la marchandise & « va-re » \\
\hline
\all{das Produkt} & \all{die Produkte} & le produit & « pro-dukt » \\
\hline
\all{die Dienstleistung} & \all{die Dienstleistungen} & la prestation de service & « dinst-leis-toung » \\
\hline
\all{der Friseur} & \all{die Friseure} & le coiffeur & « fri-zœr » \\
\hline
\all{die Bank} & \all{die Banken} & la banque & « bank » \\
\hline
\all{der Transport} & \all{die Transporte} & le transport & « trans-port » \\
\hline
\end{tabularx}
\end{center}

\begin{definition}[Bien ou service ?]
Un \cle{bien} (\all{die Ware}, \all{das Produkt}) est un objet matériel qu'on peut toucher : du pain, des chaussures, un téléphone. Un \cle{service} (\all{die Dienstleistung}) est un travail fourni par une personne : la coupe de cheveux chez le \all{Friseur}, le transport en bus, les services de la \all{Bank}.
\end{definition}

\begin{exemplebox}[Exemples traduits]
\all{Der Kunde kauft eine Ware im Geschäft.} « Le client achète une marchandise dans le magasin. »

\all{Die Dienstleistung ist teuer.} « La prestation de service est chère. »

\all{In Douala benutzen viele Menschen den Transport mit dem Motorrad.} « À Douala, beaucoup de gens utilisent le transport en moto (bend-skin). »

\all{Mein Vater geht heute zur Bank.} « Mon père va à la banque aujourd'hui. » (\all{zur} = \all{zu der})
\end{exemplebox}

\courssub{Les familles de mots : construire son vocabulaire}

\begin{definition}[Famille de mots]
Une \cle{famille de mots} (\all{die Wortfamilie}) est un ensemble de mots construits sur une même racine. Connaître la racine permet de deviner le sens de mots nouveaux. Exemple : \all{kaufen} (acheter) → \all{der Kauf} (l'achat), \all{der Käufer} (l'acheteur), \all{der Einkauf} (la course, l'achat), \all{verkaufen} (vendre), \all{der Verkäufer} (le vendeur).
\end{definition}

\begin{popfigure}
\begin{center}
\begin{tikzpicture}[mindmap, grow cyclic, every node/.style={concept, align=center}, concept color=popblue!40, text=popdark,
root concept/.append style={font=\bfseries\large, concept color=poporange!50},
level 1/.append style={level distance=3.2cm, sibling angle=72},
level 2/.append style={level distance=2.6cm, sibling angle=50}]
\node[root concept, align=center]{kaufen\\ \small acheter}
  child[concept color=popgreen!40]{ node[align=center]{der Kauf\\ \small l'achat} }
  child[concept color=popgreen!40]{ node[align=center]{der Käufer\\ \small l'acheteur} }
  child[concept color=poppurple!40]{ node[align=center]{der Einkauf\\ \small la course} 
    child[concept color=poppurple!30]{ node[align=center]{einkaufen\\ \small faire les courses} }
  }
  child[concept color=poppink!40]{ node[align=center]{verkaufen\\ \small vendre}
    child[concept color=poppink!30]{ node[align=center]{der Verkäufer\\ \small le vendeur} }
  }
  child[concept color=popgold!50]{ node[align=center]{der Kunde\\ \small le client} };
\end{tikzpicture}
\end{center}
\legende{Arbre de la famille de mots autour de la racine \all{kauf-}}
\end{popfigure}

\begin{center}
\begin{tabularx}{\linewidth}{|l|X|X|}
\hline
\rowcolor{popblueL} \textbf{Racine} & \textbf{Famille} & \textbf{Français} \\
\hline
\all{kaufen} & \all{der Kauf, der Käufer, der Einkauf, verkaufen, der Verkäufer} & achat, acheteur, course, vendre, vendeur \\
\hline
\all{dienen} & \all{der Dienst, die Dienstleistung} & service, prestation de service \\
\hline
\all{zahlen} & \all{die Zahl, die Zahlung, bezahlen} & le nombre/chiffre, le paiement, payer \\
\hline
\all{rechnen} & \all{die Rechnung, rechnen} & la facture/le calcul, calculer \\
\hline
\end{tabularx}
\end{center}

\begin{exemplebox}[La famille en action]
\all{Der Käufer bezahlt den Einkauf an der Kasse.} « L'acheteur paie ses courses à la caisse. »

\all{Der Verkäufer verkauft die Ware.} « Le vendeur vend la marchandise. »

\all{Die Zahlung ist mit Karte möglich.} « Le paiement par carte est possible. »

\all{Die Rechnung stimmt nicht.} « La facture / L'addition n'est pas juste. »
\end{exemplebox}

\courssub{Expressions utiles pour communiquer au magasin}

\begin{aretenir}[Les phrases à connaître par cœur]
\begin{itemize}
\item \all{Wie viel kostet das?} « Combien ça coûte ? »
\item \all{Was kostet dieser Pullover?} « Combien coûte ce pull ? »
\item \all{Das ist zu teuer!} « C'est trop cher ! »
\item \all{Haben Sie ein Angebot?} « Avez-vous une promotion ? »
\item \all{Ich nehme es.} « Je le prends. »
\item \all{Ich schaue mich nur um.} « Je regarde seulement. » (\all{sich umschauen}, séparable)
\item \all{Kann ich mit Karte bezahlen?} « Puis-je payer par carte ? »
\item \all{Die Rechnung, bitte!} « L'addition, s'il vous plaît ! »
\end{itemize}
\end{aretenir}

\begin{attention}[Faux amis : \all{bekommen} et \all{werden}]
Ne confondez jamais :
\begin{itemize}
\item \all{bekommen} = \textbf{recevoir, obtenir} : \all{Ich bekomme einen Rabatt.} « Je reçois une remise. »
\item \all{werden} = \textbf{devenir} : \all{Er wird Verkäufer.} « Il devient vendeur. »
\end{itemize}
\all{bekommen} ne signifie jamais « devenir » : c'est le faux ami le plus célèbre de l'allemand pour les francophones (et anglophones) !
\end{attention}

\begin{savaistu}[Négocier les prix : Cameroun et Allemagne]
Au Cameroun, on négocie (\all{handeln}) les prix au marché : c'est une tradition et presque un art ! En Allemagne, les prix dans les magasins et supermarchés sont \textbf{fixes} : on ne discute pas le prix chez le boulanger ou au supermarché. Mais il y a une exception : le \all{Flohmarkt} (le marché aux puces), où l'on peut négocier comme à Mokolo ou au marché central de Douala. Essaie cette phrase : \all{Geben Sie mir einen Rabatt?} « Vous me faites une remise ? »
\end{savaistu}

\begin{exoresolu}[Complète avec le mot qui convient]
Complète les phrases avec : \all{Kasse, Rabatt, Dienstleistung, einkaufen, Rechnung, Kunden}.

1. \all{Am Samstag \dots\ wir im Supermarkt \dots\ .}

2. \all{Der Friseur bietet eine \dots\ an.}

3. \all{Bitte zahlen Sie an der \dots\ !}

4. \all{Schüler bekommen 10 Prozent \dots\ .}

5. \all{Der Verkäufer hilft dem \dots\ .}

6. \all{Kann ich die \dots\ haben, bitte?}
\tcblower
\textbf{Solution détaillée :}

1. \all{Am Samstag \textbf{kaufen} wir im Supermarkt \textbf{ein}.} — Verbe à particule séparable : \all{kaufen} en position 2, \all{ein} en fin de phrase. « Samedi, nous faisons les courses au supermarché. »

2. \all{Der Friseur bietet eine \textbf{Dienstleistung} an.} — \all{anbieten} (séparable) : le coiffeur ne vend pas un objet, il fournit un service. « Le coiffeur propose une prestation de service. »

3. \all{Bitte zahlen Sie an der \textbf{Kasse}!} — \all{an der Kasse} = à la caisse (datif féminin). « Payez à la caisse, s'il vous plaît ! »

4. \all{Schüler bekommen 10 Prozent \textbf{Rabatt}.} — \all{Rabatt bekommen} = obtenir une remise. « Les élèves obtiennent 10 \% de réduction. »

5. \all{Der Verkäufer hilft dem \textbf{Kunden}.} — \all{helfen} + datif ; \all{der Kunde} est un nom faible : datif singulier \all{dem Kunden}. « Le vendeur aide le client. »

6. \all{Kann ich die \textbf{Rechnung} haben, bitte?} — Phrase typique au restaurant. « Puis-je avoir l'addition, s'il vous plaît ? »
\end{exoresolu}

\begin{exercice}[À toi de jouer !]
1. Traduis en allemand : « Le client paie la marchandise à la caisse. » / « Cette prestation de service est trop chère. » / « Je ne peux pas me permettre ce téléphone. »

2. Trouve la famille de mots : à partir de \all{verkaufen}, donne deux noms.

3. Classe ces mots en deux colonnes \emph{Bien} ou \emph{Service} : \all{das Brot, der Transport, die Schuhe, der Friseur, das Produkt, die Bank}.
\end{exercice}

\begin{corrige}[Exercice « À toi de jouer ! »]
1. \all{Der Kunde bezahlt die Ware an der Kasse.} / \all{Diese Dienstleistung ist zu teuer.} / \all{Ich kann mir dieses Telefon nicht leisten.}

2. \all{der Verkäufer} (le vendeur), \all{der Verkauf} (la vente).

3. \emph{Biens} : \all{das Brot, die Schuhe, das Produkt}. \emph{Services} : \all{der Transport, der Friseur, die Bank}.
\end{corrige}

\begin{aretenir}[L'essentiel de cette section]
\begin{itemize}
\item Chaque nom allemand s'apprend avec son \cle{article} et son \cle{pluriel} : \all{der Markt, die Märkte}.
\item Quatre champs lexicaux : acheter (\all{kaufen, einkaufen, bezahlen}), prix et argent (\all{der Preis, der Rabatt}), lieux et personnes (\all{das Geschäft, der Kunde}), biens et services (\all{die Ware, die Dienstleistung}).
\item Les \cle{familles de mots} multiplient ton vocabulaire : \all{kaufen → Kauf, Käufer, Einkauf, verkaufen, Verkäufer}.
\item Pièges : \all{der Kunde} est un nom faible (\all{den Kunden, dem Kunden, des Kunden}) ; \all{bekommen} = recevoir, \all{werden} = devenir ; \all{einkaufen} et \all{ausgeben} sont séparables ; \all{anbieten} (offrir/proposer) est séparable.
\end{itemize}
\end{aretenir}

\coursec{Grammaire 1 : la déclinaison de l'adjectif}

Dans la section précédente, nous avons découvert le vocabulaire de la consommation : \all{der Preis}, \all{die Ware}, \all{das Geschäft}, \all{der Kunde}, \all{die Dienstleistung}. En relisant notre texte \all{Einkaufen in Deutschland}, on remarque que ces noms sont souvent accompagnés d'adjectifs : un pain \emph{frais}, un prix \emph{bas}, une vendeuse \emph{sympathique}. Mais voilà le problème : en allemand, l'adjectif placé devant le nom change de terminaison, et ces changements obéissent à des règles précises. C'est l'une des difficultés majeures pour les francophones : nous allons la dompter méthodiquement.

\courssub{Observation : trois situations, trois comportements}

\begin{exemplebox}[Repérage dans le texte]
Observons ces groupes nominaux tirés de notre texte support :
\begin{itemize}
\item \all{das neue Radio} — « la nouvelle radio » (article défini \all{das})
\item \all{ein neues Radio} — « une nouvelle radio » (article indéfini \all{ein})
\item \all{gutes Brot} — « du bon pain » (pas d'article du tout)
\end{itemize}
L'adjectif \all{neu} apparaît sous deux formes : \textbf{neue} (après \all{das}) et \textbf{neues} (après \all{ein}) ; l'adjectif \all{gut} devient \textbf{gutes} (sans article). Constatons aussi :
\begin{itemize}
\item \all{Ich kaufe das billige Brot.} — « J'achète le pain bon marché. »
\item \all{Sie kauft ein teures Radio.} — « Elle achète une radio chère. »
\item \all{Er trinkt kaltes Wasser.} — « Il boit de l'eau froide. »
\end{itemize}
En revanche, dans \all{Das Radio ist neu.} (« La radio est neuve »), l'adjectif reste \textbf{invariable}, car il n'est pas placé devant le nom.
\end{exemplebox}

\begin{propriete}[La règle fondamentale]
\begin{enumerate}
\item L'adjectif \cle{épithète} (placé \textbf{devant le nom}) se \textbf{décline} : il prend une terminaison qui dépend du \textbf{genre}, du \textbf{nombre}, du \textbf{cas} et du \textbf{type d'article} qui précède.
\item L'adjectif \cle{attribut} (placé \textbf{après} les verbes \all{sein}, \all{werden}, \all{bleiben}) reste \textbf{invariable} : \all{Das Brot ist frisch.} « Le pain est frais. »
\end{enumerate}
\end{propriete}

\begin{attention}[Le piège numéro 1 des francophones]
En français, on dit « la radio est \emph{nouvelle} » : l'adjectif attribut s'accorde. En allemand, c'est \textbf{interdit} : on écrit \all{Das Radio ist neu.} et jamais *\all{ist neue}. L'adjectif ne se décline \textbf{que devant le nom}. De même : \all{Der Preis ist hoch.} « Le prix est élevé. », mais \all{der hohe Preis} « le prix élevé ».
\end{attention}

\courssub{Les trois déclinaisons : faible, mixte, forte}

L'allemand possède \textbf{trois tableaux} de terminaisons. Le choix dépend d'une seule question : \emph{qu'est-ce qui précède l'adjectif ?}

\begin{propriete}[Le principe directeur]
\begin{itemize}
\item Après un \textbf{article défini} (\all{der, die, das, die}) ou équivalent (\all{dieser, jener, jeder, alle}) : déclinaison \cle{faible} — seulement deux terminaisons possibles, \textbf{-e} ou \textbf{-en}. C'est la plus facile !
\item Après un \textbf{article indéfini} (\all{ein, eine}), la négation \all{kein} ou un \textbf{possessif} (\all{mein, dein, sein, unser...}) : déclinaison \cle{mixte} — l'adjectif prend \textbf{-er, -es, -e} quand l'article ne porte pas de marque, sinon \textbf{-en}.
\item \textbf{Sans article} du tout : déclinaison \cle{forte} — l'adjectif prend lui-même la terminaison de l'article défini (\all{guter Wein} = \all{der Wein}).
\end{itemize}
Idée directrice : la marque du genre et du cas doit apparaître \textbf{au moins une fois} dans le groupe nominal, soit sur l'article, soit sur l'adjectif.
\end{propriete}

\textbf{Tableau 1 : déclinaison faible} (après \all{der, die, das, dieser, jeder, alle})

\begin{center}
\begin{tabular}{|l|l|l|l|l|}
\hline
\rowcolor{popblueL} Cas & Masculin & Féminin & Neutre & Pluriel \\
\hline
Nominatif & der gut\textbf{e} Wein & die gut\textbf{e} Milch & das gut\textbf{e} Brot & die gut\textbf{en} Waren \\
\hline
Accusatif & den gut\textbf{en} Wein & die gut\textbf{e} Milch & das gut\textbf{e} Brot & die gut\textbf{en} Waren \\
\hline
Datif & dem gut\textbf{en} Wein & der gut\textbf{en} Milch & dem gut\textbf{en} Brot & den gut\textbf{en} Waren \\
\hline
Génitif & des gut\textbf{en} Weins & der gut\textbf{en} Milch & des gut\textbf{en} Brots & der gut\textbf{en} Waren \\
\hline
\end{tabular}
\end{center}

\begin{aretenir}[Astuce de mémorisation]
Déclinaison faible : \textbf{cinq cases en -e} (nominatif des trois genres + accusatif féminin et neutre), \textbf{tout le reste en -en}. Au pluriel, c'est toujours \textbf{-en}. Exemple : \all{Die guten Waren sind teuer.} « Les bonnes marchandises sont chères. »
\end{aretenir}

\textbf{Tableau 2 : déclinaison mixte} (après \all{ein, kein, mein, dein, unser...})

\begin{center}
\begin{tabular}{|l|l|l|l|l|}
\hline
\rowcolor{poporangeL} Cas & Masculin & Féminin & Neutre & Pluriel (kein/mein) \\
\hline
Nominatif & ein gut\textbf{er} Wein & eine gut\textbf{e} Milch & ein gut\textbf{es} Brot & keine gut\textbf{en} Waren \\
\hline
Accusatif & einen gut\textbf{en} Wein & eine gut\textbf{e} Milch & ein gut\textbf{es} Brot & keine gut\textbf{en} Waren \\
\hline
Datif & einem gut\textbf{en} Wein & einer gut\textbf{en} Milch & einem gut\textbf{en} Brot & keinen gut\textbf{en} Waren \\
\hline
Génitif & eines gut\textbf{en} Weins & einer gut\textbf{en} Milch & eines gut\textbf{en} Brots & keiner gut\textbf{en} Waren \\
\hline
\end{tabular}
\end{center}

La logique : \all{ein} (masculin nominatif, neutre nominatif/accusatif) ne porte \textbf{aucune marque} ; l'adjectif prend alors la marque de l'article défini : \textbf{-er} (comme \all{der}) et \textbf{-es} (comme \all{das}). Partout ailleurs, l'article porte déjà la marque (\all{einen, einem, einer, eines, keine}), donc l'adjectif prend simplement \textbf{-en}.

\textbf{Tableau 3 : déclinaison forte} (sans article)

\begin{center}
\begin{tabular}{|l|l|l|l|l|}
\hline
\rowcolor{popgreenL} Cas & Masculin & Féminin & Neutre & Pluriel \\
\hline
Nominatif & gut\textbf{er} Wein & gut\textbf{e} Milch & gut\textbf{es} Brot & gut\textbf{e} Waren \\
\hline
Accusatif & gut\textbf{en} Wein & gut\textbf{e} Milch & gut\textbf{es} Brot & gut\textbf{e} Waren \\
\hline
Datif & gut\textbf{em} Wein & gut\textbf{er} Milch & gut\textbf{em} Brot & gut\textbf{en} Waren \\
\hline
Génitif & gut\textbf{en} Weins & gut\textbf{er} Milch & gut\textbf{en} Brots & gut\textbf{er} Waren \\
\hline
\end{tabular}
\end{center}

Ici, l'adjectif « fait le travail » de l'article absent : il prend presque toutes les terminaisons de \all{der/die/das} (\textbf{-er, -e, -es, -em, -en}). Seule exception : le génitif masculin et neutre en \textbf{-en} (car le nom porte déjà le \textbf{-s} : \all{guten Weins}).

\begin{exemplebox}[Exemples en contexte marchand]
\begin{itemize}
\item \all{Der frische Fisch kommt aus Kribi.} — « Le poisson frais vient de Kribi. » (faible, nominatif masculin)
\item \all{Ich kaufe einen billigen Reis.} — « J'achète un riz bon marché. » (mixte, accusatif masculin)
\item \all{Frische Milch kostet hier 1,50 Euro.} — « Le lait frais coûte ici 1,50 euro. » (forte, nominatif féminin)
\item \all{Mit kaltem Wasser wäscht man die Wäsche.} — « Avec de l'eau froide, on lave le linge. » (forte, datif neutre)
\item \all{Die neuen Handys sind sehr teuer.} — « Les nouveaux téléphones portables sont très chers. » (faible, pluriel)
\end{itemize}
\end{exemplebox}

\begin{popfigure}
\begin{center}
\begin{tikzpicture}[font=\small]
% Titre des colonnes
\node[fill=popblue, text=white, rounded corners, minimum width=3.6cm, align=center] at (0,0) {\textbf{Faible}\\der / die / das};
\node[fill=poporange, text=white, rounded corners, minimum width=3.6cm, align=center] at (5,0) {\textbf{Mixte}\\ein / kein / mein};
\node[fill=popgreen, text=white, rounded corners, minimum width=3.6cm, align=center] at (10,0) {\textbf{Forte}\\sans article};
% Lignes de cas (Masculin / Neutre)
\node[align=center] at (0,-1.1) {N : der gut\textbf{e}\\A : den gut\textbf{en}\\D : dem gut\textbf{en}\\G : des gut\textbf{en}};
\node[align=center] at (5,-1.1) {N : ein gut\textbf{er}\\A : einen gut\textbf{en}\\D : einem gut\textbf{en}\\G : eines gut\textbf{en}};
\node[align=center] at (10,-1.1) {N : gut\textbf{er}\\A : gut\textbf{en}\\D : gut\textbf{em}\\G : gut\textbf{en}};
% Féminin / Pluriel
\node[align=center, text=popdark] at (0,-2.6) {die gut\textbf{e}\\das gut\textbf{e}\\Pl. : die gut\textbf{en}};
\node[align=center, text=popdark] at (5,-2.6) {eine gut\textbf{e}\\ein gut\textbf{es}\\Pl. : keine gut\textbf{en}};
\node[align=center, text=popdark] at (10,-2.6) {gut\textbf{e}\\gut\textbf{es}\\Pl. : gut\textbf{e}};
\node[fill=popgoldL, rounded corners, align=center, text=popdark] at (5,-4) {Règle d'or : la marque du cas apparaît au moins une fois,\\sur l'article OU sur l'adjectif.};
\end{tikzpicture}
\end{center}
\legende{Les trois déclinaisons de l'adjectif (exemple du masculin \all{gut} + \all{Wein})}
\end{popfigure}

\courssub{Cas particuliers et adjectifs invariables}

\begin{propriete}[Petites modifications orthographiques]
\begin{itemize}
\item Adjectifs en \textbf{-el} : le \textbf{e} disparaît devant la terminaison : \all{dunkel} → \all{ein dunkles Zimmer} « une chambre sombre ».
\item Adjectifs en \textbf{-er} : le \textbf{e} peut disparaître : \all{teuer} → \all{ein teures Radio} « une radio chère » ; \all{sauer} → \all{saure Milch} « du lait aigre ».
\item Adjectifs en \textbf{-en} (participes ou matériaux) : on ajoute simplement la terminaison : \all{das goldene Armband} « le bracelet en or ».
\item \all{hoch} perd son \textbf{c} devant une voyelle : \all{der hohe Preis} « le prix élevé ».
\end{itemize}
\end{propriete}

\begin{definition}[Adjectifs invariables]
Certains adjectifs, surtout des couleurs d'origine étrangère et des adjectifs de ville, ne se déclinent \textbf{jamais} : \all{rosa} « rose », \all{lila} « mauve », \all{prima} « super », \all{orange} « orange ». Exemple : \all{Sie kauft eine rosa Bluse.} « Elle achète un chemisier rose. » Les adjectifs géographiques en \textbf{-er} sont aussi invariables et prennent la majuscule : \all{die Yaoundéer Straßen} « les rues de Yaoundé », \all{der Berliner Markt} « le marché de Berlin ».
\end{definition}

\courssub{Comparaison français / allemand}

\begin{center}
\begin{tabularx}{\linewidth}{|X|X|X|}
\hline
\rowcolor{popblueL} Français & Deutsch & Remarque \\
\hline
L'adjectif s'accorde en genre et nombre : un pain frais, une radio neuve. & L'adjectif se décline en genre, nombre \textbf{et cas} : \all{ein frisches Brot}, \all{mit frischem Brot}. & L'allemand ajoute la dimension du \textbf{cas}, absente du français. \\
\hline
L'adjectif attribut s'accorde : la radio est neuve. & L'adjectif attribut est invariable : \all{Das Radio ist neu.} & Erreur typique : *\all{ist neue}. \\
\hline
La place est souvent après le nom : un prix bas. & L'adjectif épithète est toujours \textbf{avant} le nom : \all{ein niedriger Preis}. & Jamais *\all{ein Preis niedriger}. \\
\hline
\end{tabularx}
\end{center}

\courssub{Méthode : décliner un adjectif en trois questions}

\begin{methode}[Trouver la bonne terminaison]
Pour décliner correctement un adjectif épithète, pose-toi trois questions dans l'ordre :
\begin{enumerate}
\item \textbf{Y a-t-il un article ?} Article défini (\all{der/die/das}) → tableau faible. Article indéfini, \all{kein} ou possessif → tableau mixte. Rien → tableau forte.
\item \textbf{Quel est le cas ?} Sujet → nominatif ; COD → accusatif ; après \all{mit, zu, von, bei} ou COI → datif ; possession → génitif.
\item \textbf{Quel est le genre et le nombre du nom ?} Masculin, féminin, neutre ou pluriel : c'est le \textbf{nom} qui commande, jamais l'adjectif.
\end{enumerate}
Croise les réponses dans le tableau : tu obtiens la terminaison. Exemple : \all{Ich kaufe ... (frisch) Brot.} → pas d'article (forte), accusatif (COD de \all{kaufen}), neutre (\all{das Brot}) → \all{Ich kaufe frisches Brot.}
\end{methode}

\courssub{Texte d'application}

\begin{textebox}[Auf dem Markt in Yaoundé]
Am Samstag geht Familie Mbarga auf den großen Markt. Die Mutter kauft frisches Gemüse und reife Bananen. „Die roten Tomaten sind heute billig“, sagt sie. Der Vater sucht ein neues Radio, aber die teuren Geräte gefallen ihm nicht. Endlich findet er ein kleines Radio zu einem guten Preis. Die freundliche Verkäuferin macht ihm einen guten Rabatt. Der Sohn Kevin trinkt kaltes Wasser und isst ein süßes Brötchen. „Der süße Kuchen dort sieht auch gut aus!“, ruft er. Die kleine Tochter bekommt ein rosa Kleid und eine warme Jacke für die kalte Jahreszeit. Mit schweren Taschen geht die glückliche Familie nach Hause. „Gutes Essen und gute Preise — das ist unser Markt!“, sagt der Vater zufrieden.
\end{textebox}

\textbf{Glossaire} : \all{der Markt, die Märkte} « le marché » ; \all{das Gemüse, -} « les légumes » ; \all{reif} « mûr » ; \all{das Gerät, -e} « l'appareil » ; \all{der Rabatt, -e} « la remise » ; \all{das Brötchen, -} « le petit pain » ; \all{süß} « sucré » ; \all{die Tasche, -n} « le sac » ; \all{zufrieden} « satisfait ».

\textbf{Questions de compréhension} : 1) \all{Was kauft die Mutter?} 2) \all{Warum kauft der Vater zuerst kein Radio?} 3) \all{Was bekommt die Tochter?} 4) Relevez cinq adjectifs épithètes du texte et indiquez pour chacun le type de déclinaison (faible, mixte, forte).

\courssub{Exercice résolu et entraînement}

\begin{exoresolu}[Complète avec la bonne terminaison]
Énoncé : complète les phrases avec la terminaison correcte de l'adjectif entre parenthèses.
1) \all{Ich kaufe das frisch\_\_ Brot.} (frisch) 2) \all{Sie trinkt ein kalt\_\_ Wasser.} (kalt) 3) \all{Der neu\_\_ Kunde ist freundlich.} (neu) 4) \all{Wir essen gut\_\_ Suppe.} (gut, sans article) 5) \all{Die teur\_\_ Waren sind im Geschäft.} (teuer)
\tcblower
Solution détaillée :
1) Article défini \all{das} → faible ; accusatif neutre → \textbf{-e} : \all{das frische Brot}.
2) \all{ein} sans marque → mixte ; accusatif neutre → \textbf{-es} : \all{ein kaltes Wasser}.
3) \all{der} → faible ; nominatif masculin → \textbf{-e} : \all{Der neue Kunde}.
4) Sans article → forte ; accusatif féminin (\all{die Suppe}) → \textbf{-e} : \all{gute Suppe}.
5) \all{die} pluriel → faible ; nominatif pluriel → \textbf{-en} (et \all{teuer} perd son e) : \all{Die teuren Waren}.
\end{exoresolu}

\begin{exercice}[À toi de jouer]
Complète : 1) \all{Er kauft einen billig\_\_ Reis.} 2) \all{Die klein\_\_ Tochter isst ein süß\_\_ Brötchen.} 3) \all{Frisch\_\_ Milch ist gesund.} 4) \all{Mit kalt\_\_ Wasser (datif neutre)}. Puis traduis : « J'achète le pain frais. » / « Elle achète une radio chère. » / « Il boit de l'eau froide. »
\end{exercice}

\begin{corrige}[Exercice « À toi de jouer »]
1) \all{einen billigen Reis} (mixte, accusatif masculin). 2) \all{Die kleine Tochter ... ein süßes Brötchen} (faible nominatif féminin ; mixte accusatif neutre). 3) \all{Frische Milch} (forte, nominatif féminin). 4) \all{mit kaltem Wasser} (forte, datif neutre). Traductions : \all{Ich kaufe das frische Brot.} / \all{Sie kauft ein teures Radio.} / \all{Er trinkt kaltes Wasser.}
\end{corrige}

\begin{attention}[Récapitulatif des pièges]
\begin{itemize}
\item Ne décline jamais l'adjectif attribut : \all{Das Radio ist neu.} (pas *\all{neue}).
\item Après \all{ein} au masculin nominatif et au neutre, l'adjectif prend \textbf{-er / -es} : \all{ein guter Wein}, \all{ein gutes Brot} — pas *\all{ein gute Wein}.
\item Au pluriel après \all{die} ou \all{keine}, c'est toujours \textbf{-en} : \all{die guten Waren}.
\item L'adjectif se place \textbf{toujours avant} le nom en allemand : \all{ein niedriger Preis}, jamais après comme en français.
\end{itemize}
\end{attention}

\coursec{Grammaire 2 : comparatif et superlatif}

Dans la section précédente, nous avons appris à \cle{décliner l'adjectif} : \all{der billige Preis}, \all{ein neues Handy}, \all{gute Waren}. Mais quand on fait ses courses, on ne se contente pas de décrire : on \textbf{compare}. Ce pain-ci est moins cher, ce magasin-là est plus grand, cette offre est la meilleure. C'est exactement ce que font les personnages de notre texte \all{Einkaufen in Deutschland}. Pour exprimer ces comparaisons, l'allemand utilise deux formes de l'adjectif : le \cle{comparatif} (\all{Komparativ}) et le \cle{superlatif} (\all{Superlativ}). Observons d'abord, réglons ensuite.

\courssub{Observation : trois degrés de l'adjectif}

\begin{textebox}[Mini-dialogue au marché — extrait adapté du texte support]
A: „Der Markt in Yaoundé ist groß, aber der Supermarkt ist größer.“\\
B: „Ja, und die Preise?“\\
A: „Das Brot ist hier billig. Im Supermarkt ist es billiger. Aber auf dem Wochenmarkt ist es am billigsten.“\\
B: „Und die Qualität?“\\
A: „Die Ware ist gut, die Ware im Supermarkt ist besser, aber die frischen Produkte vom Bauern sind am besten.“
\end{textebox}

\textbf{Glossaire :} \all{der Preis, -e} (le prix) ; \all{billig} (bon marché) ; \all{die Qualität, -en} (la qualité) ; \all{frisch} (frais) ; \all{der Bauer, -n} (le paysan) ; \all{das Produkt, -e} (le produit).

Observons les formes en gras logique :

\begin{center}
\begin{tabularx}{\linewidth}{|X|X|X|}
\hline
\rowcolor{popblueL} \textbf{Positiv} (degré positif) & \textbf{Komparativ} (comparatif) & \textbf{Superlativ} (superlatif) \\
\hline
billig & billiger & am billigsten \\
\hline
groß & größer & am größten \\
\hline
gut & besser & am besten \\
\hline
teuer & teurer & am teuersten \\
\hline
\end{tabularx}
\end{center}

On constate trois choses :
\begin{itemize}
\item le comparatif se forme avec la terminaison \all{-er} : \all{billig} $\rightarrow$ \all{billiger} ;
\item le superlatif (prédicatif, après \all{sein}) se forme avec \all{am} + adjectif + \all{-sten} : \all{am billigsten} ;
\item certains adjectifs changent de voyelle (\all{groß} $\rightarrow$ \all{größer}) ou de radical (\all{gut} $\rightarrow$ \all{besser}).
\end{itemize}

\begin{popfigure}
\begin{tikzpicture}[
  level/.style={rectangle, rounded corners=4pt, draw=popblue, fill=popblueL, align=center, minimum width=3.6cm, minimum height=1.5cm, font=\small},
  fleche/.style={-{Stealth[length=3mm]}, thick, poporange}
]
\node[level, align=center] (pos) at (0,0){\textbf{Positiv}\\ \all{teuer}\\ (cher)};
\node[level, fill=poporangeL, draw=poporange, align=center] (komp) at (5.2,0){\textbf{Komparativ}\\ \all{teurer (als)}\\ (plus cher que)};
\node[level, fill=popgreenL, draw=popgreen, align=center] (sup) at (10.4,0){\textbf{Superlativ}\\ \all{am teuersten}\\ (le plus cher)};
\draw[fleche] (pos) -- (komp) node[midway, above, font=\footnotesize, popdark]{+ \all{-er}};
\draw[fleche] (komp) -- (sup) node[midway, above, font=\footnotesize, popdark]{\all{am} + \all{-sten}};
\end{tikzpicture}
\legende{Frise des trois degrés de l'adjectif : \all{teuer} $\rightarrow$ \all{teurer} $\rightarrow$ \all{am teuersten}}
\end{popfigure}

\courssub{Formation régulière du comparatif et du superlatif}

\begin{propriete}[Formation régulière]
\begin{itemize}
\item \textbf{Comparatif} : adjectif + \all{-er}, suivi de \all{als} (que) pour introduire le second terme : \all{Das Radio ist teuer\textbf{er} als das Handy.} (La radio est plus chère que le portable.)
\item \textbf{Superlatif} (adjectif attribut du verbe, après \all{sein}) : \all{am} + adjectif + \all{-sten} : \all{Der Fernseher ist am teuer\textbf{sten}.} (C'est la télévision qui est la plus chère.)
\item \textbf{Umlaut} : les adjectifs \textbf{monosyllabiques} dont la voyelle est \textbf{a, o, u} prennent presque toujours l'Umlaut au comparatif \textbf{et} au superlatif : \all{alt} $\rightarrow$ \all{älter} $\rightarrow$ \all{am ältesten} ; \all{groß} $\rightarrow$ \all{größer} $\rightarrow$ \all{am größten} ; \all{warm} $\rightarrow$ \all{wärmer} $\rightarrow$ \all{am wärmsten} ; \all{jung} $\rightarrow$ \all{jünger} $\rightarrow$ \all{am jüngsten} ; \all{stark} $\rightarrow$ \all{stärker} $\rightarrow$ \all{am stärksten}.
\end{itemize}
\end{propriete}

\textbf{Petites particularités orthographiques :}
\begin{itemize}
\item les adjectifs terminés par \all{-d}, \all{-t}, \all{-s}, \all{-ß}, \all{-sch}, \all{-z} prennent \all{-esten} au superlatif pour faciliter la prononciation : \all{alt} $\rightarrow$ \all{am ält\textbf{est}en}, \all{heiß} $\rightarrow$ \all{am heiß\textbf{est}en}, \all{kurz} $\rightarrow$ \all{am kürz\textbf{est}en}, \all{frisch} $\rightarrow$ \all{am frisch\textbf{est}en} ;
\item les adjectifs en \all{-er} et \all{-el} perdent souvent le \all{-e-} au comparatif : \all{teuer} $\rightarrow$ \all{teu\textbf{rer}} (et non *\all{teuerer}), \all{dunkel} $\rightarrow$ \all{dun\textbf{kler}} ;
\item \textbf{exception à l'Umlaut} : quelques adjectifs monosyllabiques en a/o/u ne prennent \emph{pas} l'Umlaut : \all{klar} $\rightarrow$ \all{klarer}, \all{voll} $\rightarrow$ \all{voller}, \all{falsch} $\rightarrow$ \all{falscher}. Il faut les apprendre par cœur.
\end{itemize}

\begin{exemplebox}[Exemples traduits]
\all{Der Supermarkt ist billiger als der Markt.} « Le supermarché est moins cher que le marché. »

\all{Das ist das beste Angebot.} « C'est la meilleure offre. »

\all{Im Juli ist es am wärmsten.} « C'est en juillet qu'il fait le plus chaud. »

\all{Mein Bruder ist älter als ich.} « Mon frère est plus âgé que moi. »

\all{Die Waren im Geschäft sind teurer als auf dem Markt.} « Les marchandises du magasin sont plus chères qu'au marché. »

\all{Douala ist größer als Bamenda, aber Yaoundé ist die größte Stadt? Nein, Douala ist am größten!} « Douala est plus grande que Bamenda, mais Yaoundé est la plus grande ville ? Non, Douala est la plus grande ! »
\end{exemplebox}

\courssub{Les adjectifs irréguliers}

Comme en français (\emph{bon} $\rightarrow$ \emph{meilleur}, et non *\emph{plus bon}), quelques adjectifs très fréquents changent complètement de radical. Ils sont à mémoriser impérativement :

\begin{center}
\begin{tabularx}{\linewidth}{|X|X|X|X|}
\hline
\rowcolor{popblueL} \textbf{Positiv} & \textbf{Komparativ} & \textbf{Superlativ} & \textbf{Français} \\
\hline
gut & besser & am besten & bon / meilleur / le meilleur \\
\hline
viel & mehr & am meisten & beaucoup / plus / le plus \\
\hline
wenig & weniger & am wenigsten & peu / moins / le moins \\
\hline
gern & lieber & am liebsten & volontiers / plutôt / de préférence \\
\hline
hoch & höher & am höchsten & haut / plus haut / le plus haut \\
\hline
nah & näher & am nächsten & proche / plus proche / le plus proche \\
\hline
\end{tabularx}
\end{center}

\begin{exemplebox}[Exemples traduits]
\all{Ich kaufe lieber auf dem Markt ein.} « Je préfère faire mes courses au marché. »

\all{Das Geschäft ist näher als der Supermarkt.} « Le magasin est plus proche que le supermarché. »

\all{Im Sommer trinkt man am meisten Wasser.} « C'est en été qu'on boit le plus d'eau. »

\all{Die Preise in Deutschland sind höher als in Kamerun.} « Les prix en Allemagne sont plus élevés qu'au Cameroun. »

\all{Ich habe weniger Geld als du.} « J'ai moins d'argent que toi. »
\end{exemplebox}

\courssub{Emplois : comparer avec \all{als}, égalité avec \all{so ... wie}}

\begin{propriete}[Les constructions de la comparaison]
\begin{itemize}
\item \textbf{Supériorité / infériorité} : comparatif + \all{als} : \all{Das Radio ist teurer \textbf{als} das alte.} (La radio est plus chère que l'ancienne.)
\item \textbf{Égalité} : \all{so + adjectif au positif + wie} : \all{Das Handy ist \textbf{so} billig \textbf{wie} das Radio.} (Le portable est aussi bon marché que la radio.)
\item \textbf{Inégalité} (négation de l'égalité) : \all{nicht so ... wie} : \all{Der Markt ist \textbf{nicht so} groß \textbf{wie} der Supermarkt.} (Le marché n'est pas aussi grand que le supermarché.)
\end{itemize}
\end{propriete}

\begin{center}
\begin{tabularx}{\linewidth}{|X|X|X|}
\hline
\rowcolor{popblueL} \textbf{Français} & \textbf{Deutsch} & \textbf{Remarque} \\
\hline
plus ... que & -er ... als & \all{teurer als} \\
\hline
moins ... que & weniger ... als & \all{weniger teuer als} (comparatif de \all{wenig}) \\
\hline
aussi ... que & so ... wie & \all{so billig wie} \\
\hline
le / la plus ... (épithète) & der / die / das ...-ste & adjectif décliné \\
\hline
le plus ... (avec sein) & am ...-sten & \all{am billigsten} \\
\hline
\end{tabularx}
\end{center}

\courssub{Le superlatif épithète : \all{der beste Preis}}

Quand le superlatif est placé \textbf{devant le nom} (épithète), il se comporte comme un adjectif ordinaire : il prend la terminaison \all{-st-} \textbf{puis la désinence de la déclinaison} vue dans la section précédente (déclinaison forte après article défini, faible après article indéfini).

\begin{center}
\begin{tabular}{|l|l|l|}
\hline
\rowcolor{popblueL} \textbf{Masculin} & \textbf{Neutre} & \textbf{Féminin / Pluriel} \\
\hline
der beste Preis & das beste Angebot & die beste Ware \\
\hline
der billigste Markt & das größte Geschäft & die besten Waren \\
\hline
ein besseres Produkt & ein billigeres Handy & eine bessere Qualität \\
\hline
\end{tabular}
\end{center}

\all{Der beste Preis ist auf dem Markt.} « Le meilleur prix est au marché. » — \all{Sie sucht das billigste Angebot.} « Elle cherche l'offre la moins chère. »

\begin{attention}[Pièges des francophones]
\begin{itemize}
\item Ne dites jamais *\all{mehr billig} sur le modèle français « plus bon marché » : le comparatif allemand se forme \textbf{toujours} avec \all{-er} : \all{billiger}, \all{interessanter}, \all{moderner}. \all{mehr} ne s'emploie que comme comparatif de \all{viel} (plus, en quantité).
\item Ne dites jamais *\all{guter} : le comparatif de \all{gut} est \all{besser} (comme « meilleur », et non *\emph{plus bon}).
\item Ne confondez pas \all{als} et \all{wenn} : \all{als} sert à \textbf{comparer} (\all{größer als} = plus grand que) ; \all{wenn} signifie « quand / si ». *\all{größer wenn} est une erreur grave.
\item Au superlatif après \all{sein}, n'oubliez pas \all{am} : \all{Er ist am besten.} (et non *\all{Er ist besten}).
\item Pour « moins ... que », on utilise \all{weniger + adjectif + als} (jamais *\all{minder}).
\end{itemize}
\end{attention}

\begin{exoresolu}[Comparer des magasins]
\textbf{Énoncé.} Complétez avec la forme correcte de l'adjectif entre parenthèses.

1. Der Supermarkt ist \dots\ als der kleine Laden. (groß)

2. Auf dem Markt sind die Früchte \dots\ . (frisch — superlatif)

3. Dieses Angebot ist \dots\ als das andere. (gut)

4. Meine Schwester kauft \dots\ im Internet ein. (gern — superlatif)

5. Der neue Fernseher ist \dots\ als der alte. (teuer)
\tcblower
\textbf{Solution détaillée.}

1. \all{Der Supermarkt ist \textbf{größer} als der kleine Laden.} Comparatif régulier de \all{groß} : + \all{-er} et Umlaut (monosyllabe en \emph{o}).

2. \all{Auf dem Markt sind die Früchte \textbf{am frischesten}.} Superlatif : \all{am} + \all{-esten} (car \all{frisch} se termine en \all{-sch}).

3. \all{Dieses Angebot ist \textbf{besser} als das andere.} Comparatif irrégulier de \all{gut} : jamais *\all{guter}.

4. \all{Meine Schwester kauft \textbf{am liebsten} im Internet ein.} Superlatif irrégulier de \all{gern} : \all{gern} $\rightarrow$ \all{lieber} $\rightarrow$ \all{am liebsten}.

5. \all{Der neue Fernseher ist \textbf{teurer} als der alte.} Adjectif en \all{-er} : le \all{-e-} disparaît au comparatif (\all{teurer}, non *\all{teuerer}).
\end{exoresolu}

\begin{savaistu}[Les Allemands, champions de la comparaison de prix]
Les Allemands adorent comparer les prix avant d'acheter ! Les sites de comparaison, les \all{Vergleichsportale}, y sont extrêmement populaires : on y compare les assurances, l'électricité, les voyages et même les supermarchés. Le mot \all{der Preisvergleich} (la comparaison de prix) fait partie du vocabulaire quotidien. Au Cameroun aussi, on compare volontiers les prix entre le marché et la boutique du quartier : \all{Wo ist es am billigsten?} — « Où est-ce le moins cher ? »
\end{savaistu}

\begin{exercice}[À vous de comparer]
1. Traduisez en allemand : a) Le marché est plus grand que le magasin. b) Cette offre est la meilleure. c) Le pain est aussi frais qu'au supermarché.

2. Formez les trois degrés : \all{alt}, \all{schnell}, \all{viel}, \all{klein}, \all{gut}, \all{wenig}.

3. Rédigez trois phrases pour comparer deux magasins de votre quartier (utilisez \all{als}, \all{so ... wie} et un superlatif).
\end{exercice}

\begin{corrige}[Exercice « À vous de comparer »]
1. a) \all{Der Markt ist größer als das Geschäft.} b) \all{Dieses Angebot ist das beste.} c) \all{Das Brot ist so frisch wie im Supermarkt.}

2. \all{alt — älter — am ältesten} ; \all{schnell — schneller — am schnellsten} ; \all{viel — mehr — am meisten} ; \all{klein — kleiner — am kleinsten} ; \all{gut — besser — am besten} ; \all{wenig — weniger — am wenigsten}.

3. Exemple de production : \all{Der Laden neben der Schule ist teurer als der Markt. Die Waren sind nicht so frisch wie auf dem Markt. Aber der Laden ist am nächsten.} (Le magasin à côté de l'école est plus cher que le marché. Les marchandises ne sont pas aussi fraîches qu'au marché. Mais le magasin est le plus proche.)
\end{corrige}

\begin{aretenir}[L'essentiel du comparatif et du superlatif]
\begin{itemize}
\item Comparatif : adjectif + \all{-er} + \all{als} ; superlatif : \all{am} + adjectif + \all{-sten} ; superlatif épithète : \all{der/die/das ...-ste} (décliné).
\item Umlaut pour les monosyllabes en a/o/u : \all{alt — älter — am ältesten}.
\item Irréguliers à connaître par cœur : \all{gut — besser — am besten} ; \all{viel — mehr — am meisten} ; \all{wenig — weniger — am wenigsten} ; \all{gern — lieber — am liebsten} ; \all{hoch — höher — am höchsten} ; \all{nah — näher — am nächsten}.
\item Égalité : \all{so ... wie} ; jamais *\all{mehr billig}, jamais *\all{guter}, et \all{als} $\neq$ \all{wenn}.
\end{itemize}
\end{aretenir}

\coursec{Conjugaison : les verbes à particule séparable}

Dans la section précédente, nous avons appris à comparer les prix et les produits grâce au comparatif et au superlatif (\all{billig -- billiger -- am billigsten}). Mais pour raconter nos achats, il ne suffit pas de comparer : il faut aussi conjuguer correctement les verbes du texte \emph{Einkaufen in Deutschland}. Or plusieurs de ces verbes présentent une particularité étonnante pour un francophone : une partie du verbe « saute » à la fin de la phrase. C'est le phénomène des \cle{verbes à particule séparable} (\all{trennbare Verben}), l'une des structures les plus caractéristiques -- et les plus évaluées -- de la langue allemande.

\courssub{Observation : la particule se détache}

\begin{textebox}[Am Samstag bei Familie Berger]
Am Samstagmorgen steht Frau Berger früh auf. Sie macht um acht Uhr das Fenster auf und hört das Radio. Dann nimmt sie ihren Einkaufskorb mit und geht zum Markt. Auf dem Markt kauft sie Obst und Gemüse ein. Die Tomaten probiert sie zuerst an, denn sie will nur frische Ware. Sie gibt heute nicht viel Geld aus, weil die Preise niedrig sind. Um elf Uhr ruft sie ihre Tochter an und lädt sie zum Kaffee ein. Nachmittags bestellt sie im Internet ein neues Buch und bezahlt es sofort mit der Karte. Der Verkäufer verspricht eine schnelle Lieferung. Am Abend kommt Frau Berger müde zurück, aber sie ist zufrieden: sie hat alles erledigt. Ihr Mann bringt die Taschen in die Küche und räumt den Kühlschrank auf. „Einkaufen ist anstrengend, aber schön“, sagt sie und lacht.
\end{textebox}

\textbf{Glossaire :} der Einkaufskorb, -körbe (le panier à provisions) ; die Ware, -n (la marchandise) ; anprobieren (essayer [un vêtement]) ; ausgeben (dépenser) ; anrufen (téléphoner à) ; einladen (inviter) ; bestellen (commander) ; bezahlen (payer) ; versprechen (promettre) ; die Lieferung, -en (la livraison) ; zurückkommen (revenir) ; aufräumen (ranger, mettre en ordre) ; erledigen (régler, accomplir) ; anstrengend (fatigant).

\textbf{Questions de compréhension :}
\begin{enumerate}
\item Wann geht Frau Berger zum Markt? (Quand Mme Berger va-t-elle au marché ?)
\item Warum gibt sie heute nicht viel Geld aus? (Pourquoi ne dépense-t-elle pas beaucoup d'argent aujourd'hui ?)
\item Was bestellt sie am Nachmittag? (Que commande-t-elle l'après-midi ?)
\end{enumerate}

\begin{exemplebox}[Observons les verbes du texte]
Comparons deux phrases :
\begin{itemize}
\item \all{Frau Berger \textbf{kauft} Obst und Gemüse \textbf{ein}.} -- « Mme Berger achète des fruits et des légumes. »
\item \all{Sie \textbf{bezahlt} das Buch sofort.} -- « Elle paie le livre immédiatement. »
\end{itemize}
Dans la première phrase, le verbe \all{einkaufen} est coupé en deux : \all{kauf-} se conjugue à la deuxième position, et la particule \all{ein-} part à la fin. Dans la seconde, \all{bezahlen} reste entier. Pourquoi ? Tout est une question de \cle{particule}.
\end{exemplebox}

\courssub{La règle : particule séparable ou inséparable ?}

\begin{definition}[Verbe à particule séparable]
Un verbe à particule séparable (\all{trennbares Verb}) est formé d'un verbe de base (\all{kaufen}) et d'une particule (\all{ein-}). Au présent et au Präteritum, dans une proposition principale, la particule \cle{se détache} du verbe conjugué et se place \cle{en toute fin de phrase}. Le verbe conjugué reste, lui, à la deuxième position.
\end{definition}

\begin{propriete}[Règle d'or : l'accent décide]
\begin{enumerate}
\item \textbf{Particules séparables} : elles portent l'\cle{accent tonique} quand on prononce le verbe : \all{EINkaufen, AUSgeben, AUFmachen, MITnehmen, ANprobieren, ZURÜCKkommen}. Ce sont notamment : \all{ein-, aus-, auf-, mit-, an-, ab-, zu-, zurück-, weg-, vor-}.
\item \textbf{Particules inséparables} : jamais accentuées, elles restent \cle{collées} au verbe : \all{be-, ge-, er-, ver-, zer-, ent-, emp-}. Exemples : \all{beZAHLEN, verKAUFEN, beSTELLEN, erLEdigen}.
\item \textbf{Particules variables} : \all{durch-, über-, um-, unter-, wieder-} sont séparables si elles sont accentuées (\all{WIEDERholen} = répéter), inséparables si l'accent tombe sur le verbe (\all{wiederHOlen} = aller rechercher).
\end{enumerate}
\end{propriete}

\begin{methode}[Comment reconnaître un verbe séparable en trois étapes]
\begin{enumerate}
\item \textbf{Prononce} le verbe à voix haute : \all{EIN-kaufen} ou \all{ein-KAU-fen} ?
\item Si l'accent tombe sur la \textbf{particule}, elle est séparable : \all{EINkaufen} $\rightarrow$ \all{Ich kaufe ein.}
\item Si l'accent tombe sur le \textbf{verbe}, la particule reste collée : \all{beZAHLEN} $\rightarrow$ \all{Ich bezahle.}
\item Astuce mémoire : les préfixes inséparables (\all{be-, ge-, er-, ver-, zer-, ent-, emp-}) ne sont jamais accentués -- apprends cette petite liste par cœur.
\end{enumerate}
\end{methode}

\courssub{La place de la particule dans la phrase}

Au présent, dans une proposition principale, la structure est rigide :

\begin{popfigure}
\begin{center}
\begin{tikzpicture}[
box/.style={draw=popblue, fill=popblueL, rounded corners=3pt, align=center, minimum height=0.9cm, font=\small},
pbox/.style={draw=poporange, fill=poporangeL, rounded corners=3pt, align=center, minimum height=0.9cm, font=\small\bfseries}]
\node[box, align=center] (s) at (0,0){SUJET\\ \all{Frau Berger}};
\node[box, right=0.5cm of s, align=center] (v){VERBE CONJUGUÉ\\ (2\textsuperscript{e} position)\\ \all{kauft}};
\node[box, right=0.5cm of v, align=center] (c){COMPLÉMENTS\\ \all{am Samstag}\\ \all{Obst und Gemüse}};
\node[pbox, right=0.5cm of c, align=center] (p){PARTICULE\\ (fin de phrase)\\ \all{ein.}};
\draw[-{Stealth}, thick, poporange] (v.south) to[bend right=25] node[below, font=\scriptsize, text=popdark]{la particule saute à la fin} (p.south);
\end{tikzpicture}
\legende{Structure de la phrase avec un verbe à particule séparable au présent.}
\end{center}
\end{popfigure}

\begin{exemplebox}[Exemples traduits]
\begin{itemize}
\item \all{Ich kaufe heute auf dem Markt ein.} -- « Je fais mes courses aujourd'hui au marché. »
\item \all{Er gibt viel Geld aus.} -- « Il dépense beaucoup d'argent. »
\item \all{Sie nimmt den Einkaufskorb mit.} -- « Elle emporte le panier. »
\item \all{Wir machen das Geschäft um neun Uhr auf.} -- « Nous ouvrons le magasin à neuf heures. »
\item \all{Der Kunde probiert die Jacke an.} -- « Le client essaie la veste. »
\item \all{Kommst du morgen zurück?} -- « Reviens-tu demain ? »
\item \all{Sie bezahlt die Rechnung.} -- « Elle paie la facture. » (inséparable : \all{be-})
\item \all{Der Verkäufer verkauft frische Ware.} -- « Le vendeur vend de la marchandise fraîche. » (inséparable : \all{ver-})
\end{itemize}
Remarque : même avec plusieurs compléments (temps, lieu), la particule reste \emph{dernière} : \all{Ich kaufe \textbf{heute auf dem Markt} ein.}
\end{exemplebox}

\courssub{Conjugaison complète au présent}

\begin{center}
\begin{tabular}{|l|l|l|}\hline
\rowcolor{popblueL} Personne & \all{einkaufen} (séparable) & \all{ausgeben} (séparable, fort) \\ \hline
\all{ich} & \all{kaufe ... ein} & \all{gebe ... aus} \\ \hline
\all{du} & \all{kaufst ... ein} & \all{gibst ... aus} \\ \hline
\all{er/sie/es} & \all{kauft ... ein} & \all{gibt ... aus} \\ \hline
\all{wir} & \all{kaufen ... ein} & \all{geben ... aus} \\ \hline
\all{ihr} & \all{kauft ... ein} & \all{gebt ... aus} \\ \hline
\all{sie/Sie} & \all{kaufen ... ein} & \all{geben ... aus} \\ \hline
\end{tabular}
\end{center}

\begin{attention}[Verbe fort + particule]
\all{ausgeben} est un verbe fort : la voyelle change (\all{e} $\rightarrow$ \all{i}) à la 2\textsuperscript{e} et 3\textsuperscript{e} personne du singulier : \all{du gibst aus, er gibt aus}. La séparation ne supprime pas l'irrégularité : on conjugue le verbe de base normalement, puis on envoie la particule à la fin.
\end{attention}

\courssub{Comparaison avec le français}

\begin{center}
\begin{tabularx}{\linewidth}{|X|X|X|}\hline
\rowcolor{popblueL} Français & Deutsch & Remarque \\ \hline
Je fais mes courses aujourd'hui. & \all{Ich kaufe heute ein.} & Un seul verbe en français, deux morceaux en allemand. \\ \hline
Il dépense beaucoup d'argent. & \all{Er gibt viel Geld aus.} & La particule modifie le sens : \all{geben} = donner, \all{ausgeben} = dépenser. \\ \hline
Elle paie la facture. & \all{Sie bezahlt die Rechnung.} & Préfixe inséparable \all{be-} : rien ne bouge. \\ \hline
Elle téléphone à sa fille. & \all{Sie ruft ihre Tochter an.} & Le français utilise un autre verbe ; l'allemand ajoute une particule. \\ \hline
\end{tabularx}
\end{center}

Le français ne connaît pas ce phénomène : nos verbes restent toujours en un seul bloc. Le phénomène le plus proche se trouve en \emph{anglais}, avec les verbes à particule (\emph{to give up} = abandonner, \emph{to take off} = décoller), où la particule change aussi le sens du verbe. En allemand, en plus, la particule se déplace physiquement en fin de phrase.

\begin{attention}[Erreurs fréquentes des francophones]
\begin{itemize}
\item Oublier la particule : *\all{Ich kaufe heute.} signifie seulement « j'achète aujourd'hui » -- le sens de « faire les courses » disparaît.
\item Coller la particule au verbe conjugué : *\all{Ich einkaufe heute.} est \textbf{faux}.
\item Séparer un préfixe inséparable : *\all{Sie zahlt die Rechnung be.} est impossible.
\item Au Bac, la place de la particule est \cle{systématiquement évaluée} au thème comme à la version : vérifie toujours la fin de tes phrases allemandes !
\end{itemize}
\end{attention}

\begin{exoresolu}[Conjugue et place la particule correctement]
Énoncé : mets les verbes entre parenthèses au présent, à la bonne place.
\begin{enumerate}
\item Frau Berger (einkaufen) am Samstag Obst.
\item Der Kunde (anprobieren) die Hose.
\item Wir (mitnehmen) den Korb zum Markt.
\item Sie (bezahlen) die Rechnung sofort.
\item Ihr (ausgeben) zu viel Geld!
\end{enumerate}
\tcblower
Solution détaillée :
\begin{enumerate}
\item \all{Frau Berger \textbf{kauft} am Samstag Obst \textbf{ein}.} -- \all{EINkaufen} est accentué sur la particule : séparation, \all{ein} en fin de phrase.
\item \all{Der Kunde \textbf{probiert} die Hose \textbf{an}.} -- \all{ANprobieren} : séparable ; le complément \all{die Hose} se place avant la particule.
\item \all{Wir \textbf{nehmen} den Korb zum Markt \textbf{mit}.} -- \all{MITnehmen} : séparable ; attention, \all{nehmen} est fort (\all{er nimmt}).
\item \all{Sie \textbf{bezahlt} die Rechnung sofort.} -- \all{be-} est inséparable : le verbe reste entier, rien à la fin.
\item \all{Ihr \textbf{gebt} zu viel Geld \textbf{aus}!} -- \all{AUSgeben} : séparable ; à \all{ihr}, pas de changement de voyelle (\all{gebt}), mais \all{du gibst / er gibt}.
\end{enumerate}
\end{exoresolu}

\begin{aretenir}[L'essentiel]
\begin{itemize}
\item Particule \cle{accentuée} (\all{ein-, aus-, auf-, mit-, an-, zurück-}) $\rightarrow$ séparable : verbe conjugué en 2\textsuperscript{e} position, particule en fin de phrase.
\item Préfixes \all{be-, ge-, er-, ver-, zer-, ent-, emp-} $\rightarrow$ inséparables, jamais accentués.
\item Particules variables (\all{durch-, über-, um-, wieder-}) : l'accent décide.
\item Les verbes forts gardent leur irrégularité : \all{er gibt aus, er nimmt mit}.
\item Test infaillible : prononce le verbe -- \all{EINkaufen} se sépare, \all{beZAHLEN} non.
\end{itemize}
\end{aretenir}

\begin{exercice}[À ton tour]
Traduis en allemand en faisant attention à la place de la particule :
\begin{enumerate}
\item Nous faisons nos courses le samedi au marché.
\item Le client dépense peu d'argent aujourd'hui.
\item Elle ouvre le magasin à huit heures.
\item J'emporte le panier et je paie la marchandise.
\item Revenez-vous demain, Madame ?
\end{enumerate}
\end{exercice}

\begin{corrige}[Exercice « À ton tour »]
\begin{enumerate}
\item \all{Wir kaufen am Samstag auf dem Markt ein.}
\item \all{Der Kunde gibt heute wenig Geld aus.}
\item \all{Sie macht das Geschäft um acht Uhr auf.}
\item \all{Ich nehme den Korb mit und bezahle die Ware.} (\all{mitnehmen} séparable, \all{bezahlen} inséparable.)
\item \all{Kommen Sie morgen zurück, Frau ...?}
\end{enumerate}
\end{corrige}

\begin{savaistu}[Pourquoi l'allemand sépare-t-il ses verbes ?]
Ce phénomène vient de l'histoire de la langue : les particules étaient autrefois de petits adverbes libres qui précisaient le sens du verbe (direction, achèvement). Avec le temps, elles se sont soudées au verbe à l'infinitif, mais ont gardé leur liberté dans la phrase conjuguée. C'est pourquoi les dictionnaires allemands notent toujours les verbes séparables avec un point médian ou un trait d'union : \all{ein·kaufen}, \all{aus·geben}. Au Cameroun, les élèves germanistes comparent souvent ce mécanisme à l'anglais \emph{phrasal verbs} qu'ils connaissent déjà -- un excellent pont entre les deux langues !
\end{savaistu}

\coursec{Expression guidée et méthode du commentaire}

Dans la section précédente, nous avons appris à conjuguer et à placer correctement les verbes à particule séparable : \all{einkaufen}, \all{mitnehmen}, \all{aufmachen}... Nous savons désormais dire \all{Ich kaufe heute auf dem Markt ein.} Cette compétence va nous servir immédiatement : il est temps de passer de la phrase isolée au \cle{texte}. Dans cette dernière section, vous allez apprendre à \textbf{commenter} le texte support, à le \textbf{résumer}, à \textbf{rédiger} un court paragraphe sur vos propres habitudes de consommation et à \textbf{jouer un dialogue} au magasin. C'est l'aboutissement de toute la leçon : comprendre, puis produire.

\courssub{La méthode du commentaire de texte}

Au lycée, commenter un texte allemand ne signifie pas le traduire : cela signifie le \cle{présenter}, en dégager les \cle{idées principales} et donner son \cle{avis personnel}. Un bon commentaire suit toujours trois étapes, comme en français.

\begin{methode}[Rédiger un commentaire de texte en trois parties]
\begin{enumerate}
\item \textbf{Introduction (Einleitung)} : présenter le texte en une ou deux phrases. Qui ? Quoi ? Où ? De quel thème s'agit-il ? Formule utile : \all{Der Text handelt von...} (le texte traite de...).
\item \textbf{Développement (Hauptteil)} : dégager deux ou trois idées principales et les appuyer par de courtes citations du texte entre guillemets allemands „...“. Utiliser les connecteurs : \all{zuerst} (d'abord), \all{dann} (ensuite), \all{außerdem} (de plus).
\item \textbf{Conclusion (Schluss)} : donner son avis personnel avec \all{meiner Meinung nach} (à mon avis) ou \all{ich finde, dass...} (je trouve que...), puis terminer par \all{schließlich} (finalement) ou une phrase de bilan.
\end{enumerate}
\end{methode}

\begin{popfigure}
\begin{center}
\begin{tikzpicture}[
  node distance=0.55cm,
  box/.style={draw=popblue, fill=popblueL, rounded corners, thick, align=center, text width=5.2cm, inner sep=6pt},
  arr/.style={-{Stealth[length=3mm]}, thick, popdark}
]
\node[box, align=center] (intro){\textbf{1. Einleitung}\\Présenter le texte et le thème\\\all{Der Text handelt von...}};
\node[box, below=of intro, fill=popgreenL, draw=popgreen, align=center] (dev){\textbf{2. Hauptteil}\\Idées principales + citations\\\all{zuerst... dann... außerdem...}};
\node[box, below=of dev, fill=poporangeL, draw=poporange, align=center] (concl){\textbf{3. Schluss}\\Avis personnel\\\all{meiner Meinung nach... schließlich...}};
\draw[arr] (intro) -- (dev);
\draw[arr] (dev) -- (concl);
\end{tikzpicture}
\end{center}
\legende{Organigramme du commentaire de texte : trois parties, trois fonctions.}
\end{popfigure}

\begin{attention}[La subordonnée : le verbe à la fin]
Quand vous donnez votre avis avec \all{weil} (parce que), \all{dass} (que) ou \all{obwohl} (bien que), le verbe conjugué allemand part \textbf{à la fin} de la subordonnée, contrairement au français :
\all{Ich kaufe gern auf dem Markt ein, weil die Preise dort billiger sind.} « J'aime faire mes courses au marché, parce que les prix y sont moins chers. »
Ne dites jamais \all{weil die Preise sind billiger} : c'est la faute la plus fréquente des francophones.
\end{attention}

\courssub{Commentaire modèle : „Familie Berger geht einkaufen“}

Voici un commentaire complet, rédigé selon la méthode, sur le texte support de la leçon. Observez les connecteurs soulignés par l'usage : ils donnent le rythme du texte.

\begin{textebox}[Kommentar : Familie Berger geht einkaufen]
Der Text handelt von der Familie Berger, die am Samstag einkaufen geht. Zuerst geht die Familie auf den Markt, weil das Obst und das Gemüse dort frischer und billiger sind. Dann kauft Frau Berger im Geschäft Brot und Milch, und Herr Berger bezahlt an der Kasse. Außerdem vergleicht die Familie die Preise: Der Markt ist billiger als das Geschäft, aber das Geschäft ist bequemer. Meiner Meinung nach ist der Text sehr interessant, weil er das Leben einer deutschen Familie zeigt. Ich finde, dass Einkaufen auf dem Markt auch in Kamerun normal ist. Schließlich lerne ich aus dem Text, dass man die Preise immer vergleichen muss.
\end{textebox}

\begin{exemplebox}[Glossaire du commentaire]
\begin{itemize}
\item \all{frisch} : frais (ici : \all{frischer}, plus frais, comparatif)
\item \all{bequem} : pratique, confortable (\all{bequemer} : plus pratique)
\item \all{vergleichen} : comparer (\all{die Familie vergleicht} : la famille compare)
\item \all{die Kasse, -n} : la caisse
\item \all{normal} : normal, courant
\end{itemize}
\end{exemplebox}

Analysons la structure : l'introduction (première phrase) présente le texte avec \all{Der Text handelt von...}. Le développement utilise \all{zuerst}, \all{dann}, \all{außerdem} et une citation implicite des idées du texte (le marché est moins cher, le magasin est plus pratique). La conclusion donne l'avis personnel avec \all{meiner Meinung nach} et \all{ich finde, dass...}, puis se termine avec \all{schließlich}. Remarquez aussi les deux comparatifs (\all{frischer und billiger}, \all{billiger als}, \all{bequemer}) et le verbe séparable bien placé : \all{die am Samstag einkaufen geht} (particule \all{ein} juste avant \all{geht} en fin de subordonnée).

\begin{exemplebox}[Deux phrases modèles à réutiliser]
\all{Meiner Meinung nach ist der Markt billiger als das Geschäft.} « À mon avis, le marché est moins cher que le magasin. »

\all{Ich kaufe gern auf dem Markt ein.} « J'aime faire mes courses au marché. »
\end{exemplebox}

\begin{savaistu}[Des expressions qui font la différence au Bac]
\all{Meiner Meinung nach} et \all{Ich finde, dass...} sont des expressions très valorisées à l'oral du Baccalauréat : elles montrent que vous savez donner un avis personnel argumenté, ce qui est exactement ce que les examinateurs attendent. Apprenez-les par cœur et placez-les dans chaque production, écrite comme orale. Attention à la construction : après \all{dass}, le verbe va à la fin : \all{Ich finde, dass der Markt billiger ist.}
\end{savaistu}

\courssub{Le résumé guidé (Zusammenfassung)}

Résumer, c'est dire l'essentiel en trois ou quatre phrases, \textbf{sans citer} et \textbf{sans donner son avis}. On garde uniquement : qui, quoi, où, et l'idée principale.

\begin{methode}[Résumer un texte en 3-4 phrases]
\begin{enumerate}
\item Repérer les personnages et le lieu : \all{Familie Berger, Markt, Geschäft}.
\item Noter 3 idées essentielles au brouillon, sans détails ni exemples.
\item Rédiger au présent, avec des phrases simples et des connecteurs (\all{zuerst}, \all{dann}, \all{schließlich}).
\item Vérifier : pas d'avis personnel, pas de citation, pas de détail inutile.
\end{enumerate}
\end{methode}

\begin{exemplebox}[Zusammenfassung modèle]
\all{Familie Berger geht am Samstag einkaufen. Zuerst kauft sie Obst und Gemüse auf dem Markt, dann Brot und Milch im Geschäft. Die Familie vergleicht die Preise: Der Markt ist billiger, das Geschäft ist bequemer.} « La famille Berger fait ses courses le samedi. D'abord elle achète des fruits et légumes au marché, puis du pain et du lait au magasin. La famille compare les prix : le marché est moins cher, le magasin est plus pratique. »
\end{exemplebox}

\begin{attention}[Résumé ou commentaire ?]
Ne confondez pas les deux exercices : dans le \textbf{résumé}, on ne donne \emph{jamais} son avis (pas de \all{meiner Meinung nach}) ; dans le \textbf{commentaire}, l'avis personnel est \emph{obligatoire} dans la conclusion. Lire la consigne avant de rédiger !
\end{attention}

\courssub{Production écrite guidée : „Meine Einkäufe“}

À votre tour ! Vous allez décrire vos habitudes de consommation en 6 à 8 phrases, en réutilisant le lexique du chapitre, au moins deux comparatifs et deux verbes à particule séparable.

\begin{methode}[Rédiger un court paragraphe en quatre étapes]
\begin{enumerate}
\item \textbf{Noter 3 idées} au brouillon : où j'achète, quoi, pourquoi (prix, qualité, habitude).
\item \textbf{Rédiger avec des connecteurs} : \all{zuerst}, \all{dann}, \all{außerdem}, \all{meiner Meinung nach}.
\item \textbf{Vérifier la place du verbe et des particules} : verbe en position 2 dans la phrase principale ; particule séparable à la fin ; verbe conjugué à la fin des subordonnées.
\item \textbf{Relire les déclinaisons} : adjectifs (\all{frisches Obst}, \all{die billigen Preise}), articles et cas (\all{auf dem Markt}, \all{in dem Geschäft}).
\end{enumerate}
\end{methode}

\begin{exemplebox}[Production modèle : Meine Einkäufe]
\all{Ich kaufe gern auf dem Markt in Yaoundé ein. Zuerst kaufe ich frisches Obst und Gemüse, weil die Preise dort billiger sind. Dann gehe ich in ein Geschäft und kaufe Brot und Reis. Der Markt ist billiger als das Geschäft, aber das Geschäft ist bequemer. Außerdem nehme ich immer eine Tasche mit. Meiner Meinung nach ist der Markt interessanter, weil man dort mit den Verkäufern spricht. Schließlich gebe ich nicht zu viel Geld aus.} « J'aime faire mes courses au marché de Yaoundé. D'abord j'achète des fruits et légumes frais, parce que les prix y sont moins chers. Ensuite je vais dans un magasin et j'achète du pain et du riz. Le marché est moins cher que le magasin, mais le magasin est plus pratique. De plus, j'apporte toujours un sac. À mon avis, le marché est plus intéressant, parce qu'on y parle avec les vendeurs. Finalement, je ne dépense pas trop d'argent. »
\end{exemplebox}

Vérifions le contrat : deux comparatifs (\all{billiger}, \all{bequemer}, \all{interessanter}), deux verbes séparables correctement placés (\all{ich kaufe ... ein}, \all{ich nehme ... mit}, \all{ich gebe ... aus}), des connecteurs, un avis personnel et des déclinaisons correctes (\all{frisches Obst}, \all{auf dem Markt}).

\courssub{Expression orale : dialogue guidé au magasin}

\begin{experience}[Jeu de rôle : Kunde und Verkäufer]
Par groupes de deux, jouez une scène au magasin ou au marché. L'élève A est le client (der Kunde), l'élève B est le vendeur (der Verkäufer). Utilisez obligatoirement les expressions rituelles suivantes, puis inversez les rôles.
\end{experience}

\begin{textebox}[Dialogue modèle : Im Geschäft]
Verkäufer: Guten Tag! Kann ich Ihnen helfen?\\
Kunde: Ja, bitte. Ich suche Brot und Milch.\\
Verkäufer: Das Brot ist dort links. Die Milch ist im Kühlschrank.\\
Kunde: Danke. Was kostet das Brot?\\
Verkäufer: Es kostet zwei Euro.\\
Kunde: Das ist billiger als auf dem Markt! Ich nehme zwei Brote.\\
Verkäufer: Gerne. Sonst noch etwas?\\
Kunde: Nein, danke, das ist alles.\\
Verkäufer: Das macht vier Euro, bitte.\\
Kunde: Hier, bitte. Auf Wiedersehen!\\
Verkäufer: Auf Wiedersehen und einen schönen Tag!
\end{textebox}

\begin{exemplebox}[Glossaire du dialogue]
\begin{itemize}
\item \all{Kann ich Ihnen helfen?} : puis-je vous aider ?
\item \all{ich suche...} : je cherche...
\item \all{Was kostet...?} : combien coûte... ?
\item \all{der Kühlschrank, die Kühlschränke} : le réfrigérateur
\item \all{Sonst noch etwas?} : et avec ceci ? (autre chose ?)
\item \all{Das macht vier Euro.} : cela fait quatre euros.
\end{itemize}
\end{exemplebox}

\begin{center}
\begin{tabularx}{\linewidth}{|X|X|X|}
\hline
\rowcolor{popblueL} Français & Deutsch & Remarque \\
\hline
Combien coûte le pain ? & Was kostet das Brot? & question rituelle d'achat \\
\hline
C'est moins cher qu'au marché. & Das ist billiger als auf dem Markt. & comparatif + \all{als} \\
\hline
Je prends deux pains. & Ich nehme zwei Brote. & \all{nehmen} : verbe fort \\
\hline
C'est tout, merci. & Das ist alles, danke. & formule de fin d'achat \\
\hline
À mon avis, c'est trop cher. & Meiner Meinung nach ist das zu teuer. & avis personnel \\
\hline
\end{tabularx}
\end{center}

\begin{exoresolu}[Compléter et transformer un dialogue]
Énoncé : complétez les répliques du client avec les éléments entre parenthèses, en respectant la place du verbe.
1. Guten Tag! Ich \dots\ Brot und Milch. (suchen)
2. Ich kaufe gern auf dem Markt \dots\ (einkaufen), weil die Preise dort \dots\ sind. (billig, comparatif)
3. Ich \dots\ zwei Brote \dots\ (mitnehmen).
\tcblower
Solution détaillée :
1. \all{Ich suche Brot und Milch.} — verbe régulier \all{suchen}, position 2.
2. \all{Ich kaufe gern auf dem Markt ein, weil die Preise dort billiger sind.} — particule \all{ein} à la fin de la phrase principale ; dans la subordonnée introduite par \all{weil}, le verbe \all{sind} va à la fin ; comparatif de \all{billig} : \all{billiger}.
3. \all{Ich nehme zwei Brote mit.} — verbe fort \all{mitnehmen} : \all{ich nehme}, particule \all{mit} renvoyée en fin de phrase.
\end{exoresolu}

\begin{exercice}[À vous de produire]
1. Rédigez votre propre paragraphe \all{Meine Einkäufe} (6-8 phrases) en suivant la méthode en quatre étapes.
2. Résumez le texte „Familie Berger geht einkaufen“ en 3 phrases, sans avis personnel.
3. Avec un camarade, écrivez puis jouez un dialogue au marché de Douala en utilisant au moins cinq expressions rituelles du tableau.
\end{exercice}

\begin{corrige}[Exercice n°2]
\all{Familie Berger geht am Samstag einkaufen. Sie kauft Obst und Gemüse auf dem Markt und Brot und Milch im Geschäft. Der Markt ist billiger, aber das Geschäft ist bequemer.} — Trois phrases, présent, aucun avis personnel, aucune citation : le contrat du résumé est respecté.
\end{corrige}

\courssub{Grille d'auto-évaluation}

Avant de rendre votre production, vérifiez chaque point et cochez mentalement :

\begin{center}
\begin{tabularx}{\linewidth}{|X|X|}
\hline
\rowcolor{popgreenL} Critère & Oui / Non \\
\hline
J'ai utilisé au moins 5 mots du vocabulaire du thème (\all{der Preis, das Geschäft, der Markt, die Ware, der Kunde, kaufen}) & \\
\hline
J'ai placé au moins 2 comparatifs (\all{billiger als}, \all{frischer}, \all{bequemer}) & \\
\hline
J'ai employé au moins 2 verbes séparables avec la particule à la fin (\all{einkaufen}, \all{mitnehmen}, \all{ausgeben}) & \\
\hline
Dans mes subordonnées (\all{weil}, \all{dass}), le verbe conjugué est à la fin & \\
\hline
J'ai donné mon avis avec \all{meiner Meinung nach} ou \all{ich finde, dass...} & \\
\hline
J'ai utilisé au moins 3 connecteurs (\all{zuerst}, \all{dann}, \all{außerdem}, \all{schließlich}) & \\
\hline
Mes adjectifs sont correctement déclinés (\all{frisches Obst}, \all{die billigen Preise}) & \\
\hline
\end{tabularx}
\end{center}

\begin{aretenir}[L'essentiel de la section]
Un commentaire allemand se construit en trois parties : introduction (\all{Der Text handelt von...}), développement avec idées et citations (\all{zuerst, dann, außerdem}), conclusion avec avis personnel (\all{meiner Meinung nach}, \all{schließlich}). Le résumé, lui, reste objectif et bref. Dans toutes vos productions, surveillez les trois points critiques du francophone : le verbe en position 2, la particule séparable à la fin de la phrase principale, et le verbe conjugué à la fin de toute subordonnée. Avec ces outils, vous savez désormais comprendre, commenter, résumer et produire sur le thème de la consommation : la leçon est accomplie.
\end{aretenir}

\newpage

\coursec{Activité d'intégration}

\courssub{Objectifs de l'activité}

À la fin de cette activité d'intégration, tu seras capable de :
\begin{itemize}
\item mobiliser le \cle{vocabulaire} de la consommation (\all{der Preis, die Ware, das Geschäft, der Kunde, die Dienstleistung, kaufen}) dans une situation de communication réelle ;
\item employer correctement le \cle{comparatif} et le \cle{superlatif} pour comparer des prix et des produits ;
\item appliquer la \cle{déclinaison de l'adjectif} dans des phrases authentiques ;
\item utiliser les \cle{verbes à particule séparable} (\all{einkaufen, mitnehmen, ausgeben, anprobieren}) avec la particule à la fin de la phrase ;
\item produire à l'écrit une affiche comparative et à l'oral un dialogue marchand--client.
\end{itemize}

\courssub{Situation}

Le club d'allemand de ton lycée à Yaoundé prépare une journée portes ouvertes. Le thème de l'année est : \all{„Einkaufen hier und dort — Konsum in Kamerun und in Deutschland“} (« Faire ses courses ici et là-bas : la consommation au Cameroun et en Allemagne »). Ton groupe doit présenter deux productions : une affiche intitulée \all{„Unsere Preise“} qui compare les prix de cinq produits au marché Mokolo de Yaoundé et dans un supermarché allemand, puis un petit jeu de rôles entre un client et un vendeur.

Voici le document de travail distribué par le professeur :

\begin{textebox}[Preise am Markt und im Supermarkt]
Familie Ngo Bell kauft jeden Samstag auf dem Markt in Yaoundé ein. Ihre Tochter Clara macht ein Austauschjahr in München. Sie schreibt: „Hier in Deutschland kaufe ich im Supermarkt ein. Die Preise sind höher als in Kamerun. Ein Kilo Tomaten kostet hier zwei Euro, aber in Yaoundé sind die Tomaten billiger. Die Mangos sind in Kamerun am billigsten und am frischesten! Hier ist das Brot billig, aber die tropischen Früchte sind teuer. Der Supermarkt ist modern und sauber, aber der Markt in Yaoundé ist lebendiger und die Verkäufer sind freundlicher. Man kann dort auch über den Preis sprechen — das geht im Supermarkt nicht!“
\end{textebox}

\begin{definition}[Glossaire du texte]
der Austausch, -e : l'échange (scolaire) ; das Austauschjahr, -e : l'année d'échange ; frisch : frais ; lebendig : vivant, animé ; die Verkäufer (Pl.) : les vendeurs ; über den Preis sprechen : discuter du prix, négocier ; das geht nicht : cela ne se fait pas.
\end{definition}

\courssub{Consignes}

Travaillez par groupes de 3 ou 4 élèves. Suivez les étapes dans l'ordre.

\textbf{Étape 1 — Compréhension du document (10 min).}
\begin{enumerate}
\item Lisez le texte de Clara à voix haute dans le groupe, chacun deux phrases.
\item Répondez par écrit, en allemand et en phrases complètes, aux questions suivantes :
\begin{enumerate}
\item \all{Wo kauft Familie Ngo Bell ein?}
\item \all{Wo kauft Clara in München ein?}
\item \all{Was ist in Kamerun am billigsten?}
\item \all{Was kann man auf dem Markt in Yaoundé machen, aber nicht im Supermarkt?}
\end{enumerate}
\end{enumerate}

\textbf{Étape 2 — L'affiche « Unsere Preise » (20 min).}
\begin{enumerate}
\setcounter{enumi}{2}
\item Choisissez cinq produits (par exemple : \all{die Tomaten, die Mangos, die Bananen, das Brot, der Reis}).
\item Sur une feuille A4, rédigez une affiche avec un tableau de prix imaginaires mais réalistes (francs CFA pour Yaoundé, euros pour Munich).
\item Écrivez sous le tableau au moins \textbf{cinq phrases comparatives} : deux avec le comparatif et \all{als}, une avec le superlatif (\all{am ...sten}), une avec un adjectif décliné devant un nom (ex. \all{die frischen Mangos}), une avec un verbe à particule séparable.
\end{enumerate}

\textbf{Étape 3 — Le jeu de rôles (15 min de préparation).}
\begin{enumerate}
\setcounter{enumi}{5}
\item Rédigez un dialogue de 10 à 14 répliques entre \all{der Kunde / die Kundin} et \all{der Verkäufer / die Verkäuferin}. Le dialogue doit contenir : les salutations, la demande de prix (\all{Was kostet...? / Wie viel kostet...?}), une comparaison de deux produits, une négociation, l'achat final, et au moins \textbf{trois verbes séparables} parmi : \all{einkaufen, mitnehmen, ausgeben, anprobieren, zurückgeben}.
\item Répartissez les rôles, répétez, puis jouez la scène devant la classe avec votre affiche comme décor.
\end{enumerate}

\textbf{Étape 4 — Évaluation croisée (10 min).}
\begin{enumerate}
\setcounter{enumi}{7}
\item Pendant les présentations des autres groupes, complétez la grille : vocabulaire du thème (2 points), comparatifs et superlatifs corrects (2 points), place des particules séparables (2 points), déclinaison des adjectifs (2 points), prononciation et fluidité (2 points).
\end{enumerate}

\courssub{Ressources à mobiliser}

\begin{aretenir}[Comparatif et superlatif : les formes à réutiliser]
\begin{itemize}
\item Comparatif : adjectif + \textbf{-er} + \all{als} : \all{Die Tomaten sind billiger als die Äpfel.} (« Les tomates sont moins chères que les pommes. »)
\item Superlatif : \all{am} + adjectif + \textbf{-sten} : \all{Die Mangos sind am billigsten.} (« Les mangues sont les moins chères. »)
\item Umlaut pour les adjectifs courts : \all{teuer $\to$ teurer}, mais \all{alt $\to$ älter $\to$ am ältesten}, \all{groß $\to$ größer $\to$ am größten}.
\item Irréguliers : \all{gut $\to$ besser $\to$ am besten} ; \all{viel $\to$ mehr $\to$ am meisten} ; \all{gern $\to$ lieber $\to$ am liebsten}.
\end{itemize}
\end{aretenir}

\begin{aretenir}[Déclinaison de l'adjectif et verbes séparables]
\begin{itemize}
\item Après l'article défini : \all{die frische\textbf{n} Mangos, der billige Reis, das deutsche Brot}.
\item Sans article : \all{frische\textbf{n} Mangos, billige\textbf{r} Reis}.
\item Verbe séparable conjugué : la particule va \textbf{en fin de phrase} : \all{Ich kaufe heute auf dem Markt \textbf{ein}.} / \all{Sie nimmt die Tomaten \textbf{mit}.} / \all{Wir geben 2000 CFA \textbf{aus}.}
\end{itemize}
\end{aretenir}

\begin{center}
\begin{tabularx}{\linewidth}{|X|X|X|}
\hline
\rowcolor{popblueL} Français & Deutsch & Remarque \\
\hline
Combien coûte... ? & \all{Wie viel kostet...? / Was kostet...?} & demande de prix \\
\hline
C'est trop cher ! & \all{Das ist zu teuer!} & négociation \\
\hline
Je prends les tomates. & \all{Ich nehme die Tomaten.} & achat \\
\hline
Vous désirez ? & \all{Was möchten Sie? / Bitte schön?} & le vendeur \\
\hline
C'est moins cher que... & \all{Das ist billiger als...} & comparaison \\
\hline
\end{tabularx}
\end{center}

\begin{attention}[Pièges à éviter]
Ne dis jamais \all{*Ich einkaufe} : le verbe est \all{ich kaufe ein} (particule en fin de phrase). Ne dis pas \all{*mehr billiger} : le comparatif allemand se forme avec \textbf{-er}, jamais avec \all{mehr}. Attention au faux ami : \all{bekommen} signifie « recevoir », pas « devenir ».
\end{attention}

\courssub{Proposition de production et corrigé commenté}

\begin{exoresolu}[Modèle d'affiche et de dialogue]
Énoncé : produis l'affiche « Unsere Preise » et un dialogue client--vendeur selon les consignes des étapes 2 et 3.
\tcblower
\textbf{Affiche modèle :}

\all{\textbf{UNSERE PREISE}}

\all{Markt Mokolo (Yaoundé): Tomaten 500 CFA, Mangos 300 CFA, Bananen 400 CFA, Brot 600 CFA, Reis 800 CFA. Supermarkt in München: Tomaten 2 Euro, Mangos 3 Euro, Bananen 1,50 Euro, Brot 1 Euro, Reis 2 Euro.}

\all{1. Die Tomaten sind auf dem Markt billiger als im Supermarkt.\\
2. Das Brot ist in Deutschland billiger als die tropischen Früchte.\\
3. Die Mangos sind in Kamerun am billigsten.\\
4. Die frischen Mangos schmecken besser als die alten Bananen.\\
5. Familie Ngo Bell kauft jeden Samstag auf dem Markt ein.}

\textbf{Dialogue modèle :}

\all{Verkäuferin: Guten Morgen! Was möchten Sie?\\
Kundin: Guten Morgen! Wie viel kosten die Tomaten?\\
Verkäuferin: Ein Kilo kostet 500 CFA.\\
Kundin: Und die Mangos?\\
Verkäuferin: Die Mangos sind billiger, nur 300 CFA. Sie sind am billigsten heute.\\
Kundin: Hmm, die Tomaten sind zu teuer! Die Bananen sind billiger als die Tomaten.\\
Verkäuferin: Aber die frischen Tomaten sind besser! Okay, für Sie: 400 CFA.\\
Kundin: Gut, ich nehme zwei Kilo. Ich kaufe heute für die ganze Familie ein.\\
Verkäuferin: Möchten Sie auch die süßen Mangos mitnehmen?\\
Kundin: Ja, ein Kilo bitte. Wie viel gebe ich insgesamt aus?\\
Verkäuferin: Das macht 1100 CFA.\\
Kundin: Hier, bitte schön. Auf Wiedersehen!\\
Verkäuferin: Danke schön, auf Wiedersehen!}

\textbf{Commentaire des choix de langue :}
\begin{itemize}
\item Comparatifs avec \all{als} : \all{billiger als die Tomaten} ; superlatif : \all{am billigsten} (adverbe, donc \all{am + -sten}).
\item Adjectifs déclinés devant le nom : \all{die frische\textbf{n} Tomaten} (article défini pluriel $\to$ \textbf{-en}), \all{die süße\textbf{n} Mangos}.
\item Particules séparables en fin de phrase : \all{kaufe ... ein}, \all{mitnehmen} (à l'infinitif après \all{möchten}, la particule reste collée), \all{gebe ... aus}.
\item Formules de politesse authentiques : \all{Was möchten Sie?}, \all{bitte schön}, \all{danke schön}.
\end{itemize}
\end{exoresolu}

\courssub{Bilan de l'activité}

Cette activité t'a permis de réinvestir l'ensemble des savoir-faire de la leçon dans une situation de communication proche de la réalité : comparer des prix comme le font chaque jour les consommateurs de Yaoundé ou de Munich, négocier comme au marché, et conclure un achat. Vérifie que tu sais désormais :
\begin{itemize}
\item citer au moins huit mots du lexique de la consommation avec leur article et leur pluriel ;
\item former un comparatif (\all{billiger als}) et un superlatif (\all{am billigsten}) sans faute ;
\item accorder un adjectif devant un nom (\all{die frischen Mangos}) ;
\item placer correctement la particule d'un verbe séparable en fin de phrase ;
\item tenir un dialogue simple d'achat en allemand, avec salutations, prix, comparaison et négociation.
\end{itemize}

\begin{experience}[Pour aller plus loin]
À la maison, rédige un court paragraphe de cinq à six phrases intitulé \all{„Mein Einkauf am Samstag“} (« Mes courses du samedi ») : décris ce que ta famille achète, où, à quel prix, en utilisant au moins un comparatif, un superlatif et deux verbes séparables. Tu le présenteras à l'oral au début du cours suivant.
\end{experience}

\newpage

\coursec{Méthodes et exercices résolus}

\courssub{Méthode 1 : Comprendre un texte sur la consommation (global puis détaillé)}

\begin{methode}[Lire et exploiter un texte de compréhension écrite]
\begin{enumerate}
\item \textbf{Première lecture globale} : lis le texte une fois sans dictionnaire. Repère le thème (ici : \all{Einkaufen und Dienstleistungen}), le type de texte (article, reportage, dialogue) et les mots déjà connus (\all{kaufen, der Preis, das Geschäft}).
\item \textbf{Repérage des mots-clés} : souligne les noms (majuscule !), les chiffres et les verbes. Les mots de la même famille t'aident : \all{kaufen -- der Käufer -- der Einkauf -- einkaufen}.
\item \textbf{Deuxième lecture détaillée} : pour chaque question, localise la phrase du texte qui contient la réponse. Ne réponds jamais « de mémoire ».
\item \textbf{Rédaction de la réponse} : réponds par une phrase complète en allemand, en reprenant les mots de la question. Respecte la place du verbe conjugué en 2\up{e} position.
\item \textbf{Vérification} : verbe conjugué à la 2\up{e} place, majuscules aux noms, bon article (\all{der Preis, die Ware, das Geschäft}).
\end{enumerate}
\end{methode}

\begin{exoresolu}[Compréhension de texte : questions sur « Einkaufen in Deutschland »]
\textbf{Énoncé.} Texte : \all{Viele Deutsche kaufen heute im Supermarkt oder im Internet ein. Die Preise im Internet sind oft billiger als im Geschäft. Aber viele Kunden gehen gern in kleine Läden, weil der Service dort besser ist. Auch Dienstleistungen wie Friseur oder Bank sind wichtig für den Alltag.}

Traduction du texte : « Beaucoup d'Allemands font aujourd'hui leurs courses au supermarché ou sur Internet. Les prix sur Internet sont souvent moins chers qu'en magasin. Mais beaucoup de clients aiment aller dans les petits commerces, parce que le service y est meilleur. Les services comme le coiffeur ou la banque sont aussi importants pour la vie quotidienne. »

\textbf{Beantworten Sie die Fragen auf Deutsch :}
\begin{enumerate}
\item \all{Wo kaufen viele Deutsche heute ein?}
\item \all{Warum gehen viele Kunden gern in kleine Läden?}
\item \all{Welche Dienstleistungen nennt der Text?}
\end{enumerate}
\tcblower
\textbf{Raisonnement et solution détaillée.}

\textbf{Question 1.} La question porte sur \all{wo} (où). On localise la première phrase : \all{Viele Deutsche kaufen heute im Supermarkt oder im Internet ein.} On reprend la structure de la question et on place le verbe conjugué en 2\up{e} position. Attention : \all{einkaufen} est un verbe à particule séparable : \all{kaufen ... ein}.

Réponse rédigée : \all{Viele Deutsche kaufen heute im Supermarkt oder im Internet ein.} « Beaucoup d'Allemands font leurs courses aujourd'hui au supermarché ou sur Internet. »

\textbf{Question 2.} Le mot-clé est \all{warum} (pourquoi) : la réponse se trouve dans la troisième phrase, introduite par \all{weil} (verbe conjugué en fin de proposition). On peut répondre avec \all{weil} ou avec \all{denn} (verbe en 2\up{e} position).

Réponse rédigée : \all{Viele Kunden gehen gern in kleine Läden, weil der Service dort besser ist.} « Beaucoup de clients aiment aller dans les petits commerces parce que le service y est meilleur. » Justification : après \all{weil}, le verbe conjugué \all{ist} se place à la fin de la subordonnée. Variante correcte : \all{..., denn der Service ist dort besser.}

\textbf{Question 3.} On cherche \all{Dienstleistungen} dans la dernière phrase : \all{Friseur oder Bank}.

Réponse rédigée : \all{Der Text nennt den Friseur und die Bank.} « Le texte cite le coiffeur et la banque. » Justification : après \all{nennen}, les noms sont à l'accusatif : \all{den Friseur} (masculin, article modifié), \all{die Bank} (féminin, forme identique au nominatif).

\textbf{Conclusion.} Chaque réponse reprend les mots de la question, forme une phrase complète et respecte la place du verbe : c'est exactement ce que le correcteur attend à l'examen.
\end{exoresolu}

\courssub{Méthode 2 : Employer le vocabulaire de la consommation}

\begin{methode}[Apprendre et réutiliser le lexique thématique]
\begin{enumerate}
\item \textbf{Apprends chaque nom avec son article et son pluriel} : \all{der Preis, -e ; das Geschäft, -e ; die Ware, -n ; der Kunde, -n ; die Dienstleistung, -en ; der Laden, "Läden}. Sans l'article, le nom est inutilisable à l'écrit.
\item \textbf{Regroupe par familles de mots} : \all{kaufen -- verkaufen -- einkaufen -- der Kauf -- der Käufer -- der Verkäufer}. Cela multiplie ton vocabulaire sans effort.
\item \textbf{Mémorise des expressions toutes faites} : \all{Wie viel kostet das?} (combien ça coûte ?), \all{Das ist zu teuer.} (c'est trop cher), \all{Ich möchte zahlen.} (je voudrais payer), \all{der gute Service} (le bon service).
\item \textbf{Réutilise le lexique dans tes propres phrases} : pour chaque mot nouveau, écris une phrase simple sur ton quotidien (marché de Yaoundé, supermarché, bendskin).
\item \textbf{Vérifie le genre} avant d'accorder l'adjectif ou l'article : \all{der billige Preis}, mais \all{das billige Geschäft}.
\end{enumerate}
\end{methode}

\begin{exoresolu}[Vocabulaire : compléter un dialogue au magasin]
\textbf{Énoncé.} Complétez le dialogue avec les mots suivants : \all{Preis -- Kunden -- Ware -- Geschäft -- Dienstleistung}.

\all{Verkäufer: Guten Tag! Kann ich Ihnen helfen?\\
Kundin: Ja, bitte. Wie hoch ist der (1) .......... dieser Jacke?\\
Verkäufer: Sie kostet nur 30 Euro. Unsere (2) .......... ist immer von guter Qualität.\\
Kundin: Das ist ein schönes (3) .......... hier, und der Service ist prima.\\
Verkäufer: Danke! Wir haben viele zufriedene (4) .......... . Ein guter Friseur ist übrigens auch eine wichtige (5) .......... in unserer Stadt.}
\tcblower
\textbf{Raisonnement et solution détaillée.}

\textbf{(1)} Après l'article défini \all{der} et la question \all{Wie hoch...?} (à combien s'élève...), il faut le nom masculin \all{Preis} : \all{Wie hoch ist der Preis dieser Jacke?} « Quel est le prix de cette veste ? »

\textbf{(2)} Le possessif \all{unsere} (féminin singulier ou pluriel) et le sens « de bonne qualité » appellent \all{Ware} (féminin) : \all{Unsere Ware ist immer von guter Qualität.} « Notre marchandise est toujours de bonne qualité. » Le verbe \all{ist} au singulier confirme le choix.

\textbf{(3)} \all{ein schönes} : adjectif avec terminaison \all{-es} après \all{ein} $\Rightarrow$ nom neutre : \all{Geschäft}. \all{Das ist ein schönes Geschäft hier.} « C'est un beau magasin ici. »

\textbf{(4)} \all{viele zufriedene} : pluriel avec adjectif en \all{-e} ; le sens « satisfaits » désigne des personnes : \all{Kunden}. \all{Wir haben viele zufriedene Kunden.} « Nous avons beaucoup de clients satisfaits. »

\textbf{(5)} \all{eine wichtige} : féminin singulier ; un coiffeur est un service : \all{Dienstleistung}. \all{Ein guter Friseur ist eine wichtige Dienstleistung.} « Un bon coiffeur est un service important. »

\textbf{Conclusion.} Pour choisir le bon mot, croise toujours deux indices : le sens de la phrase et la grammaire (article, terminaison de l'adjectif, nombre du verbe).
\end{exoresolu}

\courssub{Méthode 3 : Décliner correctement l'adjectif épithète}

\begin{methode}[Choisir la bonne terminaison de l'adjectif]
\begin{enumerate}
\item \textbf{Repère ce qui précède l'adjectif} : article défini (\all{der, die, das...}), article indéfini ou possessif (\all{ein, mein, kein...}), ou rien du tout.
\item \textbf{Après l'article défini} : terminaison \all{-e} au nominatif singulier (trois genres) et à l'accusatif féminin et neutre ; \all{-en} partout ailleurs. Ex. : \all{der billige Preis, die billige Ware, das billige Geschäft}, mais \all{den billigen Preis} (accusatif masculin).
\item \textbf{Après \all{ein} (indéfini/possessif)} : l'adjectif « montre » le genre que \all{ein} ne montre pas : \all{ein billiger Preis} (m.), \all{ein billiges Geschäft} (n.), \all{eine billige Ware} (f.).
\item \textbf{Sans article} : l'adjectif porte seul la marque du genre et du cas : \all{guter Service, gute Qualität, frisches Brot}.
\item \textbf{Adjectif après le verbe} (attribut) : JAMAIS de terminaison : \all{Der Preis ist hoch.} (et non \emph{hocher}).
\end{enumerate}
\end{methode}

\begin{exoresolu}[Grammaire : accorder les adjectifs]
\textbf{Énoncé.} Mettez l'adjectif entre parenthèses à la bonne forme.
\begin{enumerate}
\item \all{Der (teuer) .......... Pullover kostet 50 Euro.}
\item \all{Ich kaufe ein (modern) .......... Handy.}
\item \all{Die (klein) .......... Läden haben einen (gut) .......... Service.}
\item \all{Frisch .......... Brot schmeckt gut.}
\end{enumerate}
\tcblower
\textbf{Raisonnement et solution détaillée.}

\textbf{(1)} \all{Der teure Pullover kostet 50 Euro.} Après l'article défini \all{der}, nominatif masculin singulier : terminaison \all{-e}. « Le pull-over cher coûte 50 euros. »

\textbf{(2)} \all{Ich kaufe ein modernes Handy.} \all{das Handy} est neutre, ici accusatif ; \all{ein} ne marque pas le neutre, donc l'adjectif prend la marque forte \all{-es}. « J'achète un portable moderne. »

\textbf{(3)} \all{Die kleinen Läden haben einen guten Service.} Premier adjectif : pluriel après article défini \all{die} $\Rightarrow$ terminaison \all{-en} (\all{kleinen}). Deuxième adjectif : accusatif masculin après \all{einen} $\Rightarrow$ terminaison \all{-en} (\all{guten}). « Les petits commerces ont un bon service. »

\textbf{(4)} \all{Frisches Brot schmeckt gut.} Sans article devant un nom neutre (\all{das Brot}), l'adjectif prend la marque forte \all{-es}. « Le pain frais a bon goût. »

\textbf{Conclusion.} La question à te poser est toujours la même : quel article devant ? quel cas ? quel genre ? La terminaison en découle mécaniquement.
\end{exoresolu}

\courssub{Méthode 4 : Former et employer le comparatif et le superlatif}

\begin{methode}[Comparer les prix et les services]
\begin{enumerate}
\item \textbf{Comparatif} : adjectif + \all{-er} : \all{billig $\rightarrow$ billiger}, \all{teuer $\rightarrow$ teurer}. Pour comparer : \all{als} (et non « que ») : \all{Das Internet ist billiger als das Geschäft.}
\item \textbf{Superlatif attribut} (après \all{sein}) : \all{am} + adjectif + \all{-sten} : \all{Dieser Laden ist am billigsten.} « Ce magasin est le moins cher. »
\item \textbf{Superlatif épithète} (devant le nom) : adjectif + \all{-st} + terminaison d'adjectif : \all{der billigste Preis, das beste Angebot}.
\item \textbf{Irréguliers à connaître par cœur} : \all{gut -- besser -- am besten} ; \all{gern -- lieber -- am liebsten} ; \all{viel -- mehr -- am meisten} ; \all{hoch -- höher -- am höchsten}.
\item \textbf{Umlaut} pour les adjectifs monosyllabiques en a/o/u : \all{alt -- älter -- am ältesten}, \all{groß -- größer -- am größten}, \all{stark -- stärker -- am stärksten}. Mais pas d'Umlaut si la voyelle est déjà un Umlaut ou une diphtongue : \all{billig -- billiger}, \all{teuer -- teurer}.
\end{enumerate}
\end{methode}

\begin{exoresolu}[Grammaire : comparatifs et superlatifs]
\textbf{Énoncé.} Complétez avec la forme correcte de l'adjectif entre parenthèses.
\begin{enumerate}
\item \all{Im Supermarkt sind die Preise (billig) .......... als auf dem Markt.}
\item \all{Der Friseur in der Stadt ist (teuer) .......... als der Friseur im Dorf.}
\item \all{Dieses Geschäft hat das (gut) .......... Angebot.}
\item \all{Am Samstag ist der Markt (voll) .......... .}
\end{enumerate}
\tcblower
\textbf{Raisonnement et solution détaillée.}

\textbf{(1)} \all{Im Supermarkt sind die Preise billiger als auf dem Markt.} Comparatif : \all{billig + -er} ; la comparaison s'introduit par \all{als}. « Au supermarché, les prix sont moins chers qu'au marché. » Pas d'Umlaut : \all{billig} contient un \emph{i}.

\textbf{(2)} \all{Der Friseur in der Stadt ist teurer als der Friseur im Dorf.} \all{teuer + -er = teurer} (le \emph{e} central disparaît : on ne dit pas \emph{teuerer}). « Le coiffeur en ville est plus cher que le coiffeur au village. »

\textbf{(3)} \all{Dieses Geschäft hat das beste Angebot.} Superlatif irrégulier de \all{gut} : \all{best-} ; épithète devant un nom neutre avec article défini : terminaison \all{-e} $\Rightarrow$ \all{das beste Angebot}. « Ce magasin a la meilleure offre. »

\textbf{(4)} \all{Am Samstag ist der Markt am vollsten.} Superlatif attribut : \all{am + voll + -sten}. « Le samedi, le marché est le plus bondé. »

\textbf{Conclusion.} Retiens le trio : comparatif en \all{-er} + \all{als} ; superlatif attribut \all{am ... -sten} ; superlatif épithète \all{-st-} + terminaison d'adjectif.
\end{exoresolu}

\courssub{Méthode 5 : Conjuguer les verbes à particule séparable}

\begin{methode}[Placer la particule au bon endroit]
\begin{enumerate}
\item \textbf{Repère la particule séparable} : \all{ein-, auf-, an-, mit-, fern-...} (accentuée à l'oral : \emph{EINkaufen}).
\item \textbf{Phrase affirmative au présent} : verbe conjugué en 2\up{e} position, particule renvoyée À LA FIN : \all{Ich kaufe heute im Supermarkt ein.}
\item \textbf{Au Perfekt} : \all{-ge-} s'insère entre particule et radical : \all{eingekauft, aufgemacht, mitgebracht}. Ex. : \all{Wir haben gestern eingekauft.}
\item \textbf{Avec un modal ou à l'infinitif} : le verbe reste entier, à la fin : \all{Ich will heute einkaufen.}
\item \textbf{Subordonnée} : verbe conjugué entier à la fin : \all{..., weil ich heute einkaufe.}
\end{enumerate}
\end{methode}

\begin{exoresolu}[Thème guidé : traduire en allemand]
\textbf{Énoncé.} Traduisez en allemand.
\begin{enumerate}
\item « J'achète mes provisions au marché. » (faire les courses = \all{einkaufen})
\item « Le magasin ouvre à huit heures. » (\all{aufmachen})
\item « Hier, nous avons fait les courses au supermarché. »
\item « Je veux t'inviter au restaurant. » (\all{einladen})
\end{enumerate}
\tcblower
\textbf{Raisonnement et solution détaillée.}

\textbf{(1)} \all{Ich kaufe meine Lebensmittel auf dem Markt ein.} Verbe conjugué \all{kaufe} en 2\up{e} position ; particule \all{ein} à la toute fin de la phrase. « au marché » = \all{auf dem Markt} (datif de lieu). Erreur à éviter : \emph{Ich einkaufe...} — la particule doit se détacher.

\textbf{(2)} \all{Das Geschäft macht um acht Uhr auf.} \all{aufmachen} : \all{macht} en 2\up{e} position, \all{auf} en fin. L'heure s'exprime avec \all{um ... Uhr}.

\textbf{(3)} \all{Gestern haben wir im Supermarkt eingekauft.} Perfekt : auxiliaire \all{haben} (verbe d'action sans déplacement) + participe \all{eingekauft} avec \all{-ge-} inséré entre la particule et le radical. La phrase commence par \all{Gestern}, donc inversion sujet-verbe : \all{haben wir}.

\textbf{(4)} \all{Ich will dich ins Restaurant einladen.} Avec le modal \all{wollen}, l'infinitif reste entier (\all{einladen}) et se place à la fin. \all{ins Restaurant} = \all{in das} (accusatif de mouvement).

\textbf{Conclusion.} Le réflexe de survie : verbe conjugué en 2\up{e} position, particule à la fin ; si le verbe n'est pas conjugué (infinitif, subordonnée), il reste en un seul mot.
\end{exoresolu}

\courssub{Méthode 6 : Produire un court paragraphe guidé}

\begin{methode}[Rédiger 5 à 8 phrases sur ses habitudes de consommation]
\begin{enumerate}
\item \textbf{Plan simple} : où tu achètes (lieu) $\rightarrow$ ce que tu achètes (produits) $\rightarrow$ les prix (comparatif) $\rightarrow$ ton avis (service, qualité).
\item \textbf{Phrases modèles à réutiliser} : \all{Ich kaufe oft ... ein.} / \all{Die Preise sind dort ... als ...} / \all{Ich finde den Service ...} / \all{Am liebsten kaufe ich ...}
\item \textbf{Varie les débuts de phrase} (sujet, temps, lieu) mais garde le verbe en 2\up{e} position : \all{Am Wochenende gehe ich auf den Markt.}
\item \textbf{Relie les idées} avec \all{und, aber, denn, weil, auch}.
\item \textbf{Relis} : majuscules aux noms, terminaisons des adjectifs, particules séparées, verbe conjugué à la bonne place.
\end{enumerate}
\end{methode}

\begin{exoresolu}[Expression écrite : « Mes achats »]
\textbf{Énoncé.} Rédigez un texte de 5 à 7 phrases en allemand : « Wo kaufen Sie ein und warum? » (Où faites-vous vos achats et pourquoi ?). Utilisez au moins un comparatif et un verbe à particule séparable.
\tcblower
\textbf{Raisonnement.} On suit le plan : lieu $\rightarrow$ produits $\rightarrow$ comparaison de prix $\rightarrow$ avis. On vérifie chaque phrase : verbe en 2\up{e} position, particule en fin, adjectifs accordés.

\textbf{Production corrigée :}

\all{Ich kaufe meine Lebensmittel oft auf dem Markt in Yaoundé ein. Dort finde ich frisches Gemüse und gutes Obst. Die Preise auf dem Markt sind billiger als im Supermarkt. Aber im Supermarkt ist der Service schneller und die Geschäfte sind moderner. Am liebsten kaufe ich Brot und Früchte. Ich finde das Einkaufen auf dem Markt interessant, weil man mit den Verkäufern sprechen kann.}

Traduction : « J'achète souvent mes provisions au marché de Yaoundé. Là, je trouve des légumes frais et de bons fruits. Les prix au marché sont moins chers qu'au supermarché. Mais au supermarché, le service est plus rapide et les magasins sont plus modernes. J'aime surtout acheter du pain et des fruits. Je trouve que faire les courses au marché est intéressant parce qu'on peut parler avec les vendeurs. »

\textbf{Justifications grammaticales :} \all{kaufe ... ein} (particule séparée) ; \all{billiger als, schneller, moderner} (comparatifs) ; \all{frisches Gemüse, gutes Obst} (adjectifs sans article, marque forte \all{-es}) ; \all{weil man ... sprechen kann} (verbe conjugué \all{kann} à la fin de la subordonnée, infinitif \all{sprechen} juste avant).

\textbf{Conclusion.} Un paragraphe court mais correct vaut mieux qu'un long texte fautif : privilégie des structures sûres.
\end{exoresolu}

\begin{attention}[Les 5 erreurs qui coûtent des points au Bac]
\begin{enumerate}
\item \textbf{Oublier de séparer la particule} : ne dis jamais \emph{Ich einkaufe auf dem Markt}. Écris \all{Ich kaufe auf dem Markt ein.} Le verbe conjugué en 2\up{e} position, la particule à la fin.
\item \textbf{Laisser l'adjectif sans terminaison devant le nom} : \emph{der billig Preis} est faux. Écris \all{der billige Preis}, \all{ein billiger Preis}. L'adjectif épithète se décline toujours ; l'adjectif attribut (\all{Der Preis ist billig.}) jamais.
\item \textbf{Calquer le français « que » par \all{wie} au lieu de \all{als}} : « plus cher que » = \all{teurer als}. \all{wie} ne sert que pour l'égalité : \all{so teuer wie} (aussi cher que).
\item \textbf{Oublier la majuscule et l'article des noms} : \all{der Preis, die Ware, das Geschäft, die Dienstleistung, der Kunde}. Un nom sans majuscule ou avec le mauvais article est sanctionné deux fois.
\item \textbf{Faux amis et calques} : \all{bekommen} signifie « recevoir », pas « devenir » (\all{werden}) ; « faire les courses » = \all{einkaufen} (et non un calque du français) ; « le prix est bon marché » se dit \all{Der Preis ist billig / günstig}, jamais un calque de « bon marché ».
\end{enumerate}
\end{attention}

\newpage

\coursec{Exercices}

\courssub{Je vérifie mes connaissances}

\begin{exercice}[QCM : le vocabulaire de la consommation]
Entoure la bonne réponse.
\begin{enumerate}
\item \all{der Preis} signifie :
\begin{enumerate}
\item le magasin \item le prix \item le client
\end{enumerate}
\item \all{die Dienstleistung} signifie :
\begin{enumerate}
\item la marchandise \item la facture \item la prestation de service
\end{enumerate}
\item Le pluriel de \all{die Ware} est :
\begin{enumerate}
\item die Waren \item die Wares \item die Wärer
\end{enumerate}
\item \all{einkaufen} signifie :
\begin{enumerate}
\item vendre \item faire ses courses \item payer
\end{enumerate}
\item \all{der Kunde} signifie :
\begin{enumerate}
\item le vendeur \item le client \item le prix
\end{enumerate}
\end{enumerate}
\end{exercice}

\begin{exercice}[Vrai ou faux ?]
Réponds par \emph{vrai} ou \emph{faux} et justifie ta réponse par une phrase correcte en allemand.
\begin{enumerate}
\item Dans \all{Ich kaufe heute auf dem Markt ein.}, la particule \all{ein-} se place à la fin de la phrase.
\item On écrit : \all{der billig Preis}.
\item Le superlatif de \all{gut} est \all{am gutsten}.
\item \all{teurer als} correspond à « plus cher que ».
\item Dans \all{Sie nimmt das Brot mit.}, \all{mitnehmen} est un verbe à particule séparable.
\end{enumerate}
\end{exercice}

\begin{exercice}[Vocabulaire : familles de mots]
Complète le tableau avec le mot allemand correct (nom avec article et pluriel, ou verbe).
\begin{center}
\begin{tabularx}{\linewidth}{|X|X|}
\hline
\rowcolor{popblueL} Français & Deutsch \\
\hline
acheter & \ldots \\
\hline
le magasin (les magasins) & \ldots \\
\hline
la marchandise (les marchandises) & \ldots \\
\hline
le client (les clients) & \ldots \\
\hline
dépenser (de l'argent) & \ldots \\
\hline
bon marché / cher & \ldots \\
\hline
\end{tabularx}
\end{center}
\end{exercice}

\begin{exercice}[Conjugaison : verbes à particule séparable]
Conjugue le verbe entre parenthèses au \all{Präsens} et place la particule correctement.
\begin{enumerate}
\item Ich \ldots\ heute auf dem Markt \ldots\ (einkaufen).
\item Sie \ldots\ zu viel Geld \ldots\ (ausgeben).
\item \ldots\ du das Brot \ldots\ ? (mitnehmen)
\item Der Laden \ldots\ um 18 Uhr \ldots\ (zumachen).
\item Wir \ldots\ um 8 Uhr in die Stadt \ldots\ (abfahren).
\end{enumerate}
\end{exercice}

\courssub{Je m'entraîne}

\begin{exercice}[Texte à trous : les courses de Lena]
Complète le texte avec les verbes séparables conjugués au \all{Präsens} : \all{einkaufen}, \all{ausgeben}, \all{mitnehmen}, \all{anfangen}, \all{zurückkommen}. Attention à la place de la particule.
\begin{textebox}[Einkaufen am Samstag]
Am Samstag \ldots\ Lena früh in der Stadt \ldots\ . Zuerst \ldots\ sie auf dem Markt \ldots\ : Sie kauft Obst, Brot und Käse. Sie \ldots\ nicht viel Geld \ldots\ , denn die Preise sind billig. Dann \ldots\ sie eine Tüte für den Kuchen \ldots\ . Um 12 Uhr \ldots\ sie nach Hause \ldots\ .
\end{textebox}
\end{exercice}

\begin{exercice}[Déclinaison de l'adjectif]
Complète les terminaisons des adjectifs (attention au cas, au genre et à l'article).
\begin{enumerate}
\item der billig\ldots\ Preis (Nominativ)
\item ein teuer\ldots\ Radio (Akkusativ)
\item gut\ldots\ Waren (Akkusativ, Pluriel, sans article)
\item mit dem neu\ldots\ Auto (Dativ)
\item die frisch\ldots\ Brötchen (Nominativ, Pluriel)
\item ein klein\ldots\ Geschäft (Nominativ, Neutrum)
\item wegen des hoch\ldots\ Preises (Genitiv)
\item Ich kaufe das billig\ldots\ Brot. (Akkusativ)
\end{enumerate}
\end{exercice}

\begin{exercice}[Comparatif et superlatif : transformations]
Transforme chaque phrase selon le modèle : \all{Das Brot ist billig.} $\rightarrow$ \all{Die Waren sind billiger.} $\rightarrow$ \all{Das Angebot ist am billigsten.}
\begin{enumerate}
\item Das Radio ist teuer. (der Fernseher / das Handy)
\item Der Markt ist groß. (das Geschäft / das Kaufhaus)
\item Das Brot ist gut. (der Käse / der Kuchen)
\item Der Weg ist lang. (die Straße / die Reise)
\item Der Laden ist nah. (der Supermarkt / die Bäckerei)
\item Die Tasche ist schön. (der Rucksack / die Uhr)
\end{enumerate}
\end{exercice}

\begin{exercice}[Appariements : phrases à traduire]
Traduis les phrases suivantes en allemand. Utilise le vocabulaire du chapitre et fais attention à la place du verbe.
\begin{enumerate}
\item Le client achète les meilleures marchandises.
\item Elle dépense trop d'argent.
\item Ce magasin est plus cher que le marché.
\item Nous faisons nos courses le samedi.
\item Le pain est moins bon que le gâteau.
\end{enumerate}
\end{exercice}

\courssub{Je me prépare au Bac}

\begin{exercice}[Texte de compréhension : la consommation en Suisse]
Lis le texte suivant, puis réponds aux questions \textbf{en allemand, par des phrases complètes}.
\begin{textebox}[Einkaufen in der Schweiz]
In der Schweiz kaufen viele Familien ihre Lebensmittel in großen Supermärkten. Die Preise sind dort oft höher als in Deutschland oder in Frankreich. Darum fahren viele Schweizer am Wochenende über die Grenze und kaufen dort ein. Das Brot, der Käse und das Fleisch sind in den Nachbarländern billiger. Aber die Schweizer Geschäfte bieten gute Dienstleistungen: Die Verkäufer sind freundlich, und die Waren sind frisch. Viele Kunden sagen: „Die Qualität ist wichtiger als der Preis.“ Auch die kleinen Läden im Dorf haben ihre Kunden. Sie öffnen früh am Morgen und schließen spät am Abend. Am Sonntag aber bleiben fast alle Geschäfte in der Schweiz zu. Dann machen die Familien einen Spaziergang oder besuchen Freunde. (environ 130 mots)
\end{textebox}
\textbf{Fragen zum Text :}
\begin{enumerate}
\item Wo kaufen viele Schweizer Familien ihre Lebensmittel?
\item Warum fahren viele Schweizer am Wochenende über die Grenze?
\item Welche Dienstleistungen bieten die Schweizer Geschäfte?
\item Was ist für viele Kunden wichtiger als der Preis?
\item Was machen die Familien am Sonntag?
\end{enumerate}
\textbf{Fragen zur Sprache :}
\begin{enumerate}
\item Relève dans le texte deux verbes à particule séparable et recopie la phrase complète.
\item Relève un comparatif et un adjectif décliné, puis explique la terminaison de l'adjectif.
\end{enumerate}
\end{exercice}

\begin{exercice}[Thème et version]
\textbf{A. Thème.} Traduis en allemand (attention : déclinaison des adjectifs, comparatifs, verbes séparables).
\begin{enumerate}
\item Le nouveau magasin vend les meilleures marchandises de la ville.
\item Ma sœur dépense chaque semaine beaucoup d'argent au supermarché.
\item Ce pain est meilleur que le pain du marché, mais il est aussi plus cher.
\item Le client prend le fromage frais et le lait froid.
\item Nous faisons nos courses le samedi matin et nous revenons à midi.
\end{enumerate}
\textbf{B. Version.} Traduis en français.
\begin{textebox}[Am Kiosk]
Jeden Morgen kauft Herr Mbarga eine Zeitung am Kiosk. Die Zeitung ist nicht teuer, aber die Zeitschriften sind teurer. „Ich nehme die billigste Zeitung“, sagt er und gibt dem Verkäufer das Geld. Dann fährt er mit dem Bus zur Arbeit.
\end{textebox}
\end{exercice}

\begin{exercice}[Production écrite guidée : Meine Einkäufe am Wochenende]
Rédige un texte allemand de \textbf{8 à 10 lignes} (environ 60 à 80 mots) intitulé \all{Meine Einkäufe am Wochenende}.
Consignes obligatoires :
\begin{enumerate}
\item Raconte où tu fais tes courses (marché, supermarché, petit magasin à Yaoundé ou à Douala) et ce que tu achètes.
\item Emploie \textbf{au moins deux comparatifs ou superlatifs} (par exemple : \all{billiger als}, \all{am besten}).
\item Emploie \textbf{au moins deux verbes à particule séparable} conjugués, avec la particule à la bonne place (par exemple : \all{einkaufen}, \all{mitnehmen}, \all{ausgeben}).
\item Emploie \textbf{au moins deux adjectifs déclinés} devant un nom (par exemple : \all{das frische Brot}).
\end{enumerate}
Tu peux t'aider du plan suivant : 1. Wann und wo? 2. Was kaufst du? 3. Wie sind die Preise und die Qualität? 4. Deine Meinung.
\end{exercice}

\newpage

\coursec{Corrigés des exercices}

\begin{corrige}[Exercice 1 — QCM : le vocabulaire de la consommation]
\begin{enumerate}
\item \all{der Preis} signifie : \textbf{b) le prix}.
\item \all{die Dienstleistung} signifie : \textbf{c) la prestation de service}.
\item Le pluriel de \all{die Ware} est : \textbf{a) die Waren}.
\item \all{einkaufen} signifie : \textbf{b) faire ses courses}.
\item \all{der Kunde} signifie : \textbf{b) le client}.
\end{enumerate}
\textbf{Point de vigilance :} ne pas confondre \all{der Kunde} (le client) et \all{der Verkäufer} (le vendeur) ; \all{die Ware} prend un \all{-n} au pluriel, comme beaucoup de noms féminins. Attention aussi au genre : on dit \all{der Preis} (masculin) et non \all{die Preis}, alors que le français « le prix » ne donne aucun indice sur le genre allemand.
\end{corrige}

\begin{corrige}[Exercice 2 — Vrai ou faux ?]
\begin{enumerate}
\item \textbf{Vrai.} \all{Ich kaufe heute auf dem Markt ein.} — Le verbe conjugué \all{kaufe} est en deuxième position et la particule \all{ein-} va à la fin de la phrase.
\item \textbf{Faux.} On écrit : \all{der billige Preis}. — Après l'article défini \all{der} au nominatif masculin, l'adjectif prend la terminaison \all{-e} (déclinaison faible).
\item \textbf{Faux.} Le superlatif de \all{gut} est \all{am besten} (comparatif : \all{besser}). C'est un adjectif irrégulier, à apprendre par cœur.
\item \textbf{Vrai.} \all{teurer als} = « plus cher que » ; le comparatif se forme avec \all{-er} + \all{als} (attention à l'élision : \all{teuer} $\rightarrow$ \all{teurer}, le \all{-e-} disparaît).
\item \textbf{Vrai.} \all{Sie nimmt das Brot mit.} — \all{mitnehmen} est séparable : \all{nimmt} en deuxième position, \all{mit} à la fin.
\end{enumerate}
\end{corrige}

\begin{corrige}[Exercice 3 — Vocabulaire : familles de mots]
\begin{center}
\begin{tabularx}{\linewidth}{|X|X|}
\hline
\rowcolor{popblueL} Français & Deutsch \\
\hline
acheter & \all{kaufen} \\
\hline
le magasin (les magasins) & \all{das Geschäft, die Geschäfte} \\
\hline
la marchandise (les marchandises) & \all{die Ware, die Waren} \\
\hline
le client (les clients) & \all{der Kunde, die Kunden} \\
\hline
dépenser (de l'argent) & \all{Geld ausgeben} \\
\hline
bon marché / cher & \all{billig / teuer} \\
\hline
\end{tabularx}
\end{center}
\textbf{Point de vigilance :} \all{ausgeben} est un verbe à particule séparable : \all{Ich gebe viel Geld aus.} « Je dépense beaucoup d'argent. » Ne pas confondre avec \all{geben} seul, qui signifie « donner ».
\end{corrige}

\begin{corrige}[Exercice 4 — Conjugaison : verbes à particule séparable]
\begin{enumerate}
\item \all{Ich kaufe heute auf dem Markt ein.} « Je fais mes courses au marché aujourd'hui. »
\item \all{Sie gibt zu viel Geld aus.} « Elle dépense trop d'argent. » (verbe fort : \all{geben} $\rightarrow$ \all{gibt})
\item \all{Nimmst du das Brot mit?} « Tu prends le pain avec toi ? » (verbe fort : \all{nehmen} $\rightarrow$ \all{nimmst} ; à l'interrogative sans mot interrogatif, le verbe conjugué passe en première position, la particule reste à la fin.)
\item \all{Der Laden macht um 18 Uhr zu.} « Le magasin ferme à 18 heures. »
\item \all{Wir fahren um 8 Uhr in die Stadt ab.} « Nous partons en ville à 8 heures. »
\end{enumerate}
\textbf{Règle :} au \all{Präsens}, dans une phrase affirmative, le verbe conjugué occupe la deuxième position et la particule séparable se détache pour aller en fin de phrase. C'est la règle capitale de ce chapitre : la particule « encadre » la phrase avec le verbe.
\end{corrige}

\begin{corrige}[Exercice 5 — Texte à trous : les courses de Lena]
\begin{textebox}[Einkaufen am Samstag — corrigé]
Am Samstag \textbf{fängt} Lena früh in der Stadt \textbf{an}. Zuerst \textbf{kauft} sie auf dem Markt \textbf{ein} : Sie kauft Obst, Brot und Käse. Sie \textbf{gibt} nicht viel Geld \textbf{aus}, denn die Preise sind billig. Dann \textbf{nimmt} sie eine Tüte für den Kuchen \textbf{mit}. Um 12 Uhr \textbf{kommt} sie nach Hause \textbf{zurück}.
\end{textebox}
Traduction : « Le samedi, Lena commence tôt en ville. D'abord, elle fait ses courses au marché : elle achète des fruits, du pain et du fromage. Elle ne dépense pas beaucoup d'argent, car les prix sont bas. Ensuite, elle emporte un sachet pour le gâteau. À midi, elle rentre à la maison. »
\textbf{Points de vigilance :} \all{anfangen} $\rightarrow$ \all{fängt an} (verbe fort, Umlaut à la 2\up{e} et 3\up{e} personne du singulier) ; \all{ausgeben} $\rightarrow$ \all{gibt aus} ; \all{mitnehmen} $\rightarrow$ \all{nimmt mit} ; \all{zurückkommen} $\rightarrow$ \all{kommt zurück}. Remarquer aussi : quand la phrase commence par un complément (\all{Am Samstag}, \all{Um 12 Uhr}), le sujet passe après le verbe (inversion).
\end{corrige}

\begin{corrige}[Exercice 6 — Déclinaison de l'adjectif]
\begin{enumerate}
\item \all{der billig\textbf{e} Preis} — nominatif masculin après article défini : terminaison \all{-e}.
\item \all{ein teuer\textbf{es} Radio} — accusatif neutre après \all{ein} : terminaison \all{-es} (déclinaison mixte).
\item \all{gut\textbf{e} Waren} — accusatif pluriel sans article : terminaison \all{-e} (déclinaison forte).
\item \all{mit dem neu\textbf{en} Auto} — datif après article défini : toujours \all{-en}.
\item \all{die frisch\textbf{en} Brötchen} — nominatif pluriel après article défini : \all{-en}.
\item \all{ein klein\textbf{es} Geschäft} — nominatif neutre après \all{ein} : \all{-es}.
\item \all{wegen des hoch\textbf{en} Preises} — génitif masculin après article défini : \all{-en} (et le nom prend \all{-es} : \all{des Preises}).
\item \all{Ich kaufe das billig\textbf{e} Brot.} — accusatif neutre après article défini : \all{-e}.
\end{enumerate}
\textbf{Astuce :} après \all{der/die/das} au singulier, l'adjectif prend \all{-e} ou \all{-en} ; après \all{ein}, il doit « montrer » le genre que l'article ne montre pas : \all{-er} (masculin), \all{-e} (féminin), \all{-es} (neutre) au nominatif. Au datif et au génitif après article défini, c'est toujours \all{-en} : c'est le réflexe le plus sûr.
\end{corrige}

\begin{corrige}[Exercice 7 — Comparatif et superlatif : transformations]
\begin{enumerate}
\item \all{Das Radio ist teuer.} $\rightarrow$ \all{Der Fernseher ist teurer.} $\rightarrow$ \all{Das Handy ist am teuersten.}
\item \all{Der Markt ist groß.} $\rightarrow$ \all{Das Geschäft ist größer.} $\rightarrow$ \all{Das Kaufhaus ist am größten.}
\item \all{Das Brot ist gut.} $\rightarrow$ \all{Der Käse ist besser.} $\rightarrow$ \all{Der Kuchen ist am besten.}
\item \all{Der Weg ist lang.} $\rightarrow$ \all{Die Straße ist länger.} $\rightarrow$ \all{Die Reise ist am längsten.}
\item \all{Der Laden ist nah.} $\rightarrow$ \all{Der Supermarkt ist näher.} $\rightarrow$ \all{Die Bäckerei ist am nächsten.}
\item \all{Die Tasche ist schön.} $\rightarrow$ \all{Der Rucksack ist schöner.} $\rightarrow$ \all{Die Uhr ist am schönsten.}
\end{enumerate}
\textbf{Points de vigilance :} Umlaut au comparatif et au superlatif pour les adjectifs monosyllabiques en \all{a/o/u} : \all{groß} $\rightarrow$ \all{größer} $\rightarrow$ \all{am größten}, \all{lang} $\rightarrow$ \all{länger} $\rightarrow$ \all{am längsten}. Irréguliers : \all{gut} $\rightarrow$ \all{besser} $\rightarrow$ \all{am besten} ; \all{nah} $\rightarrow$ \all{näher} $\rightarrow$ \all{am nächsten}. Le superlatif épithète s'écrit \all{am + -sten} (ou \all{-esten} après \all{s, ß, z, sch} pour faciliter la prononciation : \all{am größten}, \all{am frischesten}).
\end{corrige}

\begin{corrige}[Exercice 8 — Appariements : phrases à traduire]
\begin{enumerate}
\item \all{Der Kunde kauft die besten Waren.} (« les meilleures » = superlatif décliné : pluriel accusatif après article défini $\rightarrow$ \all{-en}.)
\item \all{Sie gibt zu viel Geld aus.} (verbe séparable : particule \all{aus} en fin de phrase.)
\item \all{Dieses Geschäft ist teurer als der Markt.} (comparatif + \all{als}.)
\item \all{Wir kaufen am Samstag ein.} (verbe séparable \all{einkaufen} ; le complément de temps se place avant la particule.)
\item \all{Das Brot ist weniger gut als der Kuchen.} (« moins ... que » = \all{weniger ... als}.)
\end{enumerate}
\textbf{Barème indicatif :} 1 point par phrase (vocabulaire 0,5 ; grammaire et place des mots 0,5).
\end{corrige}

\begin{corrige}[Exercice 9 — Texte de compréhension : la consommation en Suisse]
\textbf{Fragen zum Text — réponses modèles :}
\begin{enumerate}
\item \all{Viele Schweizer Familien kaufen ihre Lebensmittel in großen Supermärkten.} « Beaucoup de familles suisses achètent leurs produits alimentaires dans de grands supermarchés. »
\item \all{Sie fahren über die Grenze, weil die Preise dort billiger sind.} (ou : \all{..., denn das Brot, der Käse und das Fleisch sind in den Nachbarländern billiger.} — remarquer : après \all{weil}, le verbe conjugué va à la fin ; après \all{denn}, il reste en deuxième position.)
\item \all{Die Schweizer Geschäfte bieten gute Dienstleistungen: Die Verkäufer sind freundlich, und die Waren sind frisch.}
\item \all{Für viele Kunden ist die Qualität wichtiger als der Preis.}
\item \all{Am Sonntag machen die Familien einen Spaziergang oder besuchen Freunde.}
\end{enumerate}
\textbf{Fragen zur Sprache — réponses modèles :}
\begin{enumerate}
\item Verbes à particule séparable relevés dans le texte : \all{... und kaufen dort ein.} (\all{einkaufen}) ; \all{Am Sonntag aber bleiben fast alle Geschäfte in der Schweiz zu.} (\all{zu bleiben}, qui s'écrit en deux mots mais fonctionne comme un verbe à particule : la particule \all{zu} va en fin de phrase). Dans les deux cas, la particule se trouve en fin de phrase.
\item Comparatif : \all{billiger} (\all{... sind in den Nachbarländern billiger}) ou \all{wichtiger als}. Adjectifs déclinés : \all{großen Supermärkten} (datif pluriel après \all{in} + article défini pluriel $\rightarrow$ \all{-en}) ; \all{gute Dienstleistungen} (accusatif pluriel sans article $\rightarrow$ \all{-e}) ; \all{die kleinen Läden} (nominatif pluriel après article défini $\rightarrow$ \all{-en}).
\end{enumerate}
\textbf{Barème indicatif :} questions de compréhension 5 $\times$ 1 pt ; questions de langue 2 $\times$ 1,5 pt. On sanctionne la phrase incomplète, le verbe mal placé ou la réponse qui ne reprend pas les éléments de la question.
\end{corrige}

\begin{corrige}[Exercice 10 — Thème et version]
\textbf{A. Thème — traductions modèles :}
\begin{enumerate}
\item \all{Das neue Geschäft verkauft die besten Waren der Stadt.}
\item \all{Meine Schwester gibt jede Woche viel Geld im Supermarkt aus.}
\item \all{Dieses Brot ist besser als das Brot vom Markt, aber es ist auch teurer.}
\item \all{Der Kunde nimmt den frischen Käse und die kalte Milch.}
\item \all{Wir kaufen am Samstagmorgen ein und kommen um Mittag zurück.}
\end{enumerate}
\textbf{Points de vigilance :} phrase 1 : \all{neue} (nominatif neutre après \all{das}) et \all{besten} (pluriel accusatif après \all{die}) ; phrase 2 : particule \all{aus} à la toute fin, après tous les compléments ; phrase 4 : \all{den frischen Käse} (accusatif masculin $\rightarrow$ \all{-en}) mais \all{die kalte Milch} (accusatif féminin $\rightarrow$ \all{-e}) ; phrase 5 : deux verbes séparables, donc deux particules finales (\all{ein} et \all{zurück}).
\textbf{B. Version — traduction modèle :}
« Chaque matin, M. Mbarga achète un journal au kiosque. Le journal n'est pas cher, mais les magazines sont plus chers. "Je prends le journal le moins cher", dit-il, et il donne l'argent au vendeur. Ensuite, il prend le bus pour aller au travail. »
\textbf{Barème indicatif :} thème 5 $\times$ 1 pt ; version 5 pts (fidélité 3 pts, qualité du français 2 pts).
\end{corrige}

\begin{corrige}[Exercice 11 — Production écrite guidée : Meine Einkäufe am Wochenende]
\textbf{Proposition de rédaction modèle :}
\begin{textebox}[Meine Einkäufe am Wochenende]
Am Samstagmorgen kaufe ich mit meiner Mutter auf dem Markt in Yaoundé ein. Der Markt ist größer und billiger als der Supermarkt. Zuerst kaufen wir frisches Obst und Gemüse: Bananen, Tomaten und Zwiebeln. Die Tomaten sind dort am frischesten. Dann nehme ich ein kleines Brot für meinen Bruder mit. Ich gebe nicht viel Geld aus, denn die Preise sind gut. Die Verkäufer sind immer freundlich. Um elf Uhr kommen wir nach Hause zurück. Ich finde: Der Markt ist besser als das Geschäft, weil die Waren dort frischer sind.
\end{textebox}
Traduction : « Samedi matin, je fais mes courses avec ma mère au marché de Yaoundé. Le marché est plus grand et moins cher que le supermarché. D'abord, nous achetons des fruits et légumes frais : bananes, tomates et oignons. Les tomates y sont les plus fraîches. Ensuite, j'emporte un petit pain pour mon frère. Je ne dépense pas beaucoup d'argent, car les prix sont bons. Les vendeurs sont toujours aimables. À onze heures, nous rentrons à la maison. Je trouve que le marché est meilleur que le magasin, parce que les marchandises y sont plus fraîches. »
\textbf{Vérification des consignes :} comparatifs et superlatifs : \all{größer und billiger als}, \all{am frischesten}, \all{besser als}, \all{frischer} ; verbes séparables : \all{kaufe ... ein}, \all{nehme ... mit}, \all{gebe ... aus}, \all{kommen ... zurück} ; adjectifs déclinés : \all{frisches Obst} (accusatif neutre sans article $\rightarrow$ \all{-es}), \all{ein kleines Brot} (accusatif neutre après \all{ein} $\rightarrow$ \all{-es}).
\textbf{Barème indicatif (10 pts) :} respect des consignes et du plan 2 pts ; comparatifs et superlatifs corrects 2 pts ; verbes séparables bien placés 2 pts ; déclinaisons des adjectifs 2 pts ; vocabulaire et orthographe (majuscules, Umlaute) 2 pts.
\textbf{Points de vigilance :} majuscule à tous les noms ; particule séparable en toute fin de phrase ; ne jamais écrire \all{mehr billiger} (double comparatif interdit) : on dit \all{billiger als}. Après \all{weil}, le verbe conjugué va à la fin : \all{..., weil die Waren dort frischer sind.}
\end{corrige}

\newpage

\coursec{Fiche bilan}

\begin{popfigure}
\begin{tikzpicture}[
  mindmap,
  grow cyclic,
  every node/.style={concept, align=center, text=white, font=\small},
  root concept/.append style={concept color=popblue, font=\bfseries\small, minimum size=3.2cm},
  level 1 concept/.append style={level distance=3.6cm, sibling angle=51.4, minimum size=2.3cm},
  level 2 concept/.append style={level distance=2.2cm, minimum size=1.8cm, font=\scriptsize},
]
\node[root concept]{Konsum und Dienstleistungen}
  child[concept color=poporange]{ node[align=center]{Texte\\ \all{Einkaufen in Deutschland}}
    child[concept color=poporange!70]{ node[align=center]{compréhension\\ globale + détaillée} }
  }
  child[concept color=popgreen]{ node[align=center]{Lexique\\ \all{kaufen, der Preis, die Ware}}
    child[concept color=popgreen!70]{ node[align=center]{articles\\ + pluriels} }
  }
  child[concept color=poppurple]{ node[align=center]{Déclinaison\\ de l'adjectif}
    child[concept color=poppurple!70]{ node {faible / mixte / forte} }
  }
  child[concept color=poppink]{ node[align=center]{Comparatif\\ Superlatif}
    child[concept color=poppink!70]{ node {\all{billiger, am billigsten}} }
  }
  child[concept color=popgold]{ node[align=center]{Verbes à\\ particule séparable}
    child[concept color=popgold!70]{ node {\all{einkaufen : ich kaufe ein}} }
  }
  child[concept color=popblue!60]{ node[align=center]{Expression\\ guidée}
    child[concept color=popblue!40]{ node[align=center]{réponses, résumé,\\ court paragraphe} }
  };
\end{tikzpicture}
\legende{Carte mentale de la leçon : la consommation de biens et de services}
\end{popfigure}

\begin{bilan}
\textbf{1. Lexique clé (à connaître avec article et pluriel)}
\begin{itemize}
  \item \all{der Konsum} (sans pluriel) : la consommation ; \all{der Kunde, -n} : le client ; \all{die Kundin, -nen} : la cliente.
  \item \all{das Geschäft, -e} : le magasin ; \all{der Laden, \"-} : la boutique ; \all{der Markt, \"-e} : le marché ; \all{der Supermarkt, \"-e} : le supermarché.
  \item \all{die Ware, -n} : la marchandise ; \all{das Produkt, -e} : le produit ; \all{die Dienstleistung, -en} : la prestation de service.
  \item \all{der Preis, -e} : le prix ; \all{das Angebot, -e} : l'offre, la promotion ; \all{der Rabatt, -e} : la réduction ; \all{die Quittung, -en} : le ticket de caisse, le reçu.
  \item \all{kaufen} : acheter ; \all{verkaufen} : vendre ; \all{bezahlen} : payer ; \all{ausgeben (gibt aus, gab aus, hat ausgegeben)} : dépenser ; \all{sparen} : économiser.
  \item Expressions utiles : \all{Wie viel kostet das?} « Combien cela coûte-t-il ? » ; \all{Das ist zu teuer.} « C'est trop cher. » ; \all{Ich nehme es.} « Je le prends. » ; \all{Ich schaue mich nur um.} « Je regarde seulement. »
\end{itemize}

\textbf{2. Grammaire à connaître par c\oe ur}
\begin{center}
\begin{tabularx}{\linewidth}{|X|X|}
\hline
\rowcolor{popblueL} \textbf{Règle} & \textbf{Exemple} \\
\hline
Adjectif épithète : déclinaison \textbf{faible} après \all{der/die/das} : terminaisons \emph{-e} ou \emph{-en}. & \all{das billige Produkt, die billigen Produkte} \\
\hline
Déclinaison \textbf{mixte} après \all{ein, kein, mein} : \emph{-er/-es/-e} au nominatif singulier, \emph{-en} ailleurs (sauf accusatif neutre \emph{-es}). & \all{ein guter Preis, kein neues Geschäft} \\
\hline
Déclinaison \textbf{forte} sans article : l'adjectif porte lui-même la marque du cas. & \all{guter Kaffee, frische Waren} \\
\hline
Adjectif attribut : \textbf{jamais} de terminaison après \all{sein/werden/bleiben}. & \all{Der Preis ist hoch.} \\
\hline
Comparatif : adjectif + \emph{-er} (+ \all{als}) ; superlatif : \all{am} + adjectif + \emph{-sten} (adverbe) ou \all{der/die/das} + adjectif + \emph{-ste} (épithète). & \all{billiger als, am billigsten, das billigste Angebot} \\
\hline
Superlatif en \emph{-esten} après \emph{-d, -t, -s, -ß, -z} : \all{kalt -- am kältesten}. & \all{Der Winter ist am kältesten.} \\
\hline
Irréguliers : \all{gut -- besser -- am besten} ; \all{gern -- lieber -- am liebsten} ; \all{viel -- mehr -- am meisten} ; \all{hoch -- höher -- am höchsten} ; \all{nah -- näher -- am nächsten}. & \all{Dieses Geschäft ist besser.} \\
\hline
Umlaut au comparatif pour les monosyllabes en \emph{a, o, u} : \all{alt, jung, kurz, lang, stark, warm, kalt...} & \all{alt -- älter -- am ältesten} \\
\hline
Adjectifs en \emph{-el, -er} : le \emph{e} disparaît ou se contracte : \all{dunkel -- dunkler} ; \all{teuer -- teurer}. & \all{Das Angebot ist teurer.} \\
\hline
Verbes à particule séparable : la particule va \textbf{en fin de phrase} au Präsens ; \emph{ge-} s'insère entre particule et radical au Perfekt. & \all{Ich kaufe heute ein. -- Ich habe eingekauft.} \\
\hline
\end{tabularx}
\end{center}

\textbf{3. Méthodes express}
\begin{itemize}
  \item \textbf{Comprendre un texte} : lire le titre et la source, repérer le thème, première lecture globale (qui ? quoi ? où ?), puis lecture détaillée avec le glossaire.
  \item \textbf{Répondre aux questions} : reprendre les mots de la question, répondre par une phrase complète, verbe conjugué en 2\textsuperscript{e} position.
  \item \textbf{Résumer} : 3 à 5 phrases au Präsens, avec ses propres mots, sans recopier de longues phrases du texte.
  \item \textbf{Traduire} : vérifier la place du verbe, la déclinaison de l'adjectif, la particule séparable en fin de phrase.
\end{itemize}
\end{bilan}

\begin{aretenir}[Checklist avant le Bac]
\begin{itemize}
  \item[$\square$] Je connais le vocabulaire de la consommation avec l'article et le pluriel de chaque nom.
  \item[$\square$] Je sais distinguer un bien (\all{die Ware}) d'un service (\all{die Dienstleistung}).
  \item[$\square$] Je sais choisir entre déclinaison faible, mixte et forte de l'adjectif.
  \item[$\square$] Je n'accorde jamais l'adjectif attribut après \all{sein} : \all{Der Preis ist hoch.}
  \item[$\square$] Je forme le comparatif avec \emph{-er} + \all{als} et le superlatif avec \all{am ...-sten}.
  \item[$\square$] Je connais par c\oe ur les comparatifs irréguliers : \all{gut -- besser -- am besten}.
  \item[$\square$] Je pense à l'Umlaut au comparatif : \all{alt -- älter}, \all{jung -- jünger}.
  \item[$\square$] Je place la particule séparable en fin de phrase : \all{Er kauft heute ein.}
  \item[$\square$] Je forme le Perfekt des verbes séparables : \all{eingekauft}, \all{mitgebracht}.
  \item[$\square$] Je mets une majuscule à tous les noms et le verbe conjugué en 2\textsuperscript{e} position.
  \item[$\square$] Je sais répondre à une question de compréhension par une phrase complète.
  \item[$\square$] Je sais rédiger un court paragraphe (5 à 8 lignes) sur mes habitudes de consommation.
\end{itemize}
\end{aretenir}

\findecours

\end{document}
